% Équations, inéquations, systèmes — Mathématiques Tle Tech — Cours signé OnBuch+
% Programme officiel MINESEC. Compiler avec : tectonic N01-equations-inequations-systemes.tex
\newcommand{\DOCMATIERE}{Mathématiques}
\newcommand{\DOCNIVEAU}{Tle Tech}
\newcommand{\DOCMODULE}{Mathématiques – classe de Terminale technique}
\newcommand{\DOCLECON}{N01}
\newcommand{\DOCTITRE}{Équations, inéquations, systèmes}
% OnBuch+ — préambule « cours signés » ENRICHI (compilé avec tectonic / XeLaTeX)
% Système visuel « Pop » d'OnBuch+ : fond crème (jamais blanc), bordures encre
% épaisses, ombres « dures » décalées, titres Archivo Black en capitales.
% Version MATHÉMATIQUES : graphes (pgfplots/tikz), boîtes pédagogiques
% (définition, propriété, méthode, attention, le savais-tu, exercice résolu,
% exercice, TP, bilan), page de garde et signature OnBuch+ ancrée en bas.
%
% Macros attendues dans main.tex AVANT \input{preamble} :
%   \DOCMATIERE  \DOCNIVEAU  \DOCMODULE  \DOCLECON (numéro)  \DOCTITRE
\documentclass[11pt]{article}

\usepackage{fontspec}
\usepackage[french]{babel}
\usepackage[a4paper,margin=2cm,top=3.4cm,bottom=2.8cm,headheight=1.4cm,footskip=1.3cm]{geometry}
\usepackage{amsmath,amssymb,amsfonts}
\usepackage{mathtools} % \xmapsto, \coloneqq... souvent utilisés par les modèles
\usepackage{cancel} % simplifications barrées
% \abs{z} (module, valeur absolue) et \norm{u} (norme), utilisés par les modèles
\providecommand{\abs}{}\let\abs\relax\DeclarePairedDelimiter{\abs}{\lvert}{\rvert}
\providecommand{\norm}{}\let\norm\relax\DeclarePairedDelimiter{\norm}{\lVert}{\rVert}
\usepackage[table]{xcolor} % [table] : fournit aussi \rowcolors (alternance de lignes)
\usepackage{pagecolor}
\usepackage{tikz}
\usetikzlibrary{arrows.meta,positioning,calc,shapes.geometric,shapes.misc,shapes.arrows,decorations.pathmorphing,decorations.pathreplacing,decorations.markings,patterns,shadows,mindmap,trees,fit,backgrounds,matrix,intersections,angles,quotes}
\usepackage{pgfplots}
\usepackage{pgf-pie} % diagrammes circulaires \pie (statistiques)
\pgfplotsset{compat=1.18}
\usepgfplotslibrary{fillbetween,groupplots} % aires entre courbes (calcul intégral)
\usepackage{enumitem}
\usepackage{fancyhdr}
% [most] charge toutes les bibliothèques tcolorbox (listings, minted, raster,
% documentation...) et gonfle inutilement la mémoire TeX sur un document long
% (~300 boîtes) : on ne charge que ce qui est réellement utilisé.
\usepackage[skins,breakable]{tcolorbox}
\usepackage{graphicx}
\usepackage{colortbl}
\usepackage{booktabs}
\usepackage{tabularx}
\usepackage{diagbox} % case d'en-tête diagonale des tableaux à double entrée
% Colonnes extensibles centrée / gauche / droite (C, L, R), souvent utilisées par les modèles
\newcolumntype{C}{>{\centering\arraybackslash}X}
\newcolumntype{L}{>{\raggedright\arraybackslash}X}
\newcolumntype{R}{>{\raggedleft\arraybackslash}X}
\usepackage{array}
\usepackage{multicol}
\usepackage{siunitx}
% parse-numbers=false : accepte aussi \SI{2,5 \times 10^{-3}}{...}
\sisetup{locale=FR,output-decimal-marker={,},group-minimum-digits=4,parse-numbers=false}
\DeclareSIUnit\ohm{\ensuremath{\symup{\Omega}}} % U+2126 absent de la police
\usepackage{qrcode}
% \degree : commande fréquemment utilisée par les modèles pour un angle ou une
% température en dehors de \SI{}{} (non fournie par les paquets chargés).
\providecommand{\degree}{^{\circ}}
% \qed (fin de démonstration), utilisé par les modèles sans amsthm
\providecommand{\qed}{\hfill$\square$}
% \barre : double barre des valeurs interdites dans un tableau de variations (tkz-tab)
\providecommand{\barre}{\ensuremath{\|}}
% environnement proof (amsthm non chargé) utilisé par les modèles
\ifdefined\proof\else\newenvironment{proof}[1][Démonstration]{\par\noindent\textit{#1.}\ }{\qed\par}\fi
\usepackage{lastpage}
\usepackage{hyperref}
\hypersetup{hidelinks}

% ============================================================
% Palette « Pop » OnBuch+ — mêmes hex que la classe P (plus_ui.dart)
% ============================================================
\definecolor{poporange}{HTML}{F59321}
\definecolor{popdark}{HTML}{E07A0C}
\definecolor{popcream}{HTML}{FFF6E8}
\definecolor{popink}{HTML}{1C1714}
\definecolor{poppurple}{HTML}{7A5AE0}
\definecolor{popgreen}{HTML}{2E9E6A}
\definecolor{poppink}{HTML}{E0457B}
\definecolor{popblue}{HTML}{2F7DE1}
\definecolor{popgold}{HTML}{C98A12}
\definecolor{popmuted}{HTML}{8A7F76}
\definecolor{popline}{HTML}{EADFD2}
\definecolor{popgray}{HTML}{8A7F76}
\colorlet{popgrey}{popgray}
% Alias tolérants (le modèle nomme parfois les couleurs différemment)
\colorlet{popwhite}{white}
\colorlet{popblack}{popink}
\colorlet{popred}{poppink}
\colorlet{popyellow}{popgold}
% Teintes claires pour fonds de tableaux / zones
\colorlet{popblueL}{popblue!10!white}
\colorlet{poporangeL}{poporange!14!white}
\colorlet{popgreenL}{popgreen!12!white}
\colorlet{poppurpleL}{poppurple!12!white}
\colorlet{poppinkL}{poppink!12!white}
\colorlet{popgoldL}{popgold!14!white}

\pagecolor{popcream}

% ============================================================
% Typographie
% ============================================================
\defaultfontfeatures{Ligatures=TeX}
\setmainfont{PlusJakartaSans-Medium.ttf}[Path=../../../fonts/,BoldFont=PlusJakartaSans-Bold.ttf]
% Maths en sans-serif (Fira Math) avec chiffres et lettres droites de Plus
% Jakarta, pour que les formules se fondent dans le texte.
\usepackage[math-style=ISO,bold-style=ISO]{unicode-math}
\setmathfont{FiraMath-Regular.otf}[Path=../../../fonts/]
\setmathfont{PlusJakartaSans-Medium.ttf}[Path=../../../fonts/,range={up/num,up/{Latin,latin},\mathup}]
\usepackage{icomma} % virgule décimale sans espace en mode math
% Caractères mathématiques Unicode absents des polices de texte (Plus Jakarta,
% Archivo Black) : rendus par leur équivalent mathématique, en texte comme en
% formule — sinon ils s'affichent en carré vide (ex. « Arithmétique dans ℤ »).
\usepackage{newunicodechar}
\newunicodechar{ℝ}{\ensuremath{\mathbb{R}}}
\newunicodechar{ℤ}{\ensuremath{\mathbb{Z}}}
\newunicodechar{ℂ}{\ensuremath{\mathbb{C}}}
\newunicodechar{ℕ}{\ensuremath{\mathbb{N}}}
\newunicodechar{ℚ}{\ensuremath{\mathbb{Q}}}
\newunicodechar{𝔻}{\ensuremath{\mathbb{D}}}
\newunicodechar{α}{\ensuremath{\alpha}}
\newunicodechar{β}{\ensuremath{\beta}}
\newunicodechar{γ}{\ensuremath{\gamma}}
\newunicodechar{δ}{\ensuremath{\delta}}
\newunicodechar{ε}{\ensuremath{\varepsilon}}
\newunicodechar{θ}{\ensuremath{\theta}}
\newunicodechar{λ}{\ensuremath{\lambda}}
\newunicodechar{μ}{\ensuremath{\mu}}
\newunicodechar{σ}{\ensuremath{\sigma}}
\newunicodechar{τ}{\ensuremath{\tau}}
\newunicodechar{φ}{\ensuremath{\varphi}}
\newunicodechar{ψ}{\ensuremath{\psi}}
\newunicodechar{ω}{\ensuremath{\omega}}
\newunicodechar{Σ}{\ensuremath{\Sigma}}
% majuscules produites par \MakeUppercase dans les titres (α-aminés -> Α)
\newunicodechar{Α}{\ensuremath{\alpha}}
\newunicodechar{Β}{\ensuremath{\beta}}
\newunicodechar{Γ}{\ensuremath{\Gamma}}
\newunicodechar{Δ}{\ensuremath{\Delta}}
\newunicodechar{Ε}{\ensuremath{\varepsilon}}
\newunicodechar{Θ}{\ensuremath{\Theta}}
\newunicodechar{Λ}{\ensuremath{\Lambda}}
\newunicodechar{Μ}{\ensuremath{\mu}}
\newunicodechar{Τ}{\ensuremath{\tau}}
\newunicodechar{Φ}{\ensuremath{\Phi}}
\newunicodechar{Ψ}{\ensuremath{\Psi}}
\newunicodechar{Ω}{\ensuremath{\Omega}}
\newunicodechar{↦}{\ensuremath{\mapsto}}
\newunicodechar{∈}{\ensuremath{\in}}
\newunicodechar{∉}{\ensuremath{\notin}}
\newunicodechar{∧}{\ensuremath{\wedge}}
\newunicodechar{⇔}{\ensuremath{\Leftrightarrow}}
\newunicodechar{⟺}{\ensuremath{\Longleftrightarrow}}
\newunicodechar{⇒}{\ensuremath{\Rightarrow}}
\newunicodechar{⟹}{\ensuremath{\Longrightarrow}}
\newunicodechar{∘}{\ensuremath{\circ}}
\newunicodechar{∪}{\ensuremath{\cup}}
\newunicodechar{∩}{\ensuremath{\cap}}
\newunicodechar{≡}{\ensuremath{\equiv}}
\newunicodechar{⊂}{\ensuremath{\subset}}
\newunicodechar{⊥}{\ensuremath{\perp}}
\newunicodechar{∃}{\ensuremath{\exists}}
\newunicodechar{∀}{\ensuremath{\forall}}
\newunicodechar{∛}{\ensuremath{\sqrt[3]{\,}}}
\newunicodechar{∜}{\ensuremath{\sqrt[4]{\,}}}
\newunicodechar{⃗}{\ensuremath{{}^{\to}}}
\newunicodechar{ⁿ}{\ifmmode^{n}\else\textsuperscript{\ensuremath{n}}\fi}
\newunicodechar{ˣ}{\ifmmode^{x}\else\textsuperscript{\ensuremath{x}}\fi}
\newunicodechar{ᵇ}{\ifmmode^{b}\else\textsuperscript{\ensuremath{b}}\fi}
\newunicodechar{ʳ}{\ifmmode^{r}\else\textsuperscript{\ensuremath{r}}\fi}
\newunicodechar{ᵘ}{\ifmmode^{u}\else\textsuperscript{\ensuremath{u}}\fi}
\newunicodechar{ʸ}{\ifmmode^{y}\else\textsuperscript{\ensuremath{y}}\fi}
\newunicodechar{ᵃ}{\ifmmode^{a}\else\textsuperscript{\ensuremath{a}}\fi}
\newunicodechar{ᵉ}{\ifmmode^{e}\else\textsuperscript{\ensuremath{e}}\fi}
\newunicodechar{ᵏ}{\ifmmode^{k}\else\textsuperscript{\ensuremath{k}}\fi}
\newunicodechar{ᵖ}{\ifmmode^{p}\else\textsuperscript{\ensuremath{p}}\fi}
\newunicodechar{ᴺ}{\ifmmode^{N}\else\textsuperscript{\ensuremath{N}}\fi}
\newunicodechar{ᶜ}{\ifmmode^{c}\else\textsuperscript{\ensuremath{c}}\fi}
\newunicodechar{ʰ}{\ifmmode^{h}\else\textsuperscript{\ensuremath{h}}\fi}
\newunicodechar{ᵠ}{\ifmmode^{q}\else\textsuperscript{\ensuremath{q}}\fi}
\newunicodechar{ₙ}{\ifmmode_{n}\else\textsubscript{\ensuremath{n}}\fi}
\newunicodechar{ₐ}{\ifmmode_{a}\else\textsubscript{\ensuremath{a}}\fi}
\newunicodechar{ᵢ}{\ifmmode_{i}\else\textsubscript{\ensuremath{i}}\fi}
\newunicodechar{ᵣ}{\ifmmode_{r}\else\textsubscript{\ensuremath{r}}\fi}
\newunicodechar{ₖ}{\ifmmode_{k}\else\textsubscript{\ensuremath{k}}\fi}
\newunicodechar{ᵦ}{\ifmmode_{\beta}\else\textsubscript{\ensuremath{\beta}}\fi}
\newunicodechar{ₓ}{\ifmmode_{x}\else\textsubscript{\ensuremath{x}}\fi}
\newfontfamily\popdisplay{ArchivoBlack-Regular.ttf}[Path=../../../fonts/]
\newfontfamily\poplogo{SpaceGrotesk-Bold.ttf}[Path=../../../fonts/]
\newfontfamily\popsemi{PlusJakartaSans-SemiBold.ttf}[Path=../../../fonts/]
\newfontfamily\popxbold{PlusJakartaSans-ExtraBold.ttf}[Path=../../../fonts/]
\newfontfamily\popxboldls{PlusJakartaSans-ExtraBold.ttf}[Path=../../../fonts/,LetterSpace=3.0]

\setlist[enumerate]{leftmargin=1.7em,itemsep=5pt,topsep=4pt}
\setlist[itemize]{leftmargin=1.7em,itemsep=4pt,topsep=4pt,label={\color{poporange}\textbullet}}
\setlength{\parindent}{0pt}
\setlength{\parskip}{5pt}
\renewcommand{\arraystretch}{1.25}

% pgfplots : style Pop par défaut
\pgfplotsset{
  every axis/.append style={axis line style={popink,line width=1pt},tick style={popink},
    grid style={popline,line width=0.5pt},label style={font=\small},tick label style={font=\footnotesize}},
  popcurve/.style={poporange,line width=1.8pt,no markers,smooth},
}
% Même style disponible dans un \draw TikZ simple (hors pgfplots)
\tikzset{popcurve/.style={poporange,line width=1.8pt}}

% ============================================================
% Pill / squiggle
% ============================================================
\newcommand{\poppill}[3]{%
  \tikz[baseline=-0.55ex]\node[draw=popink,line width=1.2pt,rounded corners=2.6pt,%
    fill=#2,inner xsep=6.5pt,inner ysep=3pt]{%
    \popxboldls\fontsize{7}{7}\selectfont\color{#3}\MakeUppercase{#1}};%
}
\newcommand{\popsquiggle}[1]{%
  \tikz[baseline=-0.4ex]{
    \draw[#1,line width=2.6pt,line cap=round]
      plot [smooth] coordinates {(0,0.09) (0.34,0.01) (0.68,0.21) (1.02,0.01) (1.36,0.10)};
  }%
}
% Mot-clé surligné (marqueur orange sous le texte)
\newcommand{\cle}[1]{{\popxbold\color{popdark}#1}}

% ============================================================
% En-tête / pied
% ============================================================
\pagestyle{fancy}
\fancyhf{}
\renewcommand{\headrulewidth}{0pt}
\renewcommand{\footrulewidth}{0pt}
\lhead{\raisebox{-0.15ex}{\poplogo\fontsize{13}{13}\selectfont\color{popink}OnBuch\color{poporange}+}}
\rhead{\poppill{\DOCMATIERE\ \textbullet\ \DOCNIVEAU\ \textbullet\ Leçon \DOCLECON}{white}{popink}}
\lfoot{\popxbold\fontsize{7.6}{7.6}\selectfont\color{popmuted}\MakeUppercase{Cours signé OnBuch+}\ \textbullet\ {\popsemi\DOCTITRE}}
\rfoot{\tikz[baseline=-0.6ex]\node[rounded corners=8.5pt,fill=popink,minimum height=17pt,minimum width=17pt,inner xsep=5pt,inner sep=0]{\popxbold\fontsize{8}{8}\selectfont\color{popcream}\thepage\,/\,\pageref*{LastPage}};}

% ============================================================
% Titres
% ============================================================
\newcommand{\chaptitre}[2]{%
  \begin{tcolorbox}[enhanced,colback=poporange,colframe=popink,boxrule=2.6pt,arc=11pt,
      drop shadow=popink,left=15pt,right=15pt,top=12pt,bottom=11pt,before skip=3pt,after skip=18pt]
    \poppill{#2}{popink}{popcream}\par\vspace{9pt}
    {\raggedright\hyphenpenalty=10000\exhyphenpenalty=10000\popdisplay\fontsize{22}{24}\selectfont\color{popcream}\MakeUppercase{#1}\par}
  \end{tcolorbox}
}
% Section numérotée : pastille orange avec le numéro + titre Archivo
\newcounter{popsec}
\newcommand{\coursec}[1]{%
  \stepcounter{popsec}\setcounter{popsub}{0}%
  \par\vspace{16pt}\noindent
  \begin{minipage}{\linewidth}
  \tikz[baseline=-0.7ex]\node[draw=popink,line width=1.6pt,fill=poporange,rounded corners=4pt,minimum size=22pt,inner sep=0]{\popdisplay\fontsize{12}{12}\selectfont\color{popcream}\thepopsec};%
  \hspace{8pt}{\popdisplay\fontsize{15}{17}\selectfont\color{popink}\MakeUppercase{#1}}\par
  \vspace{4pt}\noindent{\color{poporange}\rule{36pt}{3.6pt}}
  \end{minipage}\par\nopagebreak\vspace{8pt}
}
\newcounter{popsub}
\newcommand{\courssub}[1]{%
  \stepcounter{popsub}%
  \par\vspace{9pt}\noindent{\popxbold\fontsize{11.5}{13}\selectfont\color{popink}{\color{poporange}\thepopsec.\thepopsub}\hspace{6pt}#1}\par\nopagebreak\vspace{4pt}
}

% ============================================================
% Boîtes pédagogiques — même recette : fond blanc, bordure épaisse colorée,
% ombre dure encre, badge plein. Titre optionnel [#1].
% ============================================================
\tcbset{popbase/.style={breakable,enhanced,colback=white,boxrule=2.3pt,arc=9pt,
  shadow={3pt}{-3pt}{0pt}{popink},left=14pt,right=14pt,top=11pt,bottom=11pt,before skip=11pt,after skip=15pt,
  fontupper=\popsemi\fontsize{10.6}{14.5}\selectfont\color{popink},
  fontlower=\popsemi\fontsize{10.6}{14.5}\selectfont\color{popink}}}
% Coche dessinée (le glyphe ✓ n'existe pas dans les polices du document)
\newcommand{\popcheck}[1]{\tikz[baseline=-0.5ex]\draw[#1,line width=1.6pt,line cap=round,line join=round](-0.13,0.0)--(-0.04,-0.09)--(0.13,0.1);}
\newcommand{\popboxhead}[3]{\poppill{#1}{#2}{white}\ifx\relax#3\relax\else\hspace{6pt}{\popxbold\fontsize{10.6}{12}\selectfont\color{popink}#3}\fi\par\vspace{7pt}}

\newtcolorbox{aretenir}[1][]{popbase,colframe=popgreen,before upper={\popboxhead{À retenir}{popgreen}{#1}}}
\newtcolorbox{exemplebox}[1][]{popbase,colframe=poporange,before upper={\popboxhead{Exemple}{poporange}{#1}}}
\newtcolorbox{definition}[1][]{popbase,colframe=popblue,before upper={\popboxhead{Définition}{popblue}{#1}}}
\newtcolorbox{propriete}[1][]{popbase,colframe=poppurple,before upper={\popboxhead{Propriété}{poppurple}{#1}}}
\newtcolorbox{methode}[1][]{popbase,colframe=popgold,before upper={\popboxhead{Méthode}{popgold}{#1}}}
\newtcolorbox{attention}[1][]{popbase,colframe=poppink,before upper={\popboxhead{Attention}{poppink}{#1}}}
\newtcolorbox{experience}[1][]{popbase,colframe=popblue,colback=popblueL,before upper={\popboxhead{Expérience}{popblue}{#1}}}
% « Le savais-tu ? » : carte violette pleine (contexte, culture, Cameroun)
\newtcolorbox{savaistu}[1][]{popbase,colback=poppurple,colframe=popink,
  fontupper=\popsemi\fontsize{10.4}{14.2}\selectfont\color{white},
  before upper={\poppill{Le savais-tu\,?}{white}{poppurple}\ifx\relax#1\relax\else\hspace{6pt}{\popxbold\color{white}#1}\fi\par\vspace{7pt}}}
% Exercice résolu : énoncé en haut, \tcblower puis la solution sur fond crème
\newcounter{popexo}\newcounter{popexr}
\newtcolorbox{exoresolu}[1][]{popbase,colframe=poporange,colbacklower=popcream,
  segmentation style={popink,line width=1.2pt,dashed},
  before upper={\stepcounter{popexr}\popboxhead{Exercice résolu \thepopexr}{poporange}{#1}},
  before lower={\poppill{Solution}{popink}{popcream}\par\vspace{6pt}}}
% Exercice (non résolu en place) — la correction est dans la partie Corrigés
\newtcolorbox{exercice}[1][]{popbase,colframe=popink,
  before upper={\stepcounter{popexo}\popboxhead{Exercice \thepopexo}{popink}{#1}}}
% Corrigé
\newtcolorbox{corrige}[1][]{popbase,colframe=popgreen,colback=popgreenL,
  before upper={\popboxhead{Corrigé}{popgreen}{#1}}}
% Objectifs / prérequis / bilan
\newtcolorbox{objectifs}{popbase,colframe=popink,colback=poporangeL,
  before upper={\popboxhead{Objectifs de la leçon}{popink}{}}}
\newtcolorbox{prerequis}{popbase,colframe=popink,colback=popblueL,
  before upper={\popboxhead{Prérequis}{popblue}{}}}
\newtcolorbox{bilan}{popbase,colframe=popink,colback=white,boxrule=2.8pt,
  before upper={\popboxhead{Fiche bilan}{poporange}{L'essentiel à connaître pour l'examen}}}
% Figure encadrée avec légende
\newtcolorbox{popfigure}[1][]{enhanced,breakable=false,colback=white,colframe=popline,boxrule=1.4pt,arc=8pt,
  left=10pt,right=10pt,top=10pt,bottom=8pt,before skip=10pt,after skip=12pt,halign=center,
  lower separated=false,halign lower=center,
  fontlower=\popsemi\fontsize{9}{11.5}\selectfont\color{popmuted},
  #1}
\newcommand{\legende}[1]{\par\vspace{6pt}{\popsemi\fontsize{9}{11.5}\selectfont\color{popmuted}#1}}

% ============================================================
% Signature OnBuch+
% ============================================================
\newcommand{\poplogoinline}[1]{{\poplogo\fontsize{#1}{#1}\selectfont\color{popink}OnBuch\color{poporange}+}}
\newcommand{\signatureonbuch}{%
  \begin{tcolorbox}[enhanced,colback=popink,colframe=popink,boxrule=2.2pt,arc=10pt,
      drop shadow=poporange,left=14pt,right=14pt,top=10pt,bottom=10pt,before skip=0pt,after skip=0pt]
    \tikz[baseline=-0.7ex]\node[circle,fill=popgreen,draw=popcream,line width=1.3pt,minimum size=22pt,inner sep=0]{\popcheck{white}};%
    \hspace{9pt}\begin{minipage}[c]{0.62\linewidth}
      {\popxbold\fontsize{12}{13}\selectfont\color{popcream}Signé\ \textbullet\ {\poplogo OnBuch\color{poporange}+}}\par\vspace{2pt}
      {\raggedright\popsemi\fontsize{8}{10.5}\selectfont\color{popline}\MakeUppercase{Équipe pédagogique OnBuch+}\\\MakeUppercase{Conforme au programme officiel MINESEC}\par}
    \end{minipage}\hfill
    {\poplogo\fontsize{22}{22}\selectfont\color{popcream}OnBuch\color{poporange}+}
  \end{tcolorbox}%
}

% ============================================================
% Page de garde
% #1 titre · #2 matière · #3 niveau · #4 module · #5 n° leçon · #6 durée ·
% #7 description / objectifs (texte ou itemize)
% ============================================================
\newcommand{\pagegarde}[7]{%
  \thispagestyle{empty}
  \newgeometry{a4paper,margin=1.7cm,top=1.6cm,bottom=1.4cm}%
  \begin{tikzpicture}[remember picture,overlay]
    \fill[poporange!18] ([xshift=-3.2cm,yshift=-3.6cm]current page.north east) circle (4.6cm);
    \fill[poppurple!14] ([xshift=2.4cm,yshift=7.6cm]current page.south west) circle (3.8cm);
  \end{tikzpicture}%
  {\poplogo\fontsize{30}{30}\selectfont\color{popink}OnBuch\color{poporange}+}\hfill
  \raisebox{0.6ex}{\poppill{Cours signé \textbullet\ Premium}{popink}{popcream}}\par
  \vspace{1pt}\popsquiggle{poppurple}\par
  \vspace{8pt}
  \begin{tcolorbox}[enhanced,colback=poporange,colframe=popink,boxrule=2.8pt,arc=13pt,
      drop shadow=popink,left=18pt,right=18pt,top=12pt,bottom=13pt,after skip=10pt]
    \poppill{#2 \textbullet\ #3}{popink}{popcream}\hspace{5pt}\poppill{#4}{white}{popink}\par\vspace{8pt}
    {\popdisplay\fontsize{14}{14}\selectfont\color{popink}LEÇON #5}\par\vspace{4pt}
    {\raggedright\hyphenpenalty=10000\exhyphenpenalty=10000\popdisplay\fontsize{25}{27}\selectfont\color{popcream}\MakeUppercase{#1}\par}
    \vspace{8pt}
    {\popxbold\fontsize{9.5}{9.5}\selectfont\color{popink}\MakeUppercase{Durée conseillée : #6}}
  \end{tcolorbox}
  \begin{tcolorbox}[enhanced,colback=white,colframe=popink,boxrule=2.2pt,arc=10pt,
      drop shadow=popink,left=16pt,right=16pt,top=10pt,bottom=10pt,after skip=8pt]
    \poppill{Dans cette leçon}{poppurple}{white}\par\vspace{5pt}
    \popsemi\fontsize{9.3}{12.6}\selectfont\color{popink}#7
  \end{tcolorbox}
  \noindent\begin{minipage}[t]{0.487\linewidth}
    \begin{tcolorbox}[enhanced,colback=popgreen,colframe=popink,boxrule=2.2pt,arc=10pt,
        drop shadow=popink,left=11pt,right=11pt,top=10pt,bottom=10pt,height=3.1cm,valign=center]
      \tcbox[colback=white,colframe=popink,boxrule=1.3pt,arc=5pt,boxsep=0pt,nobeforeafter,
        left=3pt,right=3pt,top=3pt,bottom=3pt,valign=center]{\qrcode[height=1.9cm]{https://play.google.com/store/apps/details?id=com.luvvix.onbuch&hl=fr}}%
      \hspace{8pt}\begin{minipage}[c]{0.5\linewidth}
        {\popdisplay\fontsize{11}{12.5}\selectfont\color{white}\MakeUppercase{Télécharge OnBuch}}\par\vspace{4pt}
        {\raggedright\popsemi\fontsize{8.4}{11}\selectfont\color{white}Gratuit \textbullet\ Google Play\\Tuteur IA, annales, concours, résultats\par}
      \end{minipage}
    \end{tcolorbox}
  \end{minipage}\hfill
  \begin{minipage}[t]{0.487\linewidth}
    \begin{tcolorbox}[enhanced,colback=poppurple,colframe=popink,boxrule=2.2pt,arc=10pt,
        drop shadow=popink,left=11pt,right=11pt,top=10pt,bottom=10pt,height=3.1cm,valign=center]
      \tikz[baseline=-0.7ex]\node[circle,fill=white,draw=popink,line width=1.3pt,minimum size=1.6cm,inner sep=0]{\popdisplay\fontsize{24}{24}\selectfont\color{poppurple}+};%
      \hspace{8pt}\begin{minipage}[c]{0.6\linewidth}
        {\popdisplay\fontsize{11}{12.5}\selectfont\color{white}\MakeUppercase{Passe à OnBuch+}}\par\vspace{4pt}
        {\raggedright\popsemi\fontsize{8.4}{11}\selectfont\color{white}Tous les cours signés, Tuteur IA illimité, fiches PDF — onglet OnBuch+ de l'app\par}
      \end{minipage}
    \end{tcolorbox}
  \end{minipage}
  \vspace{8pt}
  \popcontenu
  \vfill
  \signatureonbuch
  \restoregeometry
  \newpage
}

% Rangée « Ce cours contient » (page de garde)
\newcommand{\poptile}[3]{% #1 couleur #2 gros texte #3 légende
  \begin{tcolorbox}[enhanced,colback=#1,colframe=popink,boxrule=2pt,arc=9pt,shadow={2.5pt}{-2.5pt}{0pt}{popink},
      left=6pt,right=6pt,top=9pt,bottom=9pt,halign=center,valign=center,height=1.75cm,width=\linewidth,nobeforeafter]
    {\popdisplay\fontsize{12.5}{13}\selectfont\color{white}\MakeUppercase{#2}}\par\vspace{3pt}
    {\centering\popsemi\fontsize{7.8}{9.5}\selectfont\color{white}#3\par}
  \end{tcolorbox}}
\newcommand{\popcontenu}{%
  \par\noindent{\popxboldls\fontsize{7.5}{7.5}\selectfont\color{popmuted}CE COURS SIGNÉ CONTIENT}\par\vspace{7pt}
  \noindent\begin{minipage}[t]{0.235\linewidth}\poptile{popblue}{Cours}{complet, illustré, pas à pas}\end{minipage}\hfill
  \begin{minipage}[t]{0.235\linewidth}\poptile{poporange}{Méthodes}{et exercices résolus}\end{minipage}\hfill
  \begin{minipage}[t]{0.235\linewidth}\poptile{poppink}{Exercices}{type examen + corrigés}\end{minipage}\hfill
  \begin{minipage}[t]{0.235\linewidth}\poptile{popgreen}{Fiche bilan}{l'essentiel à retenir}\end{minipage}\par}

% Sommaire
\newtcolorbox{sommaire}{enhanced,colback=white,colframe=popink,boxrule=2.2pt,arc=10pt,
  drop shadow=popink,left=18pt,right=18pt,top=13pt,bottom=13pt,before skip=6pt,after skip=20pt,
  before upper={\poppill{Sommaire}{poppurple}{white}\par\vspace{9pt}\popsemi\fontsize{10.6}{14.5}\selectfont\color{popink}}}

% Signature de fin de cours — poussée tout en bas de la dernière page
\newcommand{\findecours}{%
  \par\vfill\nopagebreak
  \begin{center}\popsquiggle{poporange}\end{center}\vspace{4pt}
  \signatureonbuch
}
\AtBeginDocument{\def\llbracket{\mathopen{[\mkern-3.5mu[}}\def\rrbracket{\mathclose{]\mkern-3.5mu]}}} % intervalles d'entiers (glyphe ⟦ absent de la police)
% Glyphes absents de Fira Math : redéfinis à partir de glyphes disponibles
\AtBeginDocument{%
  \def\setminus{\mathbin{\backslash}}\let\smallsetminus\setminus
  \def\vdots{\mathord{\vbox{\baselineskip3.2pt\lineskiplimit0pt\kern4pt\hbox{.}\hbox{.}\hbox{.}}}}%
  \def\ddots{\mathinner{\mkern1mu\raise6.5pt\hbox{.}\mkern2mu\raise3.6pt\hbox{.}\mkern2mu\raise0.7pt\hbox{.}\mkern1mu}}%
  \def\triangle{\mathord{\scalebox{0.9}{$\symup{\Delta}$}}}%
  \def\nabla{\mathord{\raisebox{\depth}{\scalebox{1}[-1]{$\symup{\Delta}$}}}}%
  \def\checkmark{\popcheck{popdark}}%
  \def\star{\mathbin{\ast}}\def\bigstar{\mathord{\ast}}%
  \def\diamond{\mathbin{\scalebox{0.6}{$\Diamond$}}}%
  \def\bigcup{\mathop{\mathchoice{\raisebox{-0.3em}{\scalebox{1.7}{$\cup$}}}{\scalebox{1.25}{$\cup$}}{\cup}{\cup}}\displaylimits}%
  \def\bigcap{\mathop{\mathchoice{\raisebox{-0.3em}{\scalebox{1.7}{$\cap$}}}{\scalebox{1.25}{$\cap$}}{\cap}{\cap}}\displaylimits}%
  \def\lesseqqgtr{\mathrel{\lessgtr}}\def\bigoplus{\mathop{\oplus}\displaylimits}%
}
\providecommand{\ssi}{si et seulement si} % abréviation fréquente
\AtBeginDocument{\providecommand{\wideparen}{\overparen}} % arc de cercle

%% FIGFIX-BEGIN v5 — correctifs globaux des figures (docs/onbuch-generation/tools/figfix)
\usepackage{adjustbox}\usepackage{etoolbox}\tcbuselibrary{raster}
% toute figure TikZ/pgfplots plus large que la ligne est réduite à la largeur disponible
\BeforeBeginEnvironment{tikzpicture}{\begin{adjustbox}{max width=\linewidth}}
\AfterEndEnvironment{tikzpicture}{\end{adjustbox}}
% légendes pgfplots lisibles par-dessus les courbes
\pgfplotsset{every axis/.append style={legend style={fill=white,fill opacity=0.92,text opacity=1,draw=black!15,inner sep=3pt}}}
% étiquettes d'arêtes (syntaxe edge["..."]) avec fond blanc
\tikzset{every edge quotes/.append style={fill=white,inner sep=1.2pt}}
%% FIGFIX-END
\begin{document}

\pagegarde{Équations, inéquations, systèmes}{Mathématiques}{Tle Tech}{Mathématiques – classe de Terminale technique}{N01}{5 séances (fiche de progression harmonisée)}{Cette leçon outille l'élève de Terminale technique pour résoudre les équations du second degré, les inéquations des degrés 2 et 3, ainsi que les systèmes linéaires de trois équations à trois inconnues par substitution et par le pivot de Gauss. Ces savoirs, au cœur de l'épreuve de mathématiques, permettent de modéliser et de résoudre des situations concrètes de la vie économique et technique camerounaise.
\vspace{4pt}
\begin{multicols}{2}\small
\begin{itemize}
\item À la fin de cette leçon, je sais résoudre une équation du second degré par le discriminant, la factorisation ou la forme canonique
\item À la fin de cette leçon, je sais déterminer le signe d'un trinôme du second degré et résoudre une inéquation du second degré
\item À la fin de cette leçon, je sais factoriser un polynôme de degré 3 à l'aide d'une racine évidente et dresser son tableau de signes
\item À la fin de cette leçon, je sais résoudre une inéquation de degré 3 et donner l'ensemble des solutions en intervalles
\item À la fin de cette leçon, je sais résoudre un système linéaire 3×3 par la méthode de substitution
\item À la fin de cette leçon, je sais résoudre un système linéaire 3×3 par la méthode du pivot de Gauss
\end{itemize}\end{multicols}}

\begin{sommaire}
\begin{multicols}{2}
\begin{enumerate}
\item Équations du second degré
\item Signe du trinôme et inéquations du second degré
\item Inéquations de degré 3
\item Systèmes 3×3 : méthode de substitution
\item Systèmes 3×3 : pivot de Gauss
\item Modélisation et résolution de problèmes concrets
\item Activité d'intégration
\item Méthodes et exercices résolus
\item Exercices
\item Corrigés des exercices
\item Fiche bilan
\end{enumerate}
\end{multicols}
\end{sommaire}

\begin{objectifs}
\begin{itemize}
\item À la fin de cette leçon, je sais résoudre une équation du second degré par le discriminant, la factorisation ou la forme canonique
\item À la fin de cette leçon, je sais déterminer le signe d'un trinôme du second degré et résoudre une inéquation du second degré
\item À la fin de cette leçon, je sais factoriser un polynôme de degré 3 à l'aide d'une racine évidente et dresser son tableau de signes
\item À la fin de cette leçon, je sais résoudre une inéquation de degré 3 et donner l'ensemble des solutions en intervalles
\item À la fin de cette leçon, je sais résoudre un système linéaire 3×3 par la méthode de substitution
\item À la fin de cette leçon, je sais résoudre un système linéaire 3×3 par la méthode du pivot de Gauss
\item À la fin de cette leçon, je sais traduire une situation concrète de la vie courante en équation, inéquation ou système et interpréter les solutions
\item À la fin de cette leçon, je sais rédiger et justifier une démarche de résolution complète et conclure de façon argumentée
\end{itemize}
\end{objectifs}

\begin{prerequis}
\begin{itemize}
\item Développement, factorisation et identités remarquables (classe de 3e et Première)
\item Résolution d'équations et d'inéquations du premier degré à une inconnue
\item Signe d'un produit et d'un quotient, tableaux de signes (Première)
\item Systèmes de deux équations à deux inconnues (Première)
\item Notion d'intervalles et d'ensemble de solutions
\end{itemize}
\end{prerequis}

\newpage

\coursec{Équations du second degré}

\noindent
\textbf{Situation d'accroche.} Mama Ngo Bell vend au marché Mokolo trois produits : des sacs de manioc, des régimes de plantain et des cartons de tomates. Un client achète 2 sacs de manioc, 1 régime et 3 cartons pour 45\,000~FCFA ; un autre prend 1 sac, 2 régimes et 1 carton pour 30\,000~FCFA ; un troisième paie 25\,000~FCFA pour 3 sacs et 1 carton. Quel est le prix unitaire de chaque produit ? Par ailleurs, un menuisier de Douala veut découper une plaque de contreplaqué rectangulaire dont l'aire doit dépasser 2~m², la longueur étant de 50~cm supérieure à la largeur : comment traduire et résoudre ce problème ? Ces deux situations — l'une menant à un \emph{système 3×3}, l'autre à une \emph{inéquation du second degré} — illustrent la puissance de l'algèbre pour modéliser le concret. Avant d'aborder les systèmes et les inéquations, nous devons maîtriser parfaitement l'outil de base : \textbf{l'équation du second degré}. C'est l'objet de cette première section.

\courssub{Définition et forme canonique du trinôme}

\begin{definition}[Équation du second degré]
Une \cle{équation du second degré} à une inconnue $x$ est une équation qui s'écrit sous la forme canonique :
\[
ax^{2}+bx+c=0 \qquad \text{avec } a,b,c\in\mathbb{R} \text{ et } a\neq 0.
\]
Le coefficient $a$ est le \cle{coefficient directeur}, $b$ le coefficient de $x$ et $c$ le \cle{terme constant}. L'expression $ax^{2}+bx+c$ est un \cle{trinôme du second degré}.
\end{definition}

\begin{attention}[Condition $a\neq 0$]
Si $a=0$, l'équation devient $bx+c=0$ : c'est une \emph{équation du premier degré}, pas du second. Vérifier $a\neq 0$ est la \textbf{première étape} obligatoire de toute résolution.
\end{attention}

Pour résoudre $ax^{2}+bx+c=0$, nous allons transformer le trinôme en une \emph{différence de deux carrés}. C'est la \cle{forme canonique} (ou forme normale).

\begin{experience}[Démonstration guidée : forme canonique]
Partons de $ax^{2}+bx+c$. Factorisons $a$ sur les deux premiers termes :
\begin{align*}
ax^{2}+bx+c &= a\left(x^{2}+\frac{b}{a}x\right)+c.
\end{align*}
Pour faire apparaître un carré parfait $\left(x+\frac{b}{2a}\right)^{2} = x^{2}+\frac{b}{a}x+\frac{b^{2}}{4a^{2}}$, nous ajoutons et retranchons $\frac{b^{2}}{4a^{2}}$ à l'intérieur de la parenthèse :
\begin{align*}
a\left[\left(x^{2}+\frac{b}{a}x+\frac{b^{2}}{4a^{2}}\right)-\frac{b^{2}}{4a^{2}}\right]+c
&= a\left(x+\frac{b}{2a}\right)^{2} - \frac{ab^{2}}{4a^{2}} + c \\
&= a\left(x+\frac{b}{2a}\right)^{2} - \frac{b^{2}}{4a} + c.
\end{align*}
Mettons au même dénominateur le terme constant :
\[
-\frac{b^{2}}{4a}+c = \frac{-b^{2}+4ac}{4a} = -\frac{b^{2}-4ac}{4a}.
\]
On pose \cle{$\Delta = b^{2}-4ac$} (le \cle{discriminant}). On obtient finalement :
\[
\boxed{ax^{2}+bx+c = a\left(x+\frac{b}{2a}\right)^{2} - \frac{\Delta}{4a}}
\]
\end{experience}

\begin{propriete}[Forme canonique]
Tout trinôme $ax^{2}+bx+c$ ($a\neq 0$) s'écrit de manière unique :
\[
ax^{2}+bx+c = a\left(x+\frac{b}{2a}\right)^{2} - \frac{\Delta}{4a}, \quad \text{avec } \Delta = b^{2}-4ac.
\]
Cette forme révèle le \cle{sommet} de la parabole représentative : coordonnées $\left(-\frac{b}{2a}\,;\, -\frac{\Delta}{4a}\right)$.
\end{propriete}

\courssub{Méthode de résolution algorithmique}

\begin{methode}[Résolution de $ax^{2}+bx+c=0$ ($a\neq 0$)]
\begin{enumerate}
  \item \textbf{Identifier} les coefficients $a$, $b$, $c$ et vérifier $a\neq 0$.
  \item \textbf{Calculer} le discriminant $\Delta = b^{2}-4ac$.
  \item \textbf{Discuter} le signe de $\Delta$ et conclure :
        \begin{itemize}
          \item $\Delta > 0$ : deux solutions distinctes $x_{1} = \dfrac{-b-\sqrt{\Delta}}{2a},\; x_{2} = \dfrac{-b+\sqrt{\Delta}}{2a}$.
          \item $\Delta = 0$ : une solution double $x_{0} = -\dfrac{b}{2a}$.
          \item $\Delta < 0$ : aucune solution réelle ($S = \varnothing$).
        \end{itemize}
  \item \textbf{Rédiger} l'ensemble solution $S$.
  \item \textbf{Vérifier} (si possible) avec les relations $S = -\frac{b}{a}$ et $P = \frac{c}{a}$.
\end{enumerate}
\end{methode}

\courssub{Discriminant et nombre de solutions}

La forme canonique permet de résoudre l'équation $ax^{2}+bx+c=0$ par une simple discussion sur le signe de $\Delta$.

\begin{propriete}[Théorème : Solutions selon le signe de $\Delta$]
Soit $ax^{2}+bx+c=0$ ($a\neq 0$) et $\Delta = b^{2}-4ac$.
\begin{itemize}
  \item Si $\Delta > 0$ : l'équation admet \textbf{deux solutions distinctes} :
        \[
        x_{1} = \frac{-b-\sqrt{\Delta}}{2a}, \qquad x_{2} = \frac{-b+\sqrt{\Delta}}{2a}.
        \]
  \item Si $\Delta = 0$ : l'équation admet \textbf{une solution double} (racine double) :
        \[
        x_{0} = -\frac{b}{2a}.
        \]
  \item Si $\Delta < 0$ : l'équation n'admet \textbf{aucune solution réelle} (ensemble solution vide $\varnothing$ dans $\mathbb{R}$).
\end{itemize}
\end{propriete}

\begin{experience}[Démonstration à partir de la forme canonique]
L'équation $ax^{2}+bx+c=0$ est équivalente à :
\begin{align*}
a\left(x+\frac{b}{2a}\right)^{2} - \frac{\Delta}{4a} &= 0 \\
\iff a\left(x+\frac{b}{2a}\right)^{2} &= \frac{\Delta}{4a} \\
\iff \left(x+\frac{b}{2a}\right)^{2} &= \frac{\Delta}{4a^{2}}.
\end{align*}
\begin{itemize}
  \item $\Delta > 0$ : $\frac{\Delta}{4a^{2}} > 0$. On prend racine carrée : $x+\frac{b}{2a} = \pm\frac{\sqrt{\Delta}}{2|a|}$.
        Comme le $\pm$ absorbe le signe de $a$ ($\frac{1}{|a|} = \frac{\pm 1}{a}$), on retrouve $x = \frac{-b\pm\sqrt{\Delta}}{2a}$.
  \item $\Delta = 0$ : $\left(x+\frac{b}{2a}\right)^{2}=0 \Rightarrow x=-\frac{b}{2a}$.
  \item $\Delta < 0$ : $\frac{\Delta}{4a^{2}} < 0$. Un carré ne peut être négatif dans $\mathbb{R}$. Impossible.
\end{itemize}
\end{experience}

\begin{popfigure}
\begin{tikzpicture}[
  decision/.style={diamond, draw=popblue, thick, fill=popblueL, text width=3.5cm, align=center, inner sep=4pt},
  process/.style={rectangle, draw=popgreen, thick, fill=popgreenL, text width=4.2cm, align=center, inner sep=4pt, rounded corners},
  start/.style={ellipse, draw=poporange, thick, fill=poporangeL, text width=3cm, align=center, inner sep=4pt},
  arrow/.style={-Stealth, thick, popdark}
]
\node[start, align=center] (start){Équation $ax^2+bx+c=0$\\($a\neq 0$)};
\node[decision, below=2.5cm of start] (delta) {Calculer $\Delta = b^2-4ac$};
\node[process, below left=2.2cm and 1.8cm of delta, align=center] (pos){$\Delta > 0$\\(2 solutions distinctes)\\[2pt] $x_{1,2}=\frac{-b\pm\sqrt{\Delta}}{2a}$};
\node[process, below=2.2cm of delta, align=center] (nul){$\Delta = 0$\\(1 solution double)\\[2pt] $x_0=-\frac{b}{2a}$};
\node[process, below right=2.2cm and 1.8cm of delta, align=center] (neg){$\Delta < 0$\\(Aucune solution dans $\mathbb{R}$)\\[2pt] $S=\varnothing$};

\draw[arrow] (start) -- (delta);
\draw[arrow] (delta) -| node[near start, above, sloped, popdark] {$\Delta>0$} (pos);
\draw[arrow] (delta) -- node[left, popdark] {$\Delta=0$} (nul);
\draw[arrow] (delta) -| node[near start, above, sloped, popdark] {$\Delta<0$} (neg);
\end{tikzpicture}
\legende{Arbre de décision pour la résolution d'une équation du second degré.}
\end{popfigure}

\courssub{Cas particuliers : équations incomplètes}

Lorsque $b=0$ ou $c=0$, la factorisation est souvent plus rapide que la formule du discriminant.

\begin{aretenir}[Équations incomplètes]
\begin{itemize}
  \item \textbf{Forme $ax^{2}+c=0$ ($b=0$) :} $x^{2} = -\frac{c}{a}$.
        \begin{itemize}
          \item Si $-\frac{c}{a} > 0$ : deux solutions $x = \pm\sqrt{-\frac{c}{a}}$.
          \item Si $-\frac{c}{a} = 0$ : solution double $x=0$.
          \item Si $-\frac{c}{a} < 0$ : aucune solution.
        \end{itemize}
  \item \textbf{Forme $ax^{2}+bx=0$ ($c=0$) :} Factorisation par $x$ : $x(ax+b)=0$.
        \textbf{Deux solutions systématiques :} $x=0$ et $x=-\frac{b}{a}$.
        \textit{(Ne jamais diviser par $x$ : on perdrait la solution $x=0$ !)}
\end{itemize}
\end{aretenir}

\begin{attention}[Pièges classiques]
\begin{enumerate}
  \item \textbf{Oublier $x=0$} dans $ax^{2}+bx=0$ en divisant par $x$.
  \item \textbf{Mauvais signe dans $\Delta$ :} $\Delta = b^{2}-4ac$ (moins $4ac$, pas plus).
  \item \textbf{Ordre des racines :} Par convention $x_{1} \le x_{2}$, donc $x_{1} = \frac{-b-\sqrt{\Delta}}{2a}$ (le $-$ devant la racine).
  \item \textbf{Simplification prématurée :} Dans $\frac{-b\pm\sqrt{\Delta}}{2a}$, on ne simplifie \textbf{que si} le numérateur et le dénominateur ont un facteur commun \emph{entier} (ex: $\frac{-4\pm 2}{6} = \frac{-2\pm 1}{3}$). Attention au signe de $a$ !
\end{enumerate}
\end{attention}

\courssub{Somme et produit des racines (Relations de Viète)}

Lorsque $\Delta \ge 0$, les deux racines (réelles) $x_{1}$ et $x_{2}$ vérifient des relations simples, très utiles pour vérifier un calcul ou reconstruire une équation.

\begin{propriete}[Relations de Viète]
Si $x_{1}$ et $x_{2}$ sont les solutions de $ax^{2}+bx+c=0$ ($a\neq 0$, $\Delta \ge 0$), alors :
\[
\boxed{S = x_{1}+x_{2} = -\frac{b}{a}} \qquad \text{et} \qquad \boxed{P = x_{1}x_{2} = \frac{c}{a}}.
\]
\textbf{Démonstration (exigible au Probatoire) :}
\begin{align*}
x_{1}+x_{2} &= \frac{-b-\sqrt{\Delta}}{2a} + \frac{-b+\sqrt{\Delta}}{2a} = \frac{-2b}{2a} = -\frac{b}{a}. \\
x_{1}x_{2} &= \frac{(-b-\sqrt{\Delta})(-b+\sqrt{\Delta})}{4a^{2}} = \frac{b^{2}-\Delta}{4a^{2}} = \frac{b^{2}-(b^{2}-4ac)}{4a^{2}} = \frac{4ac}{4a^{2}} = \frac{c}{a}.
\end{align*}
\end{propriete}

\begin{exemplebox}[Utilisation de $S$ et $P$]
On sait que $x=2$ est une racine de $3x^{2}-7x+2=0$. Trouver l'autre racine sans résoudre l'équation.
\par\medskip
\textbf{Solution.} $a=3,\ b=-7,\ c=2$. Produit $P = c/a = 2/3$.
Comme $x_{1}=2$, on a $x_{2} = \frac{P}{x_{1}} = \frac{2/3}{2} = \frac{1}{3}$.
Vérification : Somme $S = -b/a = 7/3$. $2 + 1/3 = 7/3$. \checkmark
\end{exemplebox}

\courssub{Équations se ramenant au second degré : les équations bicarrées}

Une \cle{équation bicarrée} a la forme $ax^{4}+bx^{2}+c=0$ ($a\neq 0$). Elle se résout par le changement de variable $X = x^{2}$ ($X \ge 0$).

\begin{methode}[Résolution d'une équation bicarrée $ax^{4}+bx^{2}+c=0$]
\begin{enumerate}
  \item Poser $X = x^{2}$ (donc $X \ge 0$). L'équation devient $aX^{2}+bX+c=0$.
  \item Résoudre cette équation du second degré en $X$ (calculer $\Delta_X$, trouver $X_1, X_2$).
  \item \textbf{Sélectionner} les racines $X_i$ qui sont $\ge 0$ (les racines $<0$ sont rejetées car $X=x^2\ge 0$).
  \item Pour chaque $X_i \ge 0$ retenu, résoudre $x^{2}=X_i$ :
        \begin{itemize}
          \item Si $X_i > 0$ : deux valeurs $x = \pm\sqrt{X_i}$.
          \item Si $X_i = 0$ : une valeur $x=0$.
        \end{itemize}
  \item Rassembler les solutions en $x$.
\end{enumerate}
\end{methode}

\begin{exemplebox}[Équation bicarrée]
Résoudre $x^{4}-5x^{2}+4=0$.
\par\medskip
\textbf{Solution.}
$X=x^{2}\ge 0$. Équation en $X$ : $X^{2}-5X+4=0$.
$\Delta_X = 25-16=9$. $\sqrt{\Delta_X}=3$.
$X_1 = \frac{5-3}{2}=1 \ge 0$ (retenu). $X_2 = \frac{5+3}{2}=4 \ge 0$ (retenu).
\begin{itemize}
  \item $X_1=1 \Rightarrow x^{2}=1 \Rightarrow x=\pm 1$.
  \item $X_2=4 \Rightarrow x^{2}=4 \Rightarrow x=\pm 2$.
\end{itemize}
Solutions : $\boxed{S = \{-2\,;\,-1\,;\,1\,;\,2\}}$.
\end{exemplebox}

\courssub{Exemple complet résolu et illustration graphique}

\begin{exoresolu}[Résolution complète de $2x^{2}-5x+3=0$]
\textbf{Énoncé.} Résoudre dans $\mathbb{R}$ l'équation $2x^{2}-5x+3=0$. Représenter graphiquement la fonction $f(x)=2x^{2}-5x+3$ et situer les racines.
\tcblower
\textbf{Solution.}
\begin{enumerate}
  \item \textbf{Identification :} $a=2\ (\neq 0),\ b=-5,\ c=3$.
  \item \textbf{Discriminant :} $\Delta = b^{2}-4ac = (-5)^{2}-4\times 2 \times 3 = 25-24 = 1$.
        $\Delta = 1 > 0 \Rightarrow$ \textbf{deux solutions distinctes}.
  \item \textbf{Racines carrées :} $\sqrt{\Delta} = 1$.
  \item \textbf{Formules :}
        \[
        x_{1} = \frac{-b-\sqrt{\Delta}}{2a} = \frac{5-1}{4} = \frac{4}{4} = 1,
        \]
        \[
        x_{2} = \frac{-b+\sqrt{\Delta}}{2a} = \frac{5+1}{4} = \frac{6}{4} = \frac{3}{2} = 1,5.
        \]
  \item \textbf{Vérification (S et P) :} $S = 1+1,5 = 2,5 = -\frac{-5}{2} = \frac{5}{2}$. $P = 1 \times 1,5 = 1,5 = \frac{3}{2}$. \checkmark
  \item \textbf{Ensemble solution :} $\boxed{S = \left\{1\,;\,\frac{3}{2}\right\}}$.
\end{enumerate}
\end{exoresolu}

\begin{popfigure}
\begin{tikzpicture}
\begin{axis}[
    width=\linewidth, height=8cm,
    axis lines=middle,
    xlabel=$x$, ylabel=$y$,
    xlabel style={below right}, ylabel style={above left},
    xtick={0,1,1.5,2,3}, ytick={-1,0,1,2,3,4,5},
    xticklabels={$0$,$1$,$1{,}5$,$2$,$3$},
    domain=-0.5:3.5, samples=200,
    xmin=-0.5, xmax=3.5, ymin=-1.2, ymax=5.5,
    grid=major, grid style={popline, dashed},
    popcurve/.style={popblue, thick},
    clip=false
]
\addplot[popcurve] {2*x^2 - 5*x + 3};
% Racines
\addplot[only marks, mark=*, mark size=2.5pt, poporange] coordinates {(1,0) (1.5,0)};
\node[poporange, anchor=north] at (axis cs:1,0) {$x_1=1$};
\node[poporange, anchor=north] at (axis cs:1.5,0) {$x_2=1,5$};
% Sommet
\addplot[only marks, mark=*, mark size=2.5pt, poppurple] coordinates {(1.25,-0.125)};
\node[poppurple, anchor=north west] at (axis cs:1.25,-0.125) {$S(1,25;-0,125)$};
% Axe x pointillés vers racines
\draw[popmuted, dashed] (axis cs:1,0) -- (axis cs:1,-1);
\draw[popmuted, dashed] (axis cs:1.5,0) -- (axis cs:1.5,-1);
\end{axis}
\end{tikzpicture}
\legende{Parabole $y=2x^2-5x+3$ : deux intersections avec l'axe des abscisses aux racines $x_1=1$ et $x_2=1,5$. Le sommet $S$ a pour abscisse $-\frac{b}{2a}=1,25$.}
\end{popfigure}

\courssub{Exercices d'application}

\begin{exercice}[Entraînement ciblé]
\begin{enumerate}
  \item Résoudre dans $\mathbb{R}$ :
        \begin{itemize}
          \item[(a)] $3x^{2}-12=0$
          \item[(b)] $x^{2}-6x+9=0$
          \item[(c)] $2x^{2}+5x-3=0$
          \item[(d)] $-x^{2}+4x-5=0$
        \end{itemize}
  \item Sans résoudre l'équation $4x^{2}-12x+5=0$, déterminer la somme $S$ et le produit $P$ de ses racines. Si l'une des racines est $\frac{1}{2}$, trouver l'autre.
  \item Résoudre les équations bicarrées :
        \begin{itemize}
          \item[(a)] $x^{4}-13x^{2}+36=0$
          \item[(b)] $2x^{4}+3x^{2}-2=0$
        \end{itemize}
  \item \textbf{Problème concret.} Un technicien en froid industriel doit dimensionner un échangeur. La surface d'échange $S$ (en m²) en fonction du débit $x$ (en m³/h) est modélisée par $S(x)=0,5x^{2}-4x+10$. Pour quel(s) débit(s) la surface d'échange vaut-elle exactement $12$~m² ?
\end{enumerate}
\end{exercice}

\begin{corrige}[Exercice n°1]
\begin{enumerate}
  \item (a) $3x^{2}=12 \Rightarrow x^{2}=4 \Rightarrow S=\{-2;2\}$. \quad (b) $\Delta=36-36=0 \Rightarrow x_0=3 \Rightarrow S=\{3\}$. \quad (c) $\Delta=25+24=49, \sqrt{\Delta}=7 \Rightarrow x_1=-3, x_2=0,5 \Rightarrow S=\{-3;0,5\}$. \quad (d) $\Delta=16-20=-4<0 \Rightarrow S=\varnothing$.
  \item $a=4,b=-12,c=5$. $S=-b/a=3$. $P=c/a=5/4=1,25$. Autre racine $= S - 1/2 = 3 - 0,5 = 2,5$ (ou $P / 0,5 = 2,5$). Vérification : $4(2,5)^2-12(2,5)+5=25-30+5=0$.
  \item (a) $X=x^2$. $X^2-13X+36=0 \Rightarrow \Delta=169-144=25$. $X_1=4, X_2=9$ (tous deux $\ge 0$). $x=\pm 2, \pm 3 \Rightarrow S=\{-3;-2;2;3\}$. \quad (b) $2X^2+3X-2=0 \Rightarrow \Delta=9+16=25$. $X_1=\frac{-3-5}{4}=-2$ (rejeté), $X_2=\frac{-3+5}{4}=0,5$ (retenu). $x=\pm\sqrt{0,5}=\pm\frac{\sqrt{2}}{2} \Rightarrow S=\{-\frac{\sqrt{2}}{2};\frac{\sqrt{2}}{2}\}$.
  \item $0,5x^2-4x+10=12 \iff 0,5x^2-4x-2=0 \iff x^2-8x-4=0$ (×2). $\Delta=64+16=80$. $\sqrt{\Delta}=4\sqrt{5}$. $x_1=4-2\sqrt{5}\approx -0,47$ (rejeté, débit négatif), $x_2=4+2\sqrt{5}\approx 8,47$ m³/h. \textbf{Réponse :} Débit $x = 4+2\sqrt{5}$ m³/h (environ 8,47 m³/h).
\end{enumerate}
\end{corrige}

\begin{savaistu}[Histoire et métier]
Le terme « discriminant » vient du latin \emph{discriminare} (séparer, distinguer) : il « discrimine » le nombre de racines. La méthode du « complément du carré » (forme canonique) était déjà connue des mathématiciens babyloniens (vers 1800 av. J.-C.) pour résoudre des problèmes de surfaces de champs. En technique, l'équation du second degré modélise les \emph{trajectoires paraboliques} (balistique, jets d'eau, câbles de pont suspendu approximés), les \emph{lois de puissance} (pertes de charge $\propto Q^2$), et l'optimisation (surface maximale pour un périmètre donné).
\end{savaistu}

\coursec{Signe du trinôme et inéquations du second degré}

\courssub{Rappel : signe d'une expression affine et tableau de signes}

Avant d'aborder le trinôme du second degré, rappelons un outil fondamental vu en classe de Première : l'étude du signe d'une expression affine $ax+b$ ($a \neq 0$). L'équation $ax+b=0$ a pour unique solution $x_0 = -\dfrac{b}{a}$.

\begin{propriete}[Signe d'une expression affine]
Soit $f(x)=ax+b$ avec $a \neq 0$ et $x_0 = -\dfrac{b}{a}$.
\begin{itemize}
    \item Si $a > 0$ : $f(x) < 0$ pour $x < x_0$, $f(x_0)=0$, $f(x) > 0$ pour $x > x_0$. La fonction est \textbf{strictement croissante}.
    \item Si $a < 0$ : $f(x) > 0$ pour $x < x_0$, $f(x_0)=0$, $f(x) < 0$ pour $x > x_0$. La fonction est \textbf{strictement décroissante}.
\end{itemize}
Autrement dit, l'expression affine est \textbf{du signe de $a$ après sa racine}.
\end{propriete}

Pour visualiser le signe d'une expression produit ou quotient, on utilise un \textbf{tableau de signes}. C'est la méthode universelle pour résoudre les inéquations produits ou quotients.

\begin{exemplebox}[Tableau de signes d'un produit d'expressions affines]
Étudions le signe de $P(x) = (2x-1)(x+3)$.
\begin{itemize}
    \item $2x-1 = 0 \iff x = \dfrac{1}{2}$ ; ce facteur est du signe de $a=2$ (positif) \textbf{après} $\dfrac{1}{2}$.
    \item $x+3 = 0 \iff x = -3$ ; ce facteur est positif \textbf{après} $-3$.
\end{itemize}
On range les racines dans l'ordre croissant : $-3 < \dfrac{1}{2}$.
\begin{center}
\begin{tabular}{|c|ccccccc|}\hline
$x$ & $-\infty$ & & $-3$ & & $\frac{1}{2}$ & & $+\infty$ \\ \hline
$2x-1$ & & $-$ & $-$ & $-$ & $0$ & $+$ & \\
$x+3$ & & $-$ & $0$ & $+$ & $+$ & $+$ & \\ \hline
$P(x)$ & & $+$ & $0$ & $-$ & $0$ & $+$ & \\ \hline
\end{tabular}
\end{center}
On en déduit : $P(x) \geq 0 \iff x \in ]-\infty; -3] \cup \left[\dfrac{1}{2}; +\infty\right[$.
\end{exemplebox}

\begin{attention}[Rédaction d'un tableau de signes]
Dans un tableau de signes, les \textbf{signes} ($+$ ou $-$) se placent dans les cases correspondant aux \textbf{intervalles}, et les \textbf{zéros} dans les cases correspondant aux \textbf{racines}. On ne met jamais de flèches dans un tableau de signes : les flèches ($\nearrow$, $\searrow$) sont réservées aux tableaux de \textbf{variations}.
\end{attention}

\courssub{Théorème : signe du trinôme du second degré}

Soit le trinôme $T(x) = ax^2+bx+c$ avec $a \neq 0$ et son discriminant $\Delta = b^2-4ac$. Lorsque $\Delta > 0$, le trinôme admet deux racines réelles distinctes :
\[ x_1 = \frac{-b-\sqrt{\Delta}}{2a}, \quad x_2 = \frac{-b+\sqrt{\Delta}}{2a}, \]
et sa forme factorisée est $T(x) = a(x-x_1)(x-x_2)$.

\begin{experience}[Démonstration guidée du théorème du signe (cas $\Delta > 0$)]
Partons de la forme factorisée $T(x) = a(x-x_1)(x-x_2)$ en supposant $x_1 < x_2$.
\begin{enumerate}
    \item \textbf{Pour $x < x_1$ :} on a $x-x_1 < 0$ et $x-x_2 < 0$. Le produit $(x-x_1)(x-x_2)$ est \textbf{positif} (règle des signes : $- \times - = +$). Donc $T(x)$ est du \textbf{signe de $a$}.
    \item \textbf{Pour $x = x_1$ ou $x = x_2$ :} l'un des facteurs est nul, donc $T(x)=0$.
    \item \textbf{Pour $x_1 < x < x_2$ :} on a $x-x_1 > 0$ et $x-x_2 < 0$. Le produit est \textbf{négatif} ($+ \times - = -$). Donc $T(x)$ est du \textbf{signe opposé à celui de $a$}.
    \item \textbf{Pour $x > x_2$ :} les deux facteurs sont positifs, le produit est positif. $T(x)$ est du \textbf{signe de $a$}.
\end{enumerate}
\emph{Conclusion :} le trinôme est du signe de $a$ \textbf{à l'extérieur} des racines et du signe opposé \textbf{entre} les racines.
\end{experience}

\begin{propriete}[Théorème du signe du trinôme — Exigible au Probatoire]
Soit $T(x)=ax^2+bx+c$ ($a \neq 0$), $\Delta = b^2-4ac$.
\begin{enumerate}
    \item \textbf{Si $\Delta > 0$} (deux racines distinctes $x_1 < x_2$) :
    \begin{itemize}
        \item $T(x)$ est du \textbf{signe de $a$} pour $x \in ]-\infty; x_1[ \cup ]x_2; +\infty[$ (à l'extérieur des racines) ;
        \item $T(x)$ est du \textbf{signe opposé à celui de $a$} pour $x \in ]x_1; x_2[$ (entre les racines) ;
        \item $T(x)=0$ pour $x=x_1$ ou $x=x_2$.
    \end{itemize}
    \item \textbf{Si $\Delta = 0$} (racine double $x_0 = -\dfrac{b}{2a}$) : $T(x) = a(x-x_0)^2$. Le trinôme est du \textbf{signe de $a$} pour tout $x \neq x_0$ et $T(x_0)=0$. Il ne change pas de signe.
    \item \textbf{Si $\Delta < 0$} (pas de racine réelle) : $T(x)$ est \textbf{strictement du signe de $a$} pour tout $x \in \mathbb{R}$ (toujours positif si $a>0$, toujours négatif si $a<0$).
\end{enumerate}
\end{propriete}

\begin{popfigure}
\begin{tikzpicture}
\begin{axis}[
    width=0.32\linewidth, height=5cm,
    axis lines=middle, xlabel=$x$, ylabel=$y$,
    xmin=-3, xmax=3, ymin=-2, ymax=3,
    xtick={-1,1}, xticklabels={$x_1$,$x_2$}, ytick=\empty,
    tick label style={font=\footnotesize},
    title={$\Delta > 0,\ a>0$}, title style={font=\small},
]
\addplot[popcurve,domain=-2.2:2.2,samples=100] {x^2-1};
\addplot[only marks, mark=*, poporange] coordinates {(-1,0) (1,0)};
\node[popblue, font=\footnotesize] at (axis cs:-2,2.2) {$+$};
\node[popblue, font=\footnotesize] at (axis cs:2,2.2) {$+$};
\node[poporange, font=\footnotesize] at (axis cs:0,-1.4) {$-$};
\end{axis}
\end{tikzpicture}
\legende{Cas $\Delta>0$, $a>0$ : la parabole coupe l'axe des abscisses en $x_1$ et $x_2$ ; le trinôme est positif à l'extérieur, négatif entre les racines.}
\end{popfigure}

\begin{popfigure}
\begin{tikzpicture}
\begin{axis}[
    width=0.32\linewidth, height=5cm,
    axis lines=middle, xlabel=$x$, ylabel=$y$,
    xmin=-3, xmax=3, ymin=-0.5, ymax=3,
    xtick={0}, xticklabels={$x_0$}, ytick=\empty,
    tick label style={font=\footnotesize},
    title={$\Delta = 0,\ a>0$}, title style={font=\small},
]
\addplot[popcurve,domain=-1.7:1.7,samples=100] {x^2};
\addplot[only marks, mark=*, poporange] coordinates {(0,0)};
\node[popblue, font=\footnotesize] at (axis cs:-1.4,2.2) {$+$};
\node[popblue, font=\footnotesize] at (axis cs:1.4,2.2) {$+$};
\end{axis}
\end{tikzpicture}
\legende{Cas $\Delta=0$, $a>0$ : la parabole est tangente à l'axe des abscisses en $x_0$ ; le trinôme reste positif et s'annule en $x_0$.}
\end{popfigure}

\begin{popfigure}
\begin{tikzpicture}
\begin{axis}[
    width=0.32\linewidth, height=5cm,
    axis lines=middle, xlabel=$x$, ylabel=$y$,
    xmin=-3, xmax=3, ymin=-0.5, ymax=3,
    xtick=\empty, ytick=\empty,
    title={$\Delta < 0,\ a>0$}, title style={font=\small},
]
\addplot[popcurve,domain=-1.4:1.4,samples=100] {x^2+1};
\node[popblue, font=\footnotesize] at (axis cs:0,2.6) {$+$ partout};
\end{axis}
\end{tikzpicture}
\legende{Cas $\Delta<0$, $a>0$ : la parabole ne rencontre pas l'axe des abscisses ; le trinôme est strictement positif pour tout $x$ réel.}
\end{popfigure}

\begin{aretenir}[Règle mnémotechnique : « signe de $a$ à l'extérieur des racines »]
Pour $a>0$ :
\begin{itemize}
    \item \textbf{$\Delta > 0$ :} $+$ \ldots $x_1$ \ldots $-$ \ldots $x_2$ \ldots $+$ ;
    \item \textbf{$\Delta = 0$ :} $+$ \ldots $x_0$ \ldots $+$ (le trinôme ne change pas de signe) ;
    \item \textbf{$\Delta < 0$ :} $+$ partout.
\end{itemize}
Si $a<0$, on inverse tous les signes.
\end{aretenir}

\courssub{Résolution des inéquations du second degré}

Une inéquation du second degré s'écrit $ax^2+bx+c > 0$, $\geq 0$, $< 0$ ou $\leq 0$. Sa résolution suit une méthode systématique en quatre étapes.

\begin{methode}[Résolution d'une inéquation du second degré — 4 étapes]
\begin{enumerate}
    \item \textbf{Calculer le discriminant} $\Delta = b^2-4ac$.
    \item \textbf{Déterminer les racines} (si $\Delta \geq 0$) : $x_1 = \dfrac{-b-\sqrt{\Delta}}{2a}$, $x_2 = \dfrac{-b+\sqrt{\Delta}}{2a}$, puis les \textbf{ranger dans l'ordre croissant}.
    \item \textbf{Appliquer la règle du signe du trinôme} (ou construire un tableau de signes) :
    \begin{itemize}
        \item $\Delta > 0$ : signe de $a$ à l'extérieur de $[x_1; x_2]$, signe opposé entre les racines ;
        \item $\Delta = 0$ : signe de $a$ partout, nul en $x_0$ ;
        \item $\Delta < 0$ : signe de $a$ partout (strictement).
    \end{itemize}
    \item \textbf{Conclure par un ensemble de solutions écrit en intervalles}, en respectant le caractère strict ($>$, $<$ : bornes ouvertes) ou large ($\geq$, $\leq$ : bornes fermées) de l'inégalité.
\end{enumerate}
\end{methode}

\begin{exoresolu}[Résolution de $-x^2+4x-3 \geq 0$]
Énoncé : résoudre dans $\mathbb{R}$ l'inéquation $-x^2+4x-3 \geq 0$.
\tcblower
Solution détaillée :
\begin{enumerate}
    \item \textbf{Identification :} $a=-1$, $b=4$, $c=-3$. Ici $a < 0$.
    \item \textbf{Discriminant :} $\Delta = 4^2 - 4\times(-1)\times(-3) = 16 - 12 = 4$. Comme $\Delta > 0$, il y a deux racines distinctes.
    \item \textbf{Racines :}
    \[ x_1 = \frac{-4-\sqrt{4}}{2\times(-1)} = \frac{-4-2}{-2} = 3, \qquad x_2 = \frac{-4+\sqrt{4}}{2\times(-1)} = \frac{-4+2}{-2} = 1. \]
    On range dans l'ordre croissant : $1 < 3$.
    \item \textbf{Signe du trinôme :} $a=-1 < 0$. Le trinôme est donc :
    \begin{itemize}
        \item \textbf{négatif} (signe de $a$) à l'extérieur de $[1;3]$ ;
        \item \textbf{positif} (signe opposé) entre $1$ et $3$ ;
        \item nul en $1$ et en $3$.
    \end{itemize}
    \item \textbf{Conclusion :} on cherche les valeurs pour lesquelles le trinôme est positif ou nul. L'ensemble des solutions est l'intervalle \textbf{fermé} entre les racines :
    \[ S = [1; 3]. \]
\end{enumerate}
\textbf{Vérification graphique :} la parabole d'équation $y=-x^2+4x-3$ est tournée vers le bas ($a<0$) et coupe l'axe des abscisses en $1$ et $3$. Elle est au-dessus de l'axe \textbf{entre} $1$ et $3$, ce qui confirme le résultat.
\end{exoresolu}

\begin{popfigure}
\begin{tikzpicture}
\begin{axis}[
    width=0.7\linewidth, height=6cm,
    axis lines=middle, xlabel=$x$, ylabel=$y$,
    xmin=-0.5, xmax=4.5, ymin=-3, ymax=1.5,
    xtick={1,3}, ytick={-3,-2,-1,1},
    tick label style={font=\small},
]
\addplot[popcurve, thick, domain=-0.5:4.5, samples=100] {-x^2+4*x-3};
\addplot[only marks, mark=*, poporange, mark size=2.5pt] coordinates {(1,0) (3,0)};
\addplot[popgreen!50!black, domain=1:3, samples=50, fill=popgreen!20, fill opacity=0.3] {-x^2+4*x-3} \closedcycle;
\node[popgreen!50!black, anchor=south] at (axis cs:2,0.6) {$y \geq 0$};
\node[poporange, anchor=north] at (axis cs:2,-0.4) {$S=[1;3]$};
\draw[dashed, popmuted] (axis cs:1,-3) -- (axis cs:1,0) (axis cs:3,-3) -- (axis cs:3,0);
\end{axis}
\end{tikzpicture}
\legende{Représentation graphique de la solution de $-x^2+4x-3 \geq 0$. La zone colorée correspond à $y \geq 0$.}
\end{popfigure}

\courssub{Inéquations à quotient ou produit : tableau de signes complet}

Lorsque l'inéquation met en jeu un quotient $\dfrac{T_1(x)}{T_2(x)} \geq 0$ ou un produit, la règle du trinôme seule ne suffit plus. Il faut construire un \textbf{tableau de signes global} en rangeant dans l'ordre croissant les zéros de tous les facteurs (numérateur et dénominateur).

\begin{attention}[Pièges majeurs — Exigences de rédaction au Probatoire]
\begin{enumerate}
    \item \textbf{Valeurs interdites :} pour un quotient, les racines du \textbf{dénominateur} sont \textbf{exclues} de l'ensemble de définition. On les signale par une double barre dans le tableau de signes. Elles ne font \textbf{jamais} partie de la solution, même pour une inégalité large ($\geq$, $\leq$).
    \item \textbf{Écriture en intervalles :} la réponse finale \textbf{doit} être donnée sous forme d'intervalles (réunion d'intervalles si nécessaire). Écrire « $x \leq 1$ ou $x \geq 3$ » est une rédaction \textbf{incomplète}. On écrit $S = ]-\infty; 1] \cup [3; +\infty[$.
    \item \textbf{Signe de $a$ :} ne pas confondre le signe du coefficient $a$ et le signe de l'expression. Relire la règle : « signe de $a$ à l'extérieur des racines ».
    \item \textbf{Ordre des racines :} toujours ranger les racines de la plus petite à la plus grande avant de remplir le tableau.
    \item \textbf{Multiplication par un nombre négatif :} multiplier ou diviser les deux membres d'une inéquation par un nombre \textbf{strictement négatif} change le sens de l'inégalité.
\end{enumerate}
\end{attention}

\begin{exemplebox}[Inéquation quotient : $\dfrac{x^2-4}{x-1} \geq 0$]
\textbf{Résolution :}
\begin{enumerate}
    \item \textbf{Ensemble de définition :} le dénominateur ne doit pas s'annuler : $x-1 \neq 0 \iff x \neq 1$. On travaille sur $\mathcal{D} = \mathbb{R} \setminus \{1\}$.
    \item \textbf{Zéros du numérateur :} $x^2-4=0 \iff (x-2)(x+2)=0 \iff x=-2$ ou $x=2$.
    \item \textbf{Signe de chaque expression :} $x^2-4$ est un trinôme avec $a=1>0$ : positif à l'extérieur de $[-2;2]$, négatif entre $-2$ et $2$. $x-1$ est positif après $1$.
    \item \textbf{Tableau de signes :} on range $-2$, $1$ (valeur interdite), $2$ dans l'ordre croissant.
\end{enumerate}
\begin{center}
\begin{tabular}{|c|ccccccccc|}\hline
$x$ & $-\infty$ & & $-2$ & & $1$ & & $2$ & & $+\infty$ \\ \hline
$x^2-4$ & & $+$ & $0$ & $-$ & $-$ & $-$ & $0$ & $+$ & \\
$x-1$ & & $-$ & $-$ & $-$ & $0$ & $+$ & $+$ & $+$ & \\ \hline
$\dfrac{x^2-4}{x-1}$ & & $-$ & $0$ & $+$ & $||$ & $-$ & $0$ & $+$ & \\ \hline
\end{tabular}
\end{center}
\textbf{Solution :} on cherche les valeurs pour lesquelles le quotient est positif ou nul. Les zéros $-2$ et $2$ sont acceptés (ils annulent le numérateur). La valeur $1$ est exclue (elle annule le dénominateur) : la borne est ouverte en $1$.
\[ S = [-2; 1[ \cup [2; +\infty[. \]
\end{exemplebox}

\courssub{Lien avec la position de la parabole}

L'étude du signe du trinôme $T(x)=ax^2+bx+c$ est \textbf{exactement équivalente} à l'étude de la position de la parabole $\mathcal{P}$ d'équation $y=ax^2+bx+c$ par rapport à l'axe des abscisses.

\begin{propriete}[Parabole et signe du trinôme]
\begin{itemize}
    \item $T(x) > 0 \iff$ la parabole est \textbf{strictement au-dessus} de l'axe des abscisses ;
    \item $T(x) = 0 \iff$ la parabole \textbf{coupe} l'axe des abscisses (cas $\Delta>0$, points d'abscisses $x_1$ et $x_2$) ou lui est \textbf{tangente} (cas $\Delta=0$, point d'abscisse $x_0$) ;
    \item $T(x) < 0 \iff$ la parabole est \textbf{strictement en dessous} de l'axe des abscisses.
\end{itemize}
\end{propriete}

\begin{savaistu}[Visualisation géométrique]
La règle « signe de $a$ à l'extérieur des racines » se voit immédiatement sur le graphique :
\begin{itemize}
    \item Si $a > 0$ (parabole tournée vers le haut) : les branches montent vers $+\infty$. Le trinôme est positif \textbf{à l'extérieur} des points d'intersection avec l'axe des abscisses, et négatif \textbf{entre} ces points (le fond de la parabole est en bas).
    \item Si $a < 0$ (parabole tournée vers le bas) : les branches descendent vers $-\infty$. Le trinôme est négatif \textbf{à l'extérieur}, et positif \textbf{entre} les points d'intersection (le sommet est en haut).
\end{itemize}
Il est donc \textbf{fortement recommandé} d'esquisser la parabole (sens d'ouverture, racines) avant de rédiger la solution finale : cela évite les erreurs de sens d'inégalité.
\end{savaistu}

\begin{exoresolu}[Problème concret : rentabilité d'une production]
Énoncé : une entreprise de menuiserie modélise son bénéfice mensuel (en milliers de FCFA) par la fonction $B(x) = -2x^2 + 40x - 150$, où $x$ est le nombre de meubles produits (en dizaines). Pour quels volumes de production le bénéfice est-il strictement positif ?
\tcblower
Solution détaillée :
\begin{enumerate}
    \item On cherche $B(x) > 0 \iff -2x^2+40x-150 > 0$.
    \item On divise les deux membres par $-2$ (nombre négatif : \textbf{le sens de l'inégalité change}) :
    \[ x^2 - 20x + 75 < 0. \]
    \item Discriminant : $\Delta = (-20)^2 - 4\times 1 \times 75 = 400 - 300 = 100$, donc $\sqrt{\Delta}=10$.
    \item Racines : $x_1 = \dfrac{20-10}{2}=5$ et $x_2 = \dfrac{20+10}{2}=15$.
    \item Le trinôme $x^2-20x+75$ a un coefficient $a=1>0$ : il est du signe de $a$ (positif) à l'extérieur des racines et négatif entre les racines.
    \item On cherche les valeurs pour lesquelles il est strictement négatif : c'est \textbf{entre} les racines, bornes exclues : $5 < x < 15$.
    \item \textbf{Interprétation :} $x$ est exprimé en dizaines de meubles. L'entreprise réalise un bénéfice strictement positif pour une production comprise entre \textbf{50 et 150 meubles par mois} (bornes exclues).
\end{enumerate}
\end{exoresolu}

\begin{exercice}[Entraînement — Inéquations du second degré]
\begin{enumerate}
    \item Résoudre dans $\mathbb{R}$ :
    \begin{enumerate}
        \item $3x^2 - 5x - 2 \geq 0$
        \item $-x^2 + 2x + 3 < 0$
        \item $4x^2 - 12x + 9 \leq 0$
        \item $x^2 + x + 1 > 0$
    \end{enumerate}
    \item Résoudre dans $\mathbb{R}$ l'inéquation $\dfrac{x^2-9}{x+2} \leq 0$.
    \item Un projectile est lancé verticalement. Sa hauteur (en mètres) au bout de $t$ secondes est $h(t) = -5t^2 + 20t + 15$. Pendant quel intervalle de temps le projectile est-il à une hauteur strictement supérieure à 20 mètres ?
\end{enumerate}
\end{exercice}

\begin{corrige}[Exercice n°1]
\begin{enumerate}
    \item \textbf{a)} $3x^2-5x-2 \geq 0$. $a=3>0$ ; $\Delta=(-5)^2-4\times 3\times(-2)=25+24=49$ ; $\sqrt{\Delta}=7$. Racines : $x_1=\dfrac{5-7}{6}=-\dfrac{1}{3}$ et $x_2=\dfrac{5+7}{6}=2$. Le trinôme est positif à l'extérieur des racines. $S = \left]-\infty; -\dfrac{1}{3}\right] \cup [2; +\infty[$.

    \textbf{b)} $-x^2+2x+3 < 0$. On multiplie par $-1$ (le sens change) : $x^2-2x-3 > 0$. $a=1>0$ ; $\Delta=4+12=16$. Racines : $x_1=\dfrac{2-4}{2}=-1$, $x_2=\dfrac{2+4}{2}=3$. Le trinôme est strictement positif à l'extérieur des racines. $S = ]-\infty; -1[ \cup ]3; +\infty[$.

    \textbf{c)} $4x^2-12x+9 \leq 0$. On reconnaît $(2x-3)^2 \leq 0$ ($\Delta = 144-144=0$, racine double $x_0=\dfrac{3}{2}$). Un carré est toujours positif ou nul ; il ne peut être négatif. Il ne s'annule qu'en $x_0$. $S = \left\{\dfrac{3}{2}\right\}$.

    \textbf{d)} $x^2+x+1 > 0$. $a=1>0$ ; $\Delta=1-4=-3<0$. Le trinôme est strictement positif pour tout réel. $S = \mathbb{R}$.

    \item \textbf{Ensemble de définition :} $x+2 \neq 0 \iff x \neq -2$. \textbf{Zéros du numérateur :} $x^2-9=0 \iff x=-3$ ou $x=3$. On range : $-3 < -2 < 3$.
    \begin{center}
    \begin{tabular}{|c|ccccccccc|}\hline
    $x$ & $-\infty$ & & $-3$ & & $-2$ & & $3$ & & $+\infty$ \\ \hline
    $x^2-9$ & & $+$ & $0$ & $-$ & $-$ & $-$ & $0$ & $+$ & \\
    $x+2$ & & $-$ & $-$ & $-$ & $0$ & $+$ & $+$ & $+$ & \\ \hline
    $\dfrac{x^2-9}{x+2}$ & & $-$ & $0$ & $+$ & $||$ & $-$ & $0$ & $+$ & \\ \hline
    \end{tabular}
    \end{center}
    On cherche un quotient négatif ou nul ; $-2$ est exclu (dénominateur), $-3$ et $3$ sont acceptés.
    $S = ]-\infty; -3] \cup ]-2; 3]$.

    \item $h(t) > 20 \iff -5t^2+20t+15 > 20 \iff -5t^2+20t-5 > 0$. On divise par $-5$ (le sens change) : $t^2-4t+1 < 0$.
    $\Delta=16-4=12$, $\sqrt{\Delta}=2\sqrt{3}$. Racines : $t_1=\dfrac{4-2\sqrt{3}}{2}=2-\sqrt{3}\approx 0,27$ et $t_2=2+\sqrt{3}\approx 3,73$. Le trinôme ($a=1>0$) est strictement négatif entre les racines.
    $S = ]2-\sqrt{3}\,;\, 2+\sqrt{3}[$. Le projectile dépasse 20 mètres de hauteur entre environ $0,27$ s et $3,73$ s après le lancer.
\end{enumerate}
\end{corrige}

\coursec{Inéquations de degré 3}

Après avoir étudié le signe des trinômes et la résolution des inéquations du second degré, nous abordons maintenant les polynômes de degré 3. La méthode repose sur la même idée : factoriser pour obtenir un produit de facteurs du premier degré (et éventuellement un trinôme du second degré ne s'annulant pas) puis construire un tableau de signes. La difficulté supplémentaire vient de la recherche d'une première racine et de la gestion des facteurs au carré qui ne changent pas de signe.

\courssub{Polynômes de degré 3 : forme générale et recherche d'une racine évidente}

\begin{definition}[Polynôme de degré 3]
Un polynôme de degré 3 est une fonction de la forme
\[
P(x)=ax^{3}+bx^{2}+cx+d,\qquad a\neq 0,\; a,b,c,d\in\mathbb{R}.
\]
\end{definition}

\begin{propriete}[Théorème du facteur]
Soit $P$ un polynôme de degré 3 et $\alpha\in\mathbb{R}$. Si $P(\alpha)=0$, alors $P(x)$ est divisible par $(x-\alpha)$ ; il existe un polynôme $Q$ de degré 2 tel que
\[
P(x)=(x-\alpha)Q(x).
\]
Cette propriété est admise au Probatoire ; elle se justifie par l'identification des coefficients ou la division euclidienne.
\end{propriete}

\begin{methode}[Recherche d'une racine évidente]
\begin{enumerate}
\item Lister les candidats rationnels : $\pm 1,\pm 2,\pm 3,\pm\frac12,\pm\frac13,\dots$ (diviseurs du terme constant sur diviseurs du coefficient dominant).
\item Évaluer $P$ pour chaque candidat jusqu'à trouver $\alpha$ tel que $P(\alpha)=0$.
\item Une fois $\alpha$ trouvé, factoriser $P(x)=(x-\alpha)(ax^{2}+b'x+c')$ par identification des coefficients ou division euclidienne.
\end{enumerate}
\end{methode}

\begin{exemplebox}[Recherche d'une racine pour $P(x)=x^{3}-3x+2$]
On teste $x=1$ : $P(1)=1-3+2=0$. Donc $1$ est une racine.
On écrit $P(x)=(x-1)(x^{2}+px+q)$. En identifiant :
\[
(x-1)(x^{2}+px+q)=x^{3}+px^{2}+qx-x^{2}-px-q=x^{3}+(p-1)x^{2}+(q-p)x-q.
\]
On compare avec $x^{3}+0x^{2}-3x+2$ :
\[
\begin{cases}
p-1=0 &\Rightarrow p=1,\\
q-p=-3 &\Rightarrow q-1=-3 \Rightarrow q=-2,\\
-q=2 &\Rightarrow q=-2 \text{ (cohérent)}.
\end{cases}
\]
Ainsi $P(x)=(x-1)(x^{2}+x-2)$.
\end{exemplebox}

\courssub{Factorisation complète et résolution de l'équation $P(x)=0$}

Le trinôme $x^{2}+x-2$ se factorise : $\Delta=1+8=9$, racines $\frac{-1\pm3}{2}$ soit $1$ et $-2$. Donc
\[
x^{2}+x-2=(x-1)(x+2).
\]
Finalement
\[
P(x)=(x-1)^{2}(x+2).
\]

\begin{propriete}[Résolution de $P(x)=0$]
Les solutions de $P(x)=0$ sont les racines de chaque facteur :
\begin{itemize}
\item $x-1=0 \Rightarrow x=1$ (racine double),
\item $x+2=0 \Rightarrow x=-2$.
\end{itemize}
L'ensemble solution est $\{-2;1\}$.
\end{propriete}

\courssub{Tableau de signes d'un polynôme de degré 3 factorisé}

Pour résoudre une inéquation, on étudie le signe de chaque facteur sur $\mathbb{R}$.

\begin{attention}[Facteur au carré]
Un facteur de la forme $(x-\alpha)^{2}$ est \textbf{toujours $\ge 0$} (nul en $\alpha$, positif ailleurs). Il \textbf{ne change pas le signe} du produit. L'oubli de cette propriété est une erreur fréquente : on ne doit pas placer un changement de signe au niveau d'une racine double.
\end{attention}

Pour $P(x)=(x+2)(x-1)^{2}$, les valeurs remarquables sont $-2$ et $1$. On construit le tableau de signes suivant (une colonne par valeur remarquable et une colonne par intervalle) :

\begin{popfigure}
\begin{center}
\begin{tabular}{|c|ccccccccc|}\hline
$x$ & $-\infty$ & & $-2$ & & $1$ & & $+\infty$ \\ \hline
$x+2$ & $-$ & $-$ & $0$ & $+$ & $+$ & $+$ & $+$ \\ \hline
$(x-1)^2$ & $+$ & $+$ & $+$ & $+$ & $0$ & $+$ & $+$ \\ \hline
$P(x)$ & $-$ & $-$ & $0$ & $+$ & $0$ & $+$ & $+$ \\ \hline
\end{tabular}
\end{center}
\legende{Tableau de signes de $P(x)=(x+2)(x-1)^2$}
\end{popfigure}

\courssub{Résolution d'inéquations de degré 3}

\begin{methode}[Résoudre une inéquation de degré 3 en 5 étapes]
\begin{enumerate}
\item Mettre l'inéquation sous la forme $P(x)\ge 0$ (ou $>0$, $\le 0$, $<0$).
\item Factoriser complètement $P(x)$ (recherche d'une racine évidente, puis factorisation du trinôme).
\item Repérer les valeurs qui annulent chaque facteur (racines simples et doubles).
\item Construire le tableau de signes en notant le signe de chaque facteur sur chaque intervalle ; rappeler qu'un facteur au carré est toujours $\ge 0$.
\item Lire la solution sur la ligne du produit $P(x)$ et l'exprimer en intervalles (en incluant ou excluant les bornes selon le sens de l'inéquation).
\end{enumerate}
\end{methode}

\begin{exoresolu}[Résoudre $x^{3}-3x+2\ge 0$]
\textbf{Énoncé :} Déterminer l'ensemble des solutions de $x^{3}-3x+2\ge 0$.

\textbf{Solution détaillée :}
\begin{enumerate}
\item L'inéquation est déjà sous la forme $P(x)\ge 0$ avec $P(x)=x^{3}-3x+2$.
\item Factorisation : $P(x)=(x-1)^{2}(x+2)$ (voir exemple précédent).
\item Racines : $-2$ (simple) et $1$ (double).
\item Tableau de signes (cf. figure ci-dessus) :
\begin{itemize}
\item Sur $]-\infty;-2[$ : $x+2<0$, $(x-1)^2>0$ $\Rightarrow$ $P(x)<0$.
\item En $-2$ : $P(-2)=0$.
\item Sur $]-2;1[$ : $x+2>0$, $(x-1)^2>0$ $\Rightarrow$ $P(x)>0$.
\item En $1$ : $P(1)=0$.
\item Sur $]1;+\infty[$ : $x+2>0$, $(x-1)^2>0$ $\Rightarrow$ $P(x)>0$.
\end{itemize}
\item L'inéquation demande $P(x)\ge 0$. On prend les intervalles où $P(x)$ est positif ou nul : $[-2;+\infty[$.
\end{enumerate}
\textbf{Conclusion :} $S=[-2;+\infty[$.
\end{exoresolu}

\begin{popfigure}
\begin{tikzpicture}
\begin{axis}[
    width=\linewidth,
    height=7cm,
    axis lines=middle,
    xlabel=$x$, ylabel=$y$,
    xmin=-3.5, xmax=2.5,
    ymin=-2, ymax=5,
    xtick={-2,1},
    xticklabels={$-2$,$1$},
    ytick={0},
    yticklabels={$0$},
    domain=-3.5:2.5,
    samples=300,
    clip=false,
    enlargelimits=0.05
]
% Zone où P(x) >= 0
\addplot[popgreen!20, domain=-2:2.5, fill=popgreen!20, draw=none] {x^3-3*x+2} \closedcycle;
\addplot[popred!20, domain=-3.5:-2, fill=popred!20, draw=none] {x^3-3*x+2} \closedcycle;
% Courbe
\addplot[popcurve, thick, popblue] {x^3-3*x+2};
% Points racines
\addplot[only marks, mark=*, mark size=2.5, popred] coordinates {(-2,0) (1,0)};
% Droite y=0
\addplot[dashed, popmuted] coordinates {(-3.5,0) (2.5,0)};
\end{axis}
\end{tikzpicture}
\legende{Courbe de $P(x)=x^{3}-3x+2$ ; zone verte : $P(x)\ge 0$}
\end{popfigure}

\courssub{Inéquations avec quotient de polynômes}

Lorsque l'inéquation contient une fraction rationnelle $\frac{P(x)}{Q(x)}\ge 0$, il faut :
\begin{itemize}
\item Déterminer les valeurs interdites (zéros du dénominateur $Q(x)$).
\item Factoriser numérateur et dénominateur.
\item Construire un tableau de signes incluant les facteurs du numérateur et du dénominateur (le dénominateur change de signe à ses zéros simples).
\item Exclure les valeurs interdites de la solution finale (double barre $||$ au tableau).
\end{itemize}

\begin{exemplebox}[Exemple : $\frac{x^{3}-3x+2}{x-3}\ge 0$]
Numérateur : $(x-1)^{2}(x+2)$. Dénominateur : $x-3$ (zéro en $3$, valeur interdite).
Tableau de signes (lignes : $x+2$, $(x-1)^2$, $x-3$, quotient) :
\begin{center}
\begin{tabular}{|c|cccccccccc|}\hline
$x$ & $-\infty$ & & $-2$ & & $1$ & & $3$ & & $+\infty$ \\ \hline
$x+2$ & $-$ & $-$ & $0$ & $+$ & $+$ & $+$ & $+$ & $+$ & $+$ \\ \hline
$(x-1)^2$ & $+$ & $+$ & $+$ & $+$ & $0$ & $+$ & $+$ & $+$ & $+$ \\ \hline
$x-3$ & $-$ & $-$ & $-$ & $-$ & $-$ & $-$ & $0$ & $+$ & $+$ \\ \hline
$\frac{P}{Q}$ & $+$ & $+$ & $0$ & $-$ & $-$ & $0$ & $||$ & $+$ & $+$ \\ \hline
\end{tabular}
\end{center}
L'inéquation est $\ge 0$. Le quotient est positif sur $]-\infty;-2]$ (nul en $-2$) et sur $]3;+\infty[$. Il est nul en $1$ (racine double du numérateur). Il est négatif sur $]-2;1[$ et $]1;3[$.
\textbf{Solution :} $S=\{-2;1\}\cup]3;+\infty[$.
\end{exemplebox}

\courssub{Modélisation : seuils de rentabilité}

\begin{savaistu}[Applications concrètes]
Dans la gestion de production, le bénéfice $B(x)$ en fonction du nombre d'unités $x$ est souvent modélisé par un polynôme de degré 3 (coûts fixes, coûts variables croissants, effets d'échelle). Résoudre $B(x)\ge 0$ permet de déterminer les seuils de rentabilité : en dessous d'un certain volume, l'entreprise perd de l'argent ; au-delà, elle devient bénéficiaire. Les inéquations de degré 3 sont donc des outils essentiels pour l'aide à la décision technique et commerciale.
\end{savaistu}

\begin{exercice}[Entraînement]
Résoudre les inéquations suivantes :
\begin{enumerate}
\item $x^{3}-4x^{2}+x+6\ge 0$ (indice : tester $x=2$).
\item $\frac{x^{3}+2x^{2}-x-2}{x+1}>0$.
\end{enumerate}
\end{exercice}

\begin{corrige}[Exercice n°1]
\begin{enumerate}
\item $P(2)=8-16+2+6=0$ $\Rightarrow$ racine $2$. Division : $P(x)=(x-2)(x^{2}-2x-3)=(x-2)(x-3)(x+1)$. Racines simples : $-1$, $2$, $3$. Tableau de signes $\Rightarrow$ $S=[-1;2]\cup[3;+\infty[$.
\item Numérateur : $x^{3}+2x^{2}-x-2$. $x=-1$ donne $-1+2+1-2=0$ $\Rightarrow$ facteur $(x+1)$. Division : $(x+1)(x^{2}+x-2)=(x+1)(x+2)(x-1)$. Dénominateur $x+1$ (interdit $-1$). Simplification : $\frac{(x+1)(x+2)(x-1)}{x+1}=(x+2)(x-1)$ pour $x\neq-1$. Inéquation $(x+2)(x-1)>0$ $\Rightarrow$ $S=]-\infty;-2[\cup]1;+\infty[$ (en excluant $-1$ déjà hors intervalle).
\end{enumerate}
\end{corrige}

\coursec{Systèmes 3×3 : méthode de substitution}

Après avoir étudié les inéquations de degré 3 et leurs ensembles de solutions, nous abordons maintenant la résolution de systèmes d'équations linéaires à trois inconnues. Cette compétence est fondamentale en technique : elle permet de modéliser des situations où trois grandeurs sont liées par trois relations linéaires (équilibre de forces dans une structure, loi des mailles dans un circuit électrique, mélange de trois composants en chimie ou en gestion de stocks). La méthode de substitution, que nous détaillons ici, est souvent la plus intuitive lorsque l'une des équations se prête aisément à l'isolation d'une inconnue.

\courssub{Définition et notion de solution}

\begin{definition}[Système linéaire 3×3 et solution]
Un \textbf{système linéaire de trois équations à trois inconnues} (notées $x$, $y$, $z$) s'écrit sous la forme canonique :
\[
\begin{cases}
a_1 x + b_1 y + c_1 z = d_1 \\
a_2 x + b_2 y + c_2 z = d_2 \\
a_3 x + b_3 y + c_3 z = d_3
\end{cases}
\]
où les coefficients $a_i, b_i, c_i, d_i \in \mathbb{R}$ ($i=1,2,3$).
Un \textbf{triplet} $(x_0 ; y_0 ; z_0) \in \mathbb{R}^3$ est une \textbf{solution} du système si et seulement si, en remplaçant $x$ par $x_0$, $y$ par $y_0$ et $z$ par $z_0$, les trois égalités sont \textbf{simultanément} vérifiées.
L'ensemble des solutions se note $S$. On a trois cas de figure :
\begin{itemize}
    \item $S = \varnothing$ : système \textbf{impossible} (aucune solution).
    \item $S = \{(x_0;y_0;z_0)\}$ : système \textbf{déterminé} (solution unique).
    \item $S$ infini : système \textbf{indéterminé} (infinité de solutions, souvent une droite ou un plan).
\end{itemize}
\end{definition}

\begin{savaistu}[Choix de la méthode — Probatoire]
La substitution est la méthode la plus naturelle \textbf{quand une inconnue a pour coefficient $1$ (ou $-1$)} dans l'une des équations. Elle évite les fractions dès le début et limite les erreurs de calcul. Si aucun coefficient n'est $\pm 1$, le pivot de Gauss (section suivante) est souvent plus efficace, plus algorithmique et privilégié pour l'automatisation (tableur, programmation). À l'examen, choisissez la méthode la plus rapide selon la forme du système.
\end{savaistu}

\courssub{Principe de la méthode de substitution}

L'idée est de \textbf{réduire progressivement le nombre d'inconnues} :
\begin{enumerate}
    \item \textbf{Isoler} une inconnue (par exemple $z$) dans l'une des trois équations (de préférence celle où son coefficient est $\pm 1$).
    \item \textbf{Substituer} cette expression de $z$ dans les \textbf{deux autres} équations. On obtient un système de \textbf{deux équations à deux inconnues} ($x$ et $y$).
    \item Résoudre ce système $2\times 2$ (par substitution à nouveau, ou par combinaison linéaire) pour trouver $x$ et $y$.
    \item \textbf{Remonter} : calculer $z$ en remplaçant $x$ et $y$ trouvés dans l'expression de l'étape 1.
    \item \textbf{Vérifier} le triplet $(x;y;z)$ dans les \textbf{trois} équations initiales.
\end{enumerate}

\begin{methode}[Résolution par substitution — 5 étapes]
\begin{enumerate}
    \item \textbf{Choisir} l'équation la plus simple et \textbf{exprimer} une inconnue en fonction des deux autres (ex: $z = \dots$).
    \item \textbf{Remplacer} cette expression dans les \textbf{deux autres} équations. On obtient un système $2\times 2$ en $x$ et $y$.
    \item \textbf{Résoudre} ce système $2\times 2$ (substitution ou combinaison) pour obtenir $x$ et $y$.
    \item \textbf{Calculer} la troisième inconnue ($z$) par \textbf{remontée} dans la formule de l'étape 1.
    \item \textbf{Vérifier} le triplet $(x;y;z)$ dans les \textbf{trois} équations du système initial. Écrire la conclusion : $S = \{(x;y;z)\}$.
\end{enumerate}
\end{methode}

\begin{popfigure}
\begin{tikzpicture}[
    node distance=1.2cm and 2.5cm,
    box/.style={draw=popblue, thick, rounded corners, fill=popblueL, minimum width=3.8cm, minimum height=1.1cm, align=center, font=\small},
    arrow/.style={-Stealth, thick, popdark},
    labelarrow/.style={midway, fill=white, inner sep=2pt, font=\footnotesize, text=popdark}
]
% Niveau 1 : Système 3x3
\node[box, align=center] (sys3){Système initial $3\times 3$\\(3 équations, 3 inconnues $x,y,z$)};

% Flèche vers bas
\draw[arrow] (sys3.south) -- ++(0,-0.6) node[labelarrow] {Étape 1 : Isoler $z$ dans Eq(1)} -- ++(0,-0.6) node[box, anchor=north] (exprz) {$z = \alpha x + \beta y + \gamma$};

% Flèches vers les deux autres équations
\draw[arrow] (exprz.east) -- ++(1.5,0) node[labelarrow, above] {Substituer dans Eq(2)} -- ++(1.5,0) node[box, anchor=west] (eq2bis) {Eq(2)' : $a'x+b'y=d'$};
\draw[arrow] (exprz.west) -- ++(-1.5,0) node[labelarrow, above] {Substituer dans Eq(3)} -- ++(-1.5,0) node[box, anchor=east] (eq3bis) {Eq(3)' : $a''x+b''y=d''$};

% Niveau 2 : Système 2x2
\node[box, below=1.5cm of eq2bis.south, xshift=-1.5cm, align=center] (sys2){Système réduit $2\times 2$\\(2 équations, 2 inconnues $x,y$)};

\draw[arrow] (eq2bis.south) -- (sys2.north east) node[labelarrow, right] {};
\draw[arrow] (eq3bis.south) -- (sys2.north west) node[labelarrow, left] {};

% Résolution 2x2
\draw[arrow] (sys2.south) -- ++(0,-0.6) node[labelarrow] {Résoudre (substitution ou combinaison)} -- ++(0,-0.6) node[box, anchor=north] (solxy) {$x = x_0$ ; $y = y_0$};

% Remontée
\draw[arrow] (solxy.south) -- ++(0,-0.6) node[labelarrow] {Remontée : $z = \alpha x_0 + \beta y_0 + \gamma$} -- ++(0,-0.6) node[box, anchor=north] (solz) {$z = z_0$};

% Solution finale
\draw[arrow] (solz.south) -- ++(0,-0.6) node[labelarrow] {Vérification dans les 3 équations} -- ++(0,-0.6) node[box, anchor=north, fill=popgreenL, draw=popgreen] (final) {Solution : $S = \{(x_0;y_0;z_0)\}$};
\end{tikzpicture}
\legende{Réduction en escalier $3\times 3 \to 2\times 2 \to 1\times 1$ puis remontée}
\end{popfigure}

\courssub{Exemple chiffré complet}

\begin{exoresolu}[Résolution par substitution]
Résoudre le système suivant :
\[
(S) \quad \begin{cases}
(1) & x + y + z = 6 \\
(2) & 2x - y + z = 3 \\
(3) & x + 2y - z = 2
\end{cases}
\]
\textbf{Solution détaillée.}

\textbf{Étape 1 : Isoler une inconnue.}
L'équation (1) a tous ses coefficients égaux à $1$. On isole $z$ (choix arbitraire, $x$ ou $y$ iraient aussi) :
\[
z = 6 - x - y \qquad \text{(Équation 1')}
\]

\textbf{Étape 2 : Substituer dans les deux autres équations.}
\begin{align*}
\text{Dans (2) :} \quad 2x - y + (6 - x - y) &= 3 \\
2x - y + 6 - x - y &= 3 \\
x - 2y &= -3 \qquad \textbf{(Eq A)} \\[1em]
\text{Dans (3) :} \quad x + 2y - (6 - x - y) &= 2 \\
x + 2y - 6 + x + y &= 2 \\
2x + 3y &= 8 \qquad \textbf{(Eq B)}
\end{align*}
On obtient le système $2\times 2$ :
\[
\begin{cases}
\text{(A)} & x - 2y = -3 \\
\text{(B)} & 2x + 3y = 8
\end{cases}
\]

\textbf{Étape 3 : Résoudre le système $2\times 2$.}
On utilise la substitution sur (A) : $x = 2y - 3$.
On remplace dans (B) :
\begin{align*}
2(2y - 3) + 3y &= 8 \\
4y - 6 + 3y &= 8 \\
7y &= 14 \\
y &= 2
\end{align*}
On retrouve $x$ : $x = 2(2) - 3 = 1$.

\textbf{Étape 4 : Remontée pour $z$.}
On utilise l'équation (1') : $z = 6 - x - y = 6 - 1 - 2 = 3$.

\textbf{Étape 5 : Vérification dans les trois équations initiales.}
\begin{itemize}
    \item (1) : $1 + 2 + 3 = 6$ \checkmark
    \item (2) : $2(1) - 2 + 3 = 2 - 2 + 3 = 3$ \checkmark
    \item (3) : $1 + 2(2) - 3 = 1 + 4 - 3 = 2$ \checkmark
\end{itemize}
Le triplet $(1;2;3)$ vérifie le système.
\[
\boxed{S = \{(1;2;3)\}}
\]
\end{exoresolu}

\courssub{Modélisation : application concrète}

\begin{exoresolu}[Problème de mélange — Chimie / Agroalimentaire]
Un œnologue souhaite obtenir \SI{100}{\liter} d'un vin titrant \SI{12}{\%} vol. d'alcool en mélangeant trois vins de titres \SI{9}{\%}, \SI{12}{\%} et \SI{15}{\%} vol. Il utilise deux fois plus de vin à \SI{9}{\%} que de vin à \SI{15}{\%}. Déterminer les volumes de chaque vin à mélanger.
\textbf{Solution.}
\emph{Variables :} Soient $x, y, z$ les volumes (en L) des vins à \SI{9}{\%}, \SI{12}{\%}, \SI{15}{\%}.
\emph{Équations :}
\begin{enumerate}
    \item Volume total : $x + y + z = 100$
    \item Quantité d'alcool : $0,09x + 0,12y + 0,15z = 12$ (car $12\% \times 100 = 12$ L d'alcool pur)
    \item Contrainte : $x = 2z$
\end{enumerate}
On substitue $x = 2z$ dans (1) et (2) :
\begin{align*}
(1') &\quad 2z + y + z = 100 \iff y + 3z = 100 \\
(2') &\quad 0,09(2z) + 0,12y + 0,15z = 12 \iff 0,18z + 0,12y + 0,15z = 12 \iff 0,12y + 0,33z = 12
\end{align*}
Système $2\times 2$ :
\[
\begin{cases}
y + 3z = 100 \\
0,12y + 0,33z = 12
\end{cases}
\]
De la première : $y = 100 - 3z$.
Substitution dans la seconde :
\begin{align*}
0,12(100 - 3z) + 0,33z &= 12 \\
12 - 0,36z + 0,33z &= 12 \\
-0,03z &= 0 \iff z = 0
\end{align*}
Alors $y = 100$ et $x = 0$.
\emph{Vérification :} $0 + 100 + 0 = 100$ L. Alcool : $0 + 12 + 0 = 12$ L. Contrainte $x=2z$ : $0=0$.
\emph{Interprétation :} Le vin à \SI{12}{\%} suffit seul ; les deux autres ne sont pas nécessaires (solution « dégénérée » mais mathématiquement correcte).
\[
\boxed{S = \{(0;100;0)\} \text{ soit } 0\,\text{L à 9\%}, 100\,\text{L à 12\%}, 0\,\text{L à 15\%}}
\]
\end{exoresolu}

\courssub{Pièges à éviter et cas particuliers}

\begin{attention}[Erreur fréquente : la substitution stérile]
\emph{Ne substituez jamais l'expression trouvée dans l'équation qui a servi à l'obtenir.}
Si vous avez isolé $z$ dans l'équation (1) pour obtenir $z = 6-x-y$, \textbf{n'insérez pas cette expression dans l'équation (1)}. Vous obtiendriez une identité $6=6$ (ou $0=0$), qui ne fait pas avancer la résolution. \textbf{Substituez toujours dans les DEUX AUTRES équations.}
\end{attention}

\begin{attention}[Cas particuliers — Signalement]
Au cours de la résolution du système $2\times 2$ (étape 3), deux cas peuvent survenir :
\begin{itemize}
    \item \textbf{Contradiction} (ex: $0 = 5$) : Le système $2\times 2$ n'a pas de solution $\implies$ le système $3\times 3$ initial est \textbf{impossible} ($S = \varnothing$).
    \item \textbf{Identité} (ex: $0 = 0$) : Le système $2\times 2$ a une infinité de solutions (une droite). Le système $3\times 3$ initial est \textbf{indéterminé} (infinité de solutions, souvent une droite dans l'espace). On exprime alors deux inconnues en fonction de la troisième (paramètre).
\end{itemize}
Ces cas seront traités plus en profondeur avec la méthode du pivot de Gauss (section suivante).
\end{attention}

\courssub{Exercices d'application}

\begin{exercice}[Entraînement à la substitution]
Résoudre les systèmes suivants par la méthode de substitution.
\begin{enumerate}
    \item $\begin{cases} x + y + z = 0 \\ 2x - y + 3z = 10 \\ -x + 2y - z = -3 \end{cases}$
    \item $\begin{cases} 2x + y - z = 1 \\ x - 2y + z = 2 \\ 3x + y - 2z = 3 \end{cases}$ \emph{(Attention au choix de l'équation de départ !)}
    \item $\begin{cases} x + 2y - z = 4 \\ 2x - y + 2z = 1 \\ 3x + y + z = 5 \end{cases}$
    \item \textbf{(Modélisation)} Un électricien analyse un circuit à trois mailles. Les courants $i_1, i_2, i_3$ (en A) satisfont :
    $\begin{cases} i_1 + i_2 + i_3 = 6 \\ 2i_1 - i_2 + i_3 = 3 \\ i_1 + 2i_2 - i_3 = 2 \end{cases}$
    Déterminer les trois courants.
\end{enumerate}
\end{exercice}

\begin{corrige}[Exercice 1]
On isole $x$ dans (1) : $x = -y - z$.
Substitution dans (2) :
\begin{align*}
2(-y-z) - y + 3z &= 10 \\
-2y - 2z - y + 3z &= 10 \\
-3y + z &= 10 \qquad \textbf{(A)}
\end{align*}
Substitution dans (3) :
\begin{align*}
-(-y-z) + 2y - z &= -3 \\
y + z + 2y - z &= -3 \\
3y &= -3 \iff y = -1
\end{align*}
Dans (A) : $-3(-1) + z = 10 \iff 3 + z = 10 \iff z = 7$.
Remontée : $x = -(-1) - 7 = -6$.
Vérification :
\begin{itemize}
    \item (1) : $(-6) + (-1) + 7 = 0$ \checkmark
    \item (2) : $2(-6) - (-1) + 3(7) = -12 + 1 + 21 = 10$ \checkmark
    \item (3) : $-(-6) + 2(-1) - 7 = 6 - 2 - 7 = -3$ \checkmark
\end{itemize}
\[
\boxed{S = \{(-6;-1;7)\}}
\]
\end{corrige}

\begin{corrige}[Exercice 2]
Le coefficient $1$ est sur $x$ dans (2). On isole $x$ dans (2) : $x = 2 + 2y - z$.
Substitution dans (1) :
\begin{align*}
2(2+2y-z) + y - z &= 1 \\
4 + 4y - 2z + y - z &= 1 \\
5y - 3z &= -3 \qquad \textbf{(A)}
\end{align*}
Substitution dans (3) :
\begin{align*}
3(2+2y-z) + y - 2z &= 3 \\
6 + 6y - 3z + y - 2z &= 3 \\
7y - 5z &= -3 \qquad \textbf{(B)}
\end{align*}
Système (A)(B) : $\begin{cases} 5y - 3z = -3 \\ 7y - 5z = -3 \end{cases}$.
De (A) : $5y = 3z - 3 \iff y = \frac{3z-3}{5}$.
Dans (B) :
\begin{align*}
7\left(\frac{3z-3}{5}\right) - 5z &= -3 \\
\frac{21z - 21}{5} - 5z &= -3 \quad | \times 5 \\
21z - 21 - 25z &= -15 \\
-4z &= 6 \\
z &= -\num{1,5}
\end{align*}
$y = \frac{3(-\num{1,5})-3}{5} = \frac{-\num{4,5}-3}{5} = \frac{-\num{7,5}}{5} = -\num{1,5}$.
$x = 2 + 2(-\num{1,5}) - (-\num{1,5}) = 2 - 3 + \num{1,5} = \num{0,5}$.
Vérification rapide : (1) $2(\num{0,5}) + (-\num{1,5}) - (-\num{1,5}) = 1 - \num{1,5} + \num{1,5} = 1$ \checkmark.
\[
\boxed{S = \{(\num{0,5};-\num{1,5};-\num{1,5})\}}
\]
\end{corrige}

\begin{corrige}[Exercice 3]
On isole $z$ dans (1) : $z = x + 2y - 4$.
Dans (2) :
\begin{align*}
2x - y + 2(x+2y-4) &= 1 \\
2x - y + 2x + 4y - 8 &= 1 \\
4x + 3y &= 9 \qquad \textbf{(A)}
\end{align*}
Dans (3) :
\begin{align*}
3x + y + (x+2y-4) &= 5 \\
4x + 3y - 4 &= 5 \\
4x + 3y &= 9 \qquad \textbf{(B)}
\end{align*}
Les équations (A) et (B) sont identiques. Le système $2\times 2$ a une infinité de solutions.
On paramètre : posons $y = t \in \mathbb{R}$.
$4x = 9 - 3t \iff x = \frac{9-3t}{4}$.
$z = \frac{9-3t}{4} + 2t - 4 = \frac{9-3t+8t-16}{4} = \frac{5t-7}{4}$.
\[
\boxed{S = \left\{\left(\frac{9-3t}{4}; t; \frac{5t-7}{4}\right) \mid t \in \mathbb{R}\right\} \text{ (Système indéterminé)}}
\]
\end{corrige}

\begin{corrige}[Exercice 4]
Le système est identique à l'exemple chiffré complet (variables $i_1, i_2, i_3$ au lieu de $x,y,z$).
On isole $i_3$ dans (1) : $i_3 = 6 - i_1 - i_2$.
Substitution dans (2) : $2i_1 - i_2 + 6 - i_1 - i_2 = 3 \iff i_1 - 2i_2 = -3$.
Substitution dans (3) : $i_1 + 2i_2 - 6 + i_1 + i_2 = 2 \iff 2i_1 + 3i_2 = 8$.
Résolution : $i_1 = 1$, $i_2 = 2$.
Remontée : $i_3 = 6 - 1 - 2 = 3$.
\[
\boxed{i_1 = \SI{1}{\ampere},\; i_2 = \SI{2}{\ampere},\; i_3 = \SI{3}{\ampere}}
\]
\end{corrige}

\coursec{Systèmes 3×3 : pivot de Gauss}

\courssub{Transition depuis la méthode de substitution}

Dans la section précédente, nous avons résolu des systèmes de trois équations à trois inconnues par la \emph{méthode de substitution} : on isole une variable dans une équation, on la remplace dans les deux autres, et l'on continue ainsi. Cette approche est naturelle et efficace lorsque l'une des équations possède un coefficient \textbf{1} ou \textbf{-1} devant une inconnue (par exemple $x = \dots$ ou $y = \dots$), ce qui évite les fractions dès le départ.

Cependant, dans de nombreuses situations techniques (maillage de circuits électriques, équilibre de structures, bilans de matière), les coefficients sont des nombres « quelconques » (entiers relatifs, décimaux) et aucune inconnue n'est isolée simplement. La substitution engendre alors des fractions complexes, sources d'erreurs de calcul. C'est là qu'intervient le \textbf{pivot de Gauss} : une méthode \emph{systématique}, \emph{algorithmique} et \emph{robuste}, qui ne nécessite jamais d'isoler une variable avant de l'éliminer. C'est d'ailleurs l'algorithme standard implémenté dans tous les logiciels de calcul numérique (MATLAB, Python/NumPy, calculatrices CAS) pour résoudre des systèmes de grande taille.

\courssub{Principes : opérations élémentaires sur les équations}

L'idée fondamentale est de transformer le système initial en un système \emph{équivalent} (même ensemble de solutions) mais plus simple, jusqu'à obtenir une forme « en escalier » (triangulaire supérieure) qui se résout immédiatement par \emph{remontée}.

\begin{definition}[Opérations élémentaires sur un système linéaire]
Soit un système d'équations linéaires. Les trois opérations suivantes, appliquées à l'une quelconque des équations, produisent un système \textbf{équivalent} (mêmes solutions) :
\begin{enumerate}
    \item \textbf{Échanger} deux équations (permuter leurs lignes).
    \item \textbf{Multiplier} une équation par un nombre \textbf{non nul} ($k \neq 0$).
    \item \textbf{Remplacer} une équation par la somme d'elle-même et d'un multiple d'une autre équation : $E_i \leftarrow E_i + k \cdot E_j$ ($i \neq j$).
\end{enumerate}
Ces opérations correspondent exactement aux opérations élémentaires sur les lignes de la matrice augmentée du système.
\end{definition}

\begin{attention}[Piège fréquent : oublie du second membre]
Lors de l'opération $E_i \leftarrow E_i + k \cdot E_j$, on doit appliquer la combinaison \textbf{à tous les termes}, y compris le second membre (la constante à droite du signe $=$).
\emph{Erreur typique :} on combine correctement les coefficients de $x, y, z$ mais on oublie d'ajouter $k \times (\text{second membre de } E_j)$ au second membre de $E_i$. Vérifiez systématiquement la colonne de droite à chaque étape.
\end{attention}

\courssub{Méthode du pivot de Gauss : algorithme en 4 étapes}

\begin{methode}[Pivot de Gauss pour un système 3×3]
\textbf{Étape 1 — Choix du premier pivot (colonne $x$).}
Repérez dans la première colonne (coefficients de $x$) un coefficient non nul. Idéalement, choisissez un \textbf{1} ou un \textbf{-1} pour simplifier les calculs (sinon, on divisera la ligne par ce pivot). Échangez cette équation avec la première ligne si nécessaire. Ce coefficient est le \cle{pivot 1}.

\textbf{Étape 2 — Élimination de $x$ dans les lignes 2 et 3.}
Pour chaque ligne $i = 2, 3$, effectuez : $E_i \leftarrow E_i - \frac{\text{coeff}_i(x)}{\text{pivot 1}} \times E_1$.
Après cette étape, les coefficients de $x$ sont nuls dans les équations (2) et (3).

\textbf{Étape 3 — Choix du second pivot (colonne $y$, ligne 2) et élimination de $y$ dans la ligne 3.}
Dans la sous-matrice formée des lignes 2-3 et colonnes $y, z$, repérez un coefficient non nul de $y$ sur la ligne 2 (échangez lignes 2 et 3 si nécessaire). C'est le \cle{pivot 2}. Éliminez $y$ de la ligne 3 : $E_3 \leftarrow E_3 - \frac{\text{coeff}_3(y)}{\text{pivot 2}} \times E_2$.
Le système est maintenant \textbf{triangulaire supérieur} (forme en escalier).

\textbf{Étape 4 — Remontée (résolution à rebours).}
\begin{enumerate}
    \item Résolvez l'équation (3) pour trouver $z$ (une seule inconnue).
    \item Substituez $z$ dans l'équation (2) pour trouver $y$.
    \item Substituez $y$ et $z$ dans l'équation (1) pour trouver $x$.
\end{enumerate}
\end{methode}

\courssub{Illustration : structure triangulaire et pivots}

Le schéma ci-dessous montre la forme finale « en escalier » visée. Les pivots (encadrés) sont les coefficients directeurs de chaque ligne ; ils sont non nuls. Tous les coefficients \emph{sous} les pivots sont nuls.

\begin{popfigure}
\begin{tikzpicture}[
    scale=0.9,
    every node/.style={font=\small},
    pivot/.style={draw=poporange, thick, rounded corners=2pt, inner sep=2pt},
    zero/.style={text=popmuted},
    coeff/.style={text=popdark}
]
% Grille de coefficients
\matrix (m) [matrix of math nodes, nodes={minimum width=1.8em, minimum height=1.8em, anchor=center}, column sep=0.5em, row sep=0.5em] {
|[pivot]| a_{11} & a_{12} & a_{13} & |[coeff]| b_1 \\
0 & |[pivot]| a_{22} & a_{23} & |[coeff]| b_2 \\
0 & 0 & |[pivot]| a_{33} & |[coeff]| b_3 \\
};
% Étiquettes colonnes
\node[above=0.2em of m-1-1] {$x$};
\node[above=0.2em of m-1-2] {$y$};
\node[above=0.2em of m-1-3] {$z$};
\node[above=0.2em of m-1-4] {Cst};
% Étiquettes lignes
\node[left=0.5em of m-1-1] {$E_1$};
\node[left=0.5em of m-2-1] {$E_2$};
\node[left=0.5em of m-3-1] {$E_3$};
% Flèches d'élimination
\draw[popblue, ->, thick] (m-1-1.south east) -- ++(0.5,-0.5) node[below right, popblue, font=\tiny] {Élimine $x$};
\draw[popblue, ->, thick] (m-2-2.south east) -- ++(0.5,-0.5) node[below right, popblue, font=\tiny] {Élimine $y$};
% Accolade triangulaire
\draw[popgreen, thick, decorate, decoration={brace, amplitude=5pt}] (m-3-1.south west) -- (m-1-3.north east) node[midway, above left=0.5em and 1em, popgreen, font=\small] {Forme triangulaire supérieure};
\end{tikzpicture}
\legende{Structure du système après triangularisation : pivots sur la diagonale, zéros en dessous.}
\end{popfigure}

\courssub{Exemple chiffré détaillé : résolution pas à pas}

Résolvons le système suivant par la méthode du pivot de Gauss :
\[
(S) \quad \begin{cases}
x + 2y + z = 8 & (E_1)\\
2x + y - z = 1 & (E_2)\\
x + y + 2z = 9 & (E_3)
\end{cases}
\]

\begin{exemplebox}[Résolution par pivot de Gauss]
\textbf{Étape 1 : Premier pivot sur $x$ (ligne 1).}
Le coefficient de $x$ dans $E_1$ est $1$ (idéal). Pivot 1 = $1$. Pas d'échange nécessaire.

\textbf{Étape 2 : Élimination de $x$ dans $E_2$ et $E_3$.}
\begin{itemize}
    \item Pour $E_2$ : coeff de $x$ est $2$. On fait $E_2 \leftarrow E_2 - 2 \times E_1$.
    \begin{align*}
        (2x + y - z) - 2(x + 2y + z) &= 1 - 2\times 8 \\
        2x + y - z - 2x - 4y - 2z &= 1 - 16 \\
        -3y - 3z &= -15 \quad \text{(nouvelle } E_2\text{)}
    \end{align*}
    \item Pour $E_3$ : coeff de $x$ est $1$. On fait $E_3 \leftarrow E_3 - 1 \times E_1$.
    \begin{align*}
        (x + y + 2z) - 1(x + 2y + z) &= 9 - 1\times 8 \\
        x + y + 2z - x - 2y - z &= 1 \\
        -y + z &= 1 \quad \text{(nouvelle } E_3\text{)}
    \end{align*}
\end{itemize}
Système intermédiaire :
\[
\begin{cases}
x + 2y + z = 8 & (E_1)\\
-3y - 3z = -15 & (E_2)\\
-y + z = 1 & (E_3)
\end{cases}
\]

\textbf{Étape 3 : Second pivot sur $y$ (ligne 2) et élimination dans $E_3$.}
Le coefficient de $y$ dans $E_2$ est $-3$ (non nul). Pivot 2 = $-3$.
Pour simplifier, on peut diviser $E_2$ par $-3$ (opération élémentaire de type 2) : $E_2 \leftarrow \frac{E_2}{-3} \implies y + z = 5$.
On élimine $y$ de $E_3$ : coeff de $y$ dans $E_3$ est $-1$. On fait $E_3 \leftarrow E_3 - \frac{-1}{1} \times E_2$ (en utilisant $E_2$ simplifiée) c'est-à-dire $E_3 \leftarrow E_3 + E_2$.
\begin{align*}
(-y + z) + (y + z) &= 1 + 5 \\
2z &= 6 \quad \text{(nouvelle } E_3\text{)}
\end{align*}
Système triangulaire final :
\[
\begin{cases}
x + 2y + z = 8 & (E_1)\\
y + z = 5 & (E_2)\\
2z = 6 & (E_3)
\end{cases}
\]

\textbf{Étape 4 : Remontée.}
\begin{align*}
\text{Depuis } E_3 &: \quad 2z = 6 \implies \boxed{z = 3}. \\
\text{Dans } E_2 &: \quad y + 3 = 5 \implies \boxed{y = 2}. \\
\text{Dans } E_1 &: \quad x + 2(2) + 3 = 8 \implies x + 7 = 8 \implies \boxed{x = 1}.
\end{align*}
\textbf{Solution :} $(x ; y ; z) = (1 ; 2 ; 3)$.
\end{exemplebox}

\courssub{Représentation matricielle : le tableau des coefficients (complément)}

Pour gagner en lisibilité et éviter de recopier les inconnues $x, y, z$ à chaque ligne, on utilise le \textbf{tableau des coefficients} (matrice augmentée). Chaque ligne correspond à une équation, les trois premières colonnes aux coefficients de $x, y, z$, la dernière colonne (souvent séparée par une barre verticale) aux seconds membres.

\begin{popfigure}
\begin{tikzpicture}[scale=0.85, every node/.style={font=\small}]
% Définition des styles
\tikzset{
    cell/.style={minimum width=2.2em, minimum height=1.8em, anchor=center, draw=popline, inner sep=1pt},
    pivot/.style={cell, fill=poporange!20, draw=poporange, thick, rounded corners=1pt},
    header/.style={cell, fill=popblue!20, draw=popblue, font=\bfseries},
    const/.style={cell, fill=popgreen!15, draw=popgreen!50!black},
    arrow/.style={->, thick, popblue},
    op/.style={font=\tiny, popblue, anchor=west}
}
% Matrice initiale
\matrix (A) [matrix of nodes, nodes={cell}, column sep=-\pgflinewidth, row sep=-\pgflinewidth] {
|[header]| $x$ & |[header]| $y$ & |[header]| $z$ & |[header]| Cst \\
|[pivot]| 1 & 2 & 1 & |[const]| 8 \\
2 & 1 & -1 & |[const]| 1 \\
1 & 1 & 2 & |[const]| 9 \\
};
% Opérations E2 <- E2 - 2E1 ; E3 <- E3 - E1
\node[op] at ($(A-2-1.west)+(-2.5cm,0)$) {$E_2 \leftarrow E_2 - 2E_1$};
\node[op] at ($(A-3-1.west)+(-2.5cm,0)$) {$E_3 \leftarrow E_3 - E_1$};
\draw[arrow] (A-1-1.south) -- ++(0,-0.3) -| (A-2-1.north);
\draw[arrow] (A-1-1.south) -- ++(0,-0.3) -| (A-3-1.north);

% Matrice intermédiaire (décalée vers la droite)
\matrix (B) [matrix of nodes, nodes={cell}, column sep=-\pgflinewidth, row sep=-\pgflinewidth, right=3.5cm of A-1-4.east, anchor=B-1-1.west] {
|[header]| $x$ & |[header]| $y$ & |[header]| $z$ & |[header]| Cst \\
|[pivot]| 1 & 2 & 1 & |[const]| 8 \\
0 & -3 & -3 & |[const]| -15 \\
0 & -1 & 1 & |[const]| 1 \\
};
\node[above=0.2cm of B-1-1, font=\small\bfseries, popdark] {Après élimination de $x$};

% Opérations E2 <- E2/(-3) ; E3 <- E3 + E2
\node[op] at ($(B-2-1.west)+(-2.5cm,0)$) {$E_2 \leftarrow E_2/(-3)$};
\node[op] at ($(B-3-1.west)+(-2.5cm,0)$) {$E_3 \leftarrow E_3 + E_2$};

% Matrice finale
\matrix (C) [matrix of nodes, nodes={cell}, column sep=-\pgflinewidth, row sep=-\pgflinewidth, right=3.5cm of B-1-4.east, anchor=C-1-1.west] {
|[header]| $x$ & |[header]| $y$ & |[header]| $z$ & |[header]| Cst \\
|[pivot]| 1 & 2 & 1 & |[const]| 8 \\
0 & |[pivot]| 1 & 1 & |[const]| 5 \\
0 & 0 & |[pivot]| 2 & |[const]| 6 \\
};
\node[above=0.2cm of C-1-1, font=\small\bfseries, popdark] {Triangulaire (pivots encadrés)};
% Flèches entre matrices
\draw[popmuted, thick, ->] (A-1-4.east) -- (B-1-1.west);
\draw[popmuted, thick, ->] (B-1-4.east) -- (C-1-1.west);
\end{tikzpicture}
\legende{Tableau des coefficients : suivi visuel des opérations de Gauss. Les pivots successifs sont mis en évidence (orange).}
\end{popfigure}

\begin{aretenir}[Lecture du tableau]
\begin{itemize}
    \item Les opérations s'écrivent en marge (ex: $L_2 \leftarrow L_2 - 2L_1$).
    \item On vise des \textbf{1} sur la diagonale (pivots normalisés) pour simplifier la remontée.
    \item Si un pivot est nul, on \textbf{échange} avec une ligne du dessous ayant un coefficient non nul dans cette colonne.
\end{itemize}
\end{aretenir}

\courssub{Comparaison : Substitution vs Pivot de Gauss}

\begin{center}
\begin{tabularx}{\linewidth}{|l|X|X|}
\hline
\rowcolor{popblueL}
\textbf{Critère} & \textbf{Méthode de substitution} & \textbf{Pivot de Gauss} \\
\hline
\textbf{Principe} & Isoler une variable, remplacer dans les autres. & Combiner les équations pour créer des zéros. \\
\hline
\textbf{Idéal quand...} & Un coefficient est $\pm 1$ (facile à isoler). & Coefficients « quelconques », pas de $\pm 1$ évident. \\
\hline
\textbf{Risque fractions} & Élevé dès la première substitution si pas de $\pm 1$. & Maîtrisé : on peut normaliser les pivots (diviser la ligne) \emph{après} élimination. \\
\hline
\textbf{Automatisation} & Difficile à programmer (choix heuristique). & Algorithme standard (utilisé par les machines). \\
\hline
\textbf{Erreurs fréquentes} & Oublier de substituer dans \emph{toutes} les équations. & Erreur de signe dans la combinaison linéaire ; oublie du second membre. \\
\hline
\textbf{Conseil Tle Tech} & Utile pour systèmes simples ou paramétrés. & \textbf{Méthode de référence} pour l'examen : systématique, fiable, notation matricielle compacte. \\
\hline
\end{tabularx}
\end{center}

\begin{savaistu}[L'algorithme des machines]
Le \textbf{pivot de Gauss} (souvent couplé à une stratégie de « pivot partiel » : choisir le plus grand coefficient en valeur absolue comme pivot pour la stabilité numérique) est l'algorithme de base de la résolution de systèmes linéaires en informatique scientifique.
\begin{itemize}
    \item Votre calculatrice graphique (mode \texttt{EQUA} $\to$ \texttt{Simul Eqn}) l'utilise.
    \item Les logiciels de CAO (Calcul de structures, circuits) résolvent des systèmes de $10\,000 \times 10\,000$ équations par des versions optimisées (décomposition $LU$, pivot de Gauss par blocs).
    \item La complexité algorithmique est en $\mathcal{O}(n^3)$ : doubler la taille du système multiplie le temps de calcul par 8.
\end{itemize}
Maîtriser le pivot de Gauss à la main, c'est comprendre ce que fait la machine « dans la boîte noire ».
\end{savaistu}

\begin{exoresolu}[Application technique : Maillage électrique]
\textbf{Énoncé.} Dans un circuit électrique à trois mailles, la loi des mailles (Kirchhoff) donne le système suivant pour les courants $i_1, i_2, i_3$ (en Ampères) :
\[
\begin{cases}
10 i_1 - 4 i_2 - 6 i_3 = 12 & (M_1)\\
-4 i_1 + 8 i_2 - 2 i_3 = 0 & (M_2)\\
-6 i_1 - 2 i_2 + 12 i_3 = -6 & (M_3)
\end{cases}
\]
Résoudre par pivot de Gauss.

\tcblower
\textbf{Solution.}
On écrit le tableau augmenté :
\[
\left(\begin{array}{ccc|c}
10 & -4 & -6 & 12 \\
-4 & 8 & -2 & 0 \\
-6 & -2 & 12 & -6
\end{array}\right)
\]

\textbf{1. Pivot 1 :} Idéalement on veut un 1. On peut diviser $L_1$ par 2 (ou 10). Divisons $L_1$ par 2 pour garder des entiers : $L_1 \leftarrow L_1/2$.
\[
\left(\begin{array}{ccc|c}
5 & -2 & -3 & 6 \\
-4 & 8 & -2 & 0 \\
-6 & -2 & 12 & -6
\end{array}\right)
\]
Élimination de $i_1$ ($L_2 \leftarrow L_2 + \frac{4}{5}L_1$ ; $L_3 \leftarrow L_3 + \frac{6}{5}L_1$). Pour éviter les fractions, on peut faire $L_2 \leftarrow 5L_2 + 4L_1$ et $L_3 \leftarrow 5L_3 + 6L_1$ (opérations combinées autorisées car on multiplie l'équation cible par 5 avant d'ajouter, ce qui ne change pas les solutions).
\begin{align*}
L_2' &= 5(-4\; 8\; -2\; |\; 0) + 4(5\; -2\; -3\; |\; 6) = (0\; 32\; -22\; |\; 24) \\
L_3' &= 5(-6\; -2\; 12\; |\; -6) + 6(5\; -2\; -3\; |\; 6) = (0\; -22\; 42\; |\; 6)
\end{align*}
Tableau :
\[
\left(\begin{array}{ccc|c}
5 & -2 & -3 & 6 \\
0 & 32 & -22 & 24 \\
0 & -22 & 42 & 6
\end{array}\right)
\]

\textbf{2. Pivot 2 :} $32$ sur $L_2$. Simplifions $L_2$ par 2 : $L_2 \leftarrow L_2/2 \to (0\; 16\; -11\; |\; 12)$.
Élimination de $i_2$ dans $L_3$ : $L_3 \leftarrow 16 L_3 + 22 L_2$ (pour éviter fractions).
$L_3' = 16(0\; -22\; 42\; |\; 6) + 22(0\; 16\; -11\; |\; 12) = (0\; 0\; 672-242\; |\; 96+264) = (0\; 0\; 430\; |\; 360)$.
Simplifions $L_3$ par 10 : $(0\; 0\; 43\; |\; 36)$.

\textbf{3. Remontée :}
$43 i_3 = 36 \implies i_3 = \frac{36}{43} \approx 0,837\,\text{A}$.
$16 i_2 - 11 i_3 = 12 \implies 16 i_2 = 12 + \frac{396}{43} = \frac{912}{43} \implies i_2 = \frac{57}{43} \approx 1,326\,\text{A}$.
$5 i_1 - 2 i_2 - 3 i_3 = 6 \implies 5 i_1 = 6 + \frac{114}{43} + \frac{108}{43} = \frac{480}{43} \implies i_1 = \frac{96}{43} \approx 2,233\,\text{A}$.

\textbf{Réponse :} $i_1 = \frac{96}{43}\,\text{A},\; i_2 = \frac{57}{43}\,\text{A},\; i_3 = \frac{36}{43}\,\text{A}$.
\end{exoresolu}

\begin{exercice}[Entraînement : Pivot de Gauss]
Résoudre les systèmes suivants par la méthode du pivot de Gauss (présenter le tableau des coefficients).
\begin{enumerate}
    \item $\begin{cases} 2x - y + z = 5 \\ x + 3y - 2z = 4 \\ 3x + y + z = 10 \end{cases}$
    \item $\begin{cases} x + y + z = 6 \\ 2x - y + 3z = 9 \\ -x + 2y - z = 0 \end{cases}$
    \item (Problème concret) Un atelier produit trois pièces A, B, C. Chaque pièce A demande 2h tour, 1h fraisage, 1h montage. Pièce B : 1h tour, 3h fraisage, 2h montage. Pièce C : 3h tour, 1h fraisage, 2h montage. Les capacités hebdomadaires sont : Tour 20h, Fraisage 18h, Montage 16h. Combien produire de chaque pièce pour saturer l'atelier ? (Modéliser par un système 3×3 et résoudre).
\end{enumerate}
\end{exercice}

\begin{corrige}[Exercice n°1]
\textbf{1.} Tableau : $\left(\begin{array}{ccc|c} 2 & -1 & 1 & 5 \\ 1 & 3 & -2 & 4 \\ 3 & 1 & 1 & 10 \end{array}\right)$.
Échanger $L_1 \leftrightarrow L_2$ (pivot 1 = 1).
$L_2 \leftarrow L_2 - 2L_1$, $L_3 \leftarrow L_3 - 3L_1$.
$\to \left(\begin{array}{ccc|c} 1 & 3 & -2 & 4 \\ 0 & -7 & 5 & -3 \\ 0 & -8 & 7 & -2 \end{array}\right)$.
Pivot 2 = $-7$ ($L_2$). Éliminer $y$ de $L_3$ : $L_3 \leftarrow 7L_3 - 8L_2$ (pour éviter les fractions).
$7L_3 = (0\; -56\; 49\; |\; -14)$ ; $8L_2 = (0\; -56\; 40\; |\; -24)$.
$L_3' = 7L_3 - 8L_2 = (0\; 0\; 9\; |\; 10) \implies 9z = 10 \implies z = \frac{10}{9}$.
Remontée :
$-7y + 5(\frac{10}{9}) = -3 \implies -7y = -3 - \frac{50}{9} = -\frac{77}{9} \implies y = \frac{11}{9}$.
$x + 3(\frac{11}{9}) - 2(\frac{10}{9}) = 4 \implies x + \frac{13}{9} = 4 \implies x = \frac{23}{9}$.
Solution : $(\frac{23}{9} ; \frac{11}{9} ; \frac{10}{9})$.

\textbf{2.} Tableau : $\left(\begin{array}{ccc|c} 1 & 1 & 1 & 6 \\ 2 & -1 & 3 & 9 \\ -1 & 2 & -1 & 0 \end{array}\right)$.
Pivot 1 = 1 ($L_1$). $L_2 \leftarrow L_2 - 2L_1$, $L_3 \leftarrow L_3 + L_1$.
$\to \left(\begin{array}{ccc|c} 1 & 1 & 1 & 6 \\ 0 & -3 & 1 & -3 \\ 0 & 3 & 0 & 6 \end{array}\right)$.
$L_2 \leftarrow L_2/(-3) \to (0\; 1\; -1/3\; |\; 1)$.
$L_3 \leftarrow L_3 - 3L_2 \to (0\; 0\; 1\; |\; 3) \implies z=3$.
$y - 1 = 1 \implies y=2$.
$x + 2 + 3 = 6 \implies x=1$.
Solution : $(1; 2; 3)$.

\textbf{3.} Modélisation : $x, y, z$ nombres de pièces A, B, C.
$\begin{cases} 2x + y + 3z = 20 \\ x + 3y + z = 18 \\ x + 2y + 2z = 16 \end{cases}$.
Tableau : $\left(\begin{array}{ccc|c} 2 & 1 & 3 & 20 \\ 1 & 3 & 1 & 18 \\ 1 & 2 & 2 & 16 \end{array}\right)$.
$L_1 \leftrightarrow L_2$ (pivot 1 = 1).
$L_2 \leftarrow L_2 - 2L_1$, $L_3 \leftarrow L_3 - L_1$.
$\to \left(\begin{array}{ccc|c} 1 & 3 & 1 & 18 \\ 0 & -5 & 1 & -16 \\ 0 & -1 & 1 & -2 \end{array}\right)$.
$L_2 \leftrightarrow L_3$ (pivot 2 = -1 plus simple).
$L_3 \leftarrow L_3 - 5L_2$ (attention signes : $L_2$ est maintenant $0\; -1\; 1\; |\; -2$).
$L_3 : (0\; 0\; -4\; |\; -6) \implies z = 1,5$.
$-y + 1,5 = -2 \implies y = 3,5$.
$x + 3(3,5) + 1,5 = 18 \implies x = 6$.
Solution : 6 pièces A, 3,5 pièces B, 1,5 pièces C. (Ici nombres non entiers $\to$ problème de planification continue, en réalité il faut de la programmation linéaire en nombres entiers).
\end{corrige}

\coursec{Modélisation et résolution de problèmes concrets}

Après avoir maîtrisé les techniques de résolution des systèmes $3 \times 3$ par substitution et par pivot de Gauss, ainsi que le traitement des équations et inéquations du second et du troisième degré, nous abordons maintenant l'objectif final : \textbf{mettre ces outils au service de la résolution de problèmes réels}. En classe de Terminale technique, il ne suffit pas de trouver $x$, $y$, $z$ ; il faut savoir \textbf{traduire} une situation professionnelle (devis, stock, dimensionnement, budget) en langage mathématique, \textbf{résoudre} rigoureusement, puis \textbf{interpréter} et \textbf{valider} la réponse dans le contexte d'origine.

\courssub{Les cinq étapes de la modélisation}

Toute démarche de modélisation suit un processus structuré en cinq étapes. Ne les sautez jamais : c'est la garantie d'une solution exploitable par un chef de chantier, un comptable ou un ingénieur.

\begin{methode}[Modéliser un problème concret en 5 étapes]
\begin{enumerate}
    \item \textbf{Analyser la situation} : Lire, relire, identifier les grandeurs connues, les inconnues recherchées et les contraintes (budget, dimensions, positivité, entier naturel).
    \item \textbf{Choisir les inconnues} : Les nommer explicitement avec leur unité (ex. : « $x$ : prix unitaire du riz en FCFA/kg »).
    \item \textbf{Traduire les contraintes en équations ou inéquations} : Écrire le système ou l'inéquation qui modélise le problème. Vérifier le nombre d'équations = nombre d'inconnues (pour un système linéaire).
    \item \textbf{Résoudre} : Appliquer la méthode adaptée (Gauss, substitution, discriminant, factorisation, étude de signe). Présenter les calculs clairement.
    \item \textbf{Interpréter et valider} : Revenir au texte. Les solutions sont-elles positives ? Entières ? Cohérentes avec les grandeurs physiques (longueur, prix, nombre d'objets) ? Rédiger la conclusion finale en français courant.
\end{enumerate}
\end{methode}

\begin{popfigure}
\begin{tikzpicture}[
    node distance=1.2cm and 2.5cm,
    block/.style={rectangle, draw=popblue, thick, fill=popblueL, text width=2.8cm, text centered, rounded corners, minimum height=1.2cm, font=\small},
    arrow/.style={-Stealth, thick, popdark},
    decision/.style={diamond, draw=poporange, thick, fill=poporangeL, text width=2.5cm, text centered, inner sep=2pt, font=\small}
]
\node[block, align=center] (sit){1. Situation réelle\\Énoncé, données, contraintes};
\node[block, right=of sit, align=center] (inc){2. Choix des inconnues\\$x, y, z$ + unités};
\node[block, right=of inc, align=center] (eq){3. Modèle mathématique\\Système / Équation / Inéquation};
\node[block, right=of eq, align=center] (res){4. Résolution\\Gauss, $\Delta$, signe, factorisation};
\node[decision, below=of res, align=center] (val){5. Validation\\Solutions acceptables ?};
\node[block, right=of res, align=center] (inter){Interprétation\\Réponse finale contextualisée};

\draw[arrow] (sit) -- (inc);
\draw[arrow] (inc) -- (eq);
\draw[arrow] (eq) -- (res);
\draw[arrow] (res) -- (val);
\draw[arrow] (val) -- node[left, font=\tiny, text=popgreen] {OUI} (inter);
\draw[arrow] (val.west) -- ++(-1.5,0) |- node[near start, above, font=\tiny, text=popred] {NON} (eq.west);
\end{tikzpicture}
\legende{Organigramme des 5 étapes de la modélisation : la validation peut imposer un retour à la modélisation (étape 3).}
\end{popfigure}

\courssub{Problèmes de prix et de commerce : systèmes 3×3}

Les situations d'achat/vente multiples (marché, cantine, quincaillerie, stock de pièces détachées) conduisent naturellement à des systèmes linéaires de trois équations à trois inconnues. La méthode du pivot de Gauss est alors la plus efficace et la moins sujette aux erreurs de calcul.

\begin{exoresolu}[Problème des prix au marché Mokolo]
\textbf{Énoncé :} Au marché Mokolo, un revendeur achète trois types de denrées : du riz, des haricots et de l'huile. Il constate que :
\begin{itemize}
    \item 2 kg de riz, 3 kg de haricots et 1 litre d'huile coûtent 4\,200 FCFA.
    \item 1 kg de riz, 2 kg de haricots et 2 litres d'huile coûtent 3\,800 FCFA.
    \item 3 kg de riz, 1 kg de haricots et 1 litre d'huile coûtent 3\,700 FCFA.
\end{itemize}
Déterminer le prix unitaire (en FCFA) de chaque denrée.
\tcblower
\textbf{Solution détaillée :}

\noindent\textbf{1. Analyse et choix des inconnues.}
Soit $x$ le prix du kg de riz (FCFA/kg), $y$ le prix du kg de haricots (FCFA/kg), $z$ le prix du litre d'huile (FCFA/L). $x, y, z \in \mathbb{R}_+^*$.

\noindent\textbf{2. Traduction en système.}
\[
\begin{cases}
2x + 3y + z = 4200 & (L_1)\\
x + 2y + 2z = 3800 & (L_2)\\
3x + y + z = 3700 & (L_3)
\end{cases}
\]

\noindent\textbf{3. Résolution par pivot de Gauss (matrice augmentée).}
\[
\left(\begin{array}{ccc|c}
2 & 3 & 1 & 4200 \\
1 & 2 & 2 & 3800 \\
3 & 1 & 1 & 3700
\end{array}\right)
\]
On permute $L_1$ et $L_2$ pour avoir un pivot $1$ en haut à gauche :
\[
\xrightarrow{L_1 \leftrightarrow L_2}
\left(\begin{array}{ccc|c}
1 & 2 & 2 & 3800 \\
2 & 3 & 1 & 4200 \\
3 & 1 & 1 & 3700
\end{array}\right)
\]
On annule la première colonne sous le pivot :
\[
\xrightarrow{\substack{L_2 \leftarrow L_2 - 2L_1 \\ L_3 \leftarrow L_3 - 3L_1}}
\left(\begin{array}{ccc|c}
1 & 2 & 2 & 3800 \\
0 & -1 & -3 & -3400 \\
0 & -5 & -5 & -7700
\end{array}\right)
\]
On simplifie $L_2$ et $L_3$ (division par $-1$ et $-5$) :
\[
\xrightarrow{\substack{L_2 \leftarrow -L_2 \\ L_3 \leftarrow L_3 / (-5)}}
\left(\begin{array}{ccc|c}
1 & 2 & 2 & 3800 \\
0 & 1 & 3 & 3400 \\
0 & 1 & 1 & 1540
\end{array}\right)
\]
On annule sous le second pivot :
\[
\xrightarrow{L_3 \leftarrow L_3 - L_2}
\left(\begin{array}{ccc|c}
1 & 2 & 2 & 3800 \\
0 & 1 & 3 & 3400 \\
0 & 0 & -2 & -1860
\end{array}\right)
\]
On normalise $L_3$ :
\[
\xrightarrow{L_3 \leftarrow L_3 / (-2)}
\left(\begin{array}{ccc|c}
1 & 2 & 2 & 3800 \\
0 & 1 & 3 & 3400 \\
0 & 0 & 1 & 930
\end{array}\right)
\]

\noindent\textbf{4. Remontée.}
$z = 930$ FCFA/L.
$y + 3z = 3400 \implies y = 3400 - 2790 = 610$ FCFA/kg.
$x + 2y + 2z = 3800 \implies x = 3800 - 1220 - 1860 = 720$ FCFA/kg.

\noindent\textbf{5. Interprétation et validation.}
Les trois prix sont positifs et entiers, ce qui est cohérent.
\textbf{Conclusion :} Le riz coûte \textbf{720 FCFA/kg}, les haricots \textbf{610 FCFA/kg}, l'huile \textbf{930 FCFA/L}.
\end{exoresolu}

\begin{attention}[Piège : solutions mathématiquement exactes mais contextuellement impossibles]
Un système $3 \times 3$ a toujours une solution unique (si le déterminant est non nul), mais cette solution peut être \textbf{négative} ou \textbf{non entière}.
\begin{itemize}
    \item \textbf{Prix négatif :} Impossible. Cela signale une erreur dans l'énoncé, la modélisation ou les données du problème.
    \item \textbf{Nombre d'objets fractionnaire :} Si les inconnues représentent des quantités d'objets indivisibles (nombre de machines, de sacs, de personnes), la solution doit être un entier naturel. Si le système donne $x = 12,5$, il faut conclure : « Aucune solution entière ne convient, le problème n'a pas de solution dans la réalité » ou revoir les hypothèses (ex. : objets divisibles comme le sable, le ciment).
\end{itemize}
\textbf{À l'examen :} Si vous trouvez $x = -50$, écrivez : « Solution rejetée car un prix ne peut pas être négatif. » Ne laissez jamais une valeur aberrante sans commentaire.
\end{attention}

\courssub{Problèmes de dimensions et d'aires : inéquations du second degré}

Le dimensionnement de pièces (tôles, contreplaqué, terrain, vitrage) conduit souvent à une contrainte de surface minimale ou maximale, ce qui se traduit par une inéquation du second degré.

\begin{exoresolu}[Dimensions d'une plaque de contreplaqué]
\textbf{Énoncé :} Un menuisier doit découper une plaque rectangulaire de contreplaqué dont la longueur dépasse la largeur de 50 cm. La surface de la plaque doit être \textbf{au moins} 1,5 m² pour réaliser un panneau de porte. Déterminer les dimensions minimales (en cm) de la plaque.
\tcblower
\textbf{Solution détaillée :}

\noindent\textbf{1. Inconnues et unités.}
Soit $x$ la largeur en \textbf{centimètres} ($x > 0$).
La longueur est $x + 50$ (en cm).
La surface $S$ en cm² : $S = x(x + 50)$.
La contrainte : $S \ge 1,5 \text{ m}^2 = 15\,000 \text{ cm}^2$. (Attention à la conversion d'unités !)

\noindent\textbf{2. Inéquation.}
\[
x(x + 50) \ge 15\,000 \iff x^2 + 50x - 15\,000 \ge 0
\]
Soit $f(x) = x^2 + 50x - 15\,000$. C'est un trinôme du second degré, $a=1>0$.

\noindent\textbf{3. Résolution (étude du signe).}
$\Delta = 50^2 - 4(1)(-15\,000) = 2\,500 + 60\,000 = 62\,500 = 250^2$.
Racines :
$x_1 = \frac{-50 - 250}{2} = -150$ (rejetée car $x>0$)
$x_2 = \frac{-50 + 250}{2} = 100$.

Tableau de signes (pour $x > 0$) :
\begin{center}
\begin{tabular}{|c|ccccc|}\hline
$x$ & $0$ & & $100$ & & $+\infty$ \\ \hline
$f(x)$ & $-$ & & $0$ & $+$ & \\ \hline
\end{tabular}
\end{center}
$f(x) \ge 0 \iff x \in [100, +\infty[$.

\noindent\textbf{4. Interprétation.}
La largeur doit être \textbf{au minimum 100 cm} (1 m).
La longueur correspondante est $100 + 50 = 150$ cm (1,5 m).
Surface minimale : $100 \times 150 = 15\,000$ cm² = 1,5 m². \textbf{Condition respectée.}

\noindent\textbf{Conclusion :} Le menuisier doit prévoir une plaque d'au moins \textbf{100 cm de large sur 150 cm de long}.
\end{exoresolu}

\courssub{Problèmes de contraintes budgétaires et de production : inéquations du troisième degré}

La modélisation de coûts de production (coûts fixes + coûts variables non linéaires) ou de volumes (caisses, citernes) peut mener à des inéquations de degré 3. La méthode de résolution reste l'étude de signe via factorisation (racine évidente + trinôme).

\begin{exoresolu}[Production minimale pour rentabilité]
\textbf{Énoncé :} Un atelier de fabrication de tuyaux PVC a un bénéfice net (en milliers de FCFA) modélisé par $B(x) = -x^3 + 6x^2 - 11x + 6$, où $x$ est le nombre de centaines de tuyaux produits ($x \ge 0$). Pour quelle production l'atelier est-il bénéficiaire (bénéfice $> 0$) ?
\tcblower
\textbf{Solution détaillée :}

\noindent\textbf{1. Modélisation.}
On cherche $B(x) > 0 \iff -x^3 + 6x^2 - 11x + 6 > 0$.
On multiplie par $-1$ (ce qui change le sens de l'inéquation) :
\[
x^3 - 6x^2 + 11x - 6 < 0
\]
Soit $P(x) = x^3 - 6x^2 + 11x - 6$.

\noindent\textbf{2. Factorisation (recherche d'une racine évidente).}
Diviseurs du terme constant ($-6$) : $\pm 1, \pm 2, \pm 3, \pm 6$.
$P(1) = 1 - 6 + 11 - 6 = 0$. \quad $x=1$ est une racine.
Division euclidienne de $P(x)$ par $(x-1)$ :
$P(x) = (x-1)(x^2 - 5x + 6)$.
Le trinôme se factorise : $x^2 - 5x + 6 = (x-2)(x-3)$.
Donc $P(x) = (x-1)(x-2)(x-3)$.

\noindent\textbf{3. Tableau de signes (racines 1, 2, 3 ; domaine $x \ge 0$).}
\begin{center}
\begin{tabular}{|c|ccccccccc|}\hline
$x$ & $0$ & & $1$ & & $2$ & & $3$ & & $+\infty$ \\ \hline
$x-1$ & $-$ & & $0$ & $+$ & $+$ & $+$ & $+$ & $+$ & $+$ \\
$x-2$ & $-$ & $-$ & $-$ & $-$ & $0$ & $+$ & $+$ & $+$ & $+$ \\
$x-3$ & $-$ & $-$ & $-$ & $-$ & $-$ & $-$ & $0$ & $+$ & $+$ \\ \hline
$P(x)$ & $-$ & & $0$ & $+$ & $0$ & $-$ & $0$ & $+$ & \\ \hline
\end{tabular}
\end{center}
$P(x) < 0$ pour $x \in [0, 1[ \cup ]2, 3[$ (on restreint à $x \ge 0$).

\noindent\textbf{4. Interprétation.}
$x$ est en centaines de tuyaux.
L'atelier est bénéficiaire pour une production :
\begin{itemize}
    \item comprise entre \textbf{0 et 100 tuyaux} (exclu),
    \item ou comprise entre \textbf{200 et 300 tuyaux} (exclus).
\end{itemize}
Entre 100 et 200, et au-delà de 300, l'atelier est déficitaire.
\textbf{Conclusion :} Pour être rentable, l'atelier doit produire \textbf{moins de 100 tuyaux ou entre 200 et 300 tuyaux}.
\end{exoresolu}

\courssub{Interprétation critique et validation : la rédaction finale}

La résolution mathématique n'est que la moitié du travail. La seconde moitié, tout aussi notée au Probatoire, est la \textbf{rédaction justifiée}.

\begin{aretenir}[Rédaction d'une démarche complète (attendus Probatoire)]
Une réponse complète comporte obligatoirement :
\begin{enumerate}
    \item \textbf{L'énoncé de la méthode} : « Je pose $x$ = ... Je traduis par le système... / l'inéquation... »
    \item \textbf{Les calculs lisibles et alignés} : Pivot de Gauss présenté en matrice, calcul de $\Delta$, factorisation, tableau de signes tracé à la règle.
    \item \textbf{Le rejet explicite des solutions parasites} : « $x = -3$ est rejeté car une longueur est positive. » « $x = 12,5$ est rejeté car on ne peut pas acheter un demi-sac de ciment. »
    \item \textbf{La vérification dans l'énoncé} : Remplacer les valeurs trouvées dans les phrases de l'énoncé (pas seulement dans les équations).
    \item \textbf{La conclusion rédigée en français} : « Le prix du kg de riz est de 720 FCFA. » « La largeur minimale est de 1 m. »
\end{enumerate}
\end{aretenir}

\begin{savaistu}[Les mathématiques au service de l'ingénierie et de la gestion au Cameroun]
Au Cameroun, les techniciens supérieurs (BTS, HND, DUT) utilisent quotidiennement ces outils :
\begin{itemize}
    \item \textbf{Gestion des stocks / Achats :} Les services approvisionnement des grandes entreprises (CDC, SONARA, brasseries, cimenteries) résolvent des systèmes $n \times n$ pour optimiser les commandes multi-fournisseurs sous contraintes de budget et de capacité de stockage.
    \item \textbf{Dimensionnement structures :} Les bureaux d'études (génie civil, métallique) résolvent des équations du second degré pour calculer les sections de poutres, l'armature des poteaux, ou les dimensions de fondations (surface portante).
    \item \textbf{Contrôle de coûts / Seuil de rentabilité :} Les comptables de gestion modélisent les coûts fixes/variables (souvent polynômes degré 2 ou 3) et déterminent les seuils de rentabilité par résolution d'inéquations.
    \item \textbf{Électricité / Électrotechnique :} La loi des nœuds (Kirchhoff) dans un circuit à 3 mailles donne un système $3 \times 3$ pour trouver les intensités.
\end{itemize}
Maîtriser la modélisation en Terminale technique, c'est acquérir le réflexe professionnel : \textbf{Problème réel $\to$ Modèle math $\to$ Solution math $\to$ Décision technique}.
\end{savaistu}

\begin{exercice}[Entraînement : Modélisation complète]
\begin{enumerate}
    \item \textbf{Cantine scolaire.} La cantine achète du riz, du poisson et de l'huile.
    \begin{itemize}
        \item 5 kg de riz + 2 kg de poisson + 3 L d'huile = 14\,500 FCFA.
        \item 3 kg de riz + 4 kg de poisson + 2 L d'huile = 16\,000 FCFA.
        \item 2 kg de riz + 1 kg de poisson + 4 L d'huile = 10\,500 FCFA.
    \end{itemize}
    Trouver les prix unitaires par pivot de Gauss.
    \item \textbf{Atelier métallique.} On veut fabriquer un bac rectangulaire sans couvercle, à base carrée de côté $x$ (dm), de hauteur $h = x - 2$ (dm). Le volume doit être \textbf{supérieur ou égal à} 9 dm³. Déterminer les dimensions minimales du bac (inéquation degré 3).
    \item \textbf{Budget carburant.} Un camion consomme $C(v) = 0,005v^2 - 0,3v + 10$ litres/100km à la vitesse $v$ (km/h). Le chauffeur a 50 litres pour faire 300 km. Quelle plage de vitesse doit-il respecter ? (Inéquation second degré).
\end{enumerate}
\end{exercice}

\begin{corrige}[Exercice 1 -- Cantine scolaire]
\textbf{Inconnues :} $x$ (riz), $y$ (poisson), $z$ (huile) en FCFA/kg ou L.
\textbf{Système :}
\[
\begin{cases}
5x + 2y + 3z = 14500 & (L_1)\\
3x + 4y + 2z = 16000 & (L_2)\\
2x + y + 4z = 10500 & (L_3)
\end{cases}
\]

\textbf{Résolution par pivot de Gauss :}
Matrice augmentée :
\[
\left(\begin{array}{ccc|c}
5 & 2 & 3 & 14500 \\
3 & 4 & 2 & 16000 \\
2 & 1 & 4 & 10500
\end{array}\right)
\]
$L_1 \leftrightarrow L_3$ (pivot 2 ou 1 ? Mieux : $L_3$ a 1 en 2ème col. Faisons $L_1 \leftrightarrow L_3$ pour avoir 2 en (1,1) ou divisons. Choisissons $L_1 \leftarrow L_1/5$ ou permutation. Utilisons $L_1 \leftrightarrow L_3$ pour avoir 2, puis divisons.)
\[
\xrightarrow{L_1 \leftrightarrow L_3}
\left(\begin{array}{ccc|c}
2 & 1 & 4 & 10500 \\
3 & 4 & 2 & 16000 \\
5 & 2 & 3 & 14500
\end{array}\right)
\xrightarrow{L_1 \leftarrow L_1/2}
\left(\begin{array}{ccc|c}
1 & 0.5 & 2 & 5250 \\
3 & 4 & 2 & 16000 \\
5 & 2 & 3 & 14500
\end{array}\right)
\]
\[
\xrightarrow{\substack{L_2 \leftarrow L_2 - 3L_1 \\ L_3 \leftarrow L_3 - 5L_1}}
\left(\begin{array}{ccc|c}
1 & 0.5 & 2 & 5250 \\
0 & 2.5 & -4 & 250 \\
0 & -0.5 & -7 & -11750
\end{array}\right)
\xrightarrow{L_2 \leftarrow L_2/2.5}
\left(\begin{array}{ccc|c}
1 & 0.5 & 2 & 5250 \\
0 & 1 & -1.6 & 100 \\
0 & -0.5 & -7 & -11750
\end{array}\right)
\]
\[
\xrightarrow{\substack{L_1 \leftarrow L_1 - 0.5L_2 \\ L_3 \leftarrow L_3 + 0.5L_2}}
\left(\begin{array}{ccc|c}
1 & 0 & 2.8 & 5200 \\
0 & 1 & -1.6 & 100 \\
0 & 0 & -7.8 & -11700
\end{array}\right)
\xrightarrow{L_3 \leftarrow L_3 / (-7.8)}
\left(\begin{array}{ccc|c}
1 & 0 & 2.8 & 5200 \\
0 & 1 & -1.6 & 100 \\
0 & 0 & 1 & 1500
\end{array}\right)
\]
Remontée :
$z = 1500$
$y - 1,6(1500) = 100 \implies y = 100 + 2400 = 2500$
$x + 2,8(1500) = 5200 \implies x = 5200 - 4200 = 1000$

\textbf{Vérification :}
$5(1000) + 2(2500) + 3(1500) = 5000 + 5000 + 4500 = 14500$ \checkmark
$3(1000) + 4(2500) + 2(1500) = 3000 + 10000 + 3000 = 16000$ \checkmark
$2(1000) + 1(2500) + 4(1500) = 2000 + 2500 + 6000 = 10500$ \checkmark

\textbf{Conclusion :} Riz \textbf{1000 FCFA/kg}, Poisson \textbf{2500 FCFA/kg}, Huile \textbf{1500 FCFA/L}.
\end{corrige}

\begin{corrige}[Exercice 2 -- Bac métallique]
\textbf{Inconnue :} $x$ longueur du côté de la base carrée (dm), $x > 2$ (car hauteur $x-2 > 0$).
\textbf{Volume :} $V(x) = x^2 \cdot (x-2) = x^3 - 2x^2$.
\textbf{Contrainte :} $V(x) \ge 9 \iff x^3 - 2x^2 - 9 \ge 0$.
Soit $P(x) = x^3 - 2x^2 - 9$.
Recherche racine : diviseurs de 9 : $\pm 1, \pm 3, \pm 9$.
$P(3) = 27 - 18 - 9 = 0$. \quad $x=3$ est racine.
Division par $(x-3)$ : $P(x) = (x-3)(x^2 + x + 3)$.
Le trinôme $x^2+x+3$ a $\Delta = 1 - 12 = -11 < 0$, il est toujours $>0$ (car $a=1>0$).
Signe de $P(x)$ = Signe de $(x-3)$.
$P(x) \ge 0 \iff x - 3 \ge 0 \iff x \ge 3$.

\textbf{Interprétation :} $x \ge 3$ dm. La hauteur $h = x-2 \ge 1$ dm.
\textbf{Conclusion :} Le bac doit avoir une base d'au moins \textbf{3 dm de côté} et une hauteur d'au moins \textbf{1 dm} (Volume minimal 9 dm³).
\end{corrige}

\begin{corrige}[Exercice 3 -- Carburant]
\textbf{Données :} Conso $C(v) = 0,005v^2 - 0,3v + 10$ L/100km. Distance 300 km $\to$ 3 $\times$ 100km.
Carburant dispo : 50 L.
\textbf{Inéquation :} $3 \times C(v) \le 50$
\[
3(0,005v^2 - 0,3v + 10) \le 50
\]
\[
0,015v^2 - 0,9v + 30 \le 50
\]
\[
0,015v^2 - 0,9v - 20 \le 0
\]
Multiplier par 1000 pour éviter les décimales :
\[
15v^2 - 900v - 20000 \le 0
\]
Diviser par 5 :
\[
3v^2 - 180v - 4000 \le 0
\]
\textbf{Résolution :} $a=3>0$. $\Delta = (-180)^2 - 4(3)(-4000) = 32400 + 48000 = 80400$.
$\sqrt{\Delta} = \sqrt{80400} = \sqrt{400 \times 201} = 20\sqrt{201} \approx 283,55$.
Racines :
$v_1 = \frac{180 - 283,55}{6} \approx -17,26$ (rejetée, $v>0$)
$v_2 = \frac{180 + 283,55}{6} \approx 77,26$

Tableau de signes (trinôme $>0$ hors racines) :
$3v^2 - 180v - 4000 \le 0 \iff v \in [v_1, v_2]$.
Comme $v>0$ : $v \in [0, 77,26]$.

\textbf{Conclusion :} Le chauffeur doit rouler à une vitesse comprise entre \textbf{0 et 77 km/h} (arrondi à l'entier inférieur pour sécurité : 77 km/h max).
\end{corrige}

\newpage

\coursec{Activité d'intégration}

\coursec{Équations, inéquations, systèmes : activité d'intégration}

\courssub{Objectifs de l'activité}

À l'issue de cette activité d'intégration, tu seras capable de :
\begin{itemize}
\item traduire une situation concrète de la vie scolaire et économique camerounaise par un \cle{système linéaire} de trois équations à trois inconnues ;
\item résoudre ce système par la méthode de \cle{substitution}, puis vérifier le résultat par le \cle{pivot de Gauss} ;
\item modéliser une contrainte budgétaire par une \cle{inéquation} et déterminer l'ensemble de ses solutions ;
\item rédiger et justifier une démarche complète, une réponse argumentée et une conclusion exploitable par un décideur (le trésorier de la coopérative).
\end{itemize}

\begin{attention}[Compétences évaluées]
Cette activité mobilise les deux savoir-faire officiels du programme : \emph{appliquer les notions de l'unité à des situations concrètes de la vie courante} et \emph{rédiger et justifier une démarche, une réponse ou une conclusion}. La qualité de la rédaction compte autant que l'exactitude des calculs.
\end{attention}

\courssub{Situation}

La coopérative scolaire du Lycée Technique de Nkol-Eton (Yaoundé) revend aux élèves, à prix coûtant, trois articles : des \textbf{cahiers} de 200 pages, des \textbf{stylos} bleus et des \textbf{rames de papier} format A4. Le trésorier, élève en Terminale technique, a égaré la facture détaillée des prix unitaires. Il ne dispose plus que des bordereaux de trois commandes passées au cours du premier trimestre auprès des fournisseurs du quartier Mokolo.

\begin{center}
\begin{tabularx}{\linewidth}{|l|X|X|X|X|}
\hline
\rowcolor{popblueL} \textbf{Commande} & \textbf{Cahiers} & \textbf{Stylos} & \textbf{Rames} & \textbf{Montant total (FCFA)} \\
\hline
Commande n\textsuperscript{o}1 (septembre) & 100 & 50 & 20 & 90\,000 \\
\hline
Commande n\textsuperscript{o}2 (octobre) & 60 & 100 & 10 & 65\,000 \\
\hline
Commande n\textsuperscript{o}3 (novembre) & 40 & 30 & 40 & 86\,000 \\
\hline
\end{tabularx}
\end{center}

On note $x$ le prix unitaire d'un cahier, $y$ celui d'un stylo et $z$ celui d'une rame de papier, en FCFA.

Pour la rentrée du deuxième trimestre, la coopérative dispose d'un \textbf{budget maximal de 150\,000 FCFA}. Le bureau a décidé d'acheter des stylos à 500 FCFA l'unité (prix désormais connu) et des rames à 2\,000 FCFA l'unité, en quantités déjà fixées : 100 stylos et 20 rames. Le reste du budget servira à acheter des cahiers, dont le prix unitaire est celui trouvé grâce aux trois commandes. Une contrainte s'impose : pour équiper les élèves de six classes, il faut \textbf{au moins 80 cahiers}. Le trésorier doit déterminer le nombre de cahiers qu'il est possible de commander, vérifier que la contrainte des 80 cahiers peut être respectée, et proposer la commande optimale (le plus grand nombre entier de cahiers possible).

\courssub{Consignes}

Travail en groupes de trois ou quatre élèves. Chaque groupe rédige un rapport unique.

\begin{enumerate}
\item Traduire chacune des trois commandes par une équation reliant $x$, $y$ et $z$. Écrire le système $(S)$ de trois équations à trois inconnues obtenu.
\item Résoudre le système $(S)$ par la méthode de \textbf{substitution} : exprimer une inconnue en fonction des deux autres à partir d'une équation, puis se ramener à un système de deux équations à deux inconnues. Conclure sur les trois prix unitaires.
\item \textbf{Vérifier} le résultat en résolvant à nouveau $(S)$ par la méthode du \textbf{pivot de Gauss} (opérations élémentaires sur les équations, forme triangulaire, remontée). Les deux méthodes donnent-elles les mêmes prix ?
\item Vérifier que les prix trouvés sont cohérents avec la réalité du marché (un cahier coûte-t-il plus cher qu'un stylo ?) et avec chacun des trois bordereaux.
\item En utilisant le prix du cahier trouvé, traduire la contrainte budgétaire du deuxième trimestre par une inéquation d'inconnue $n$ (nombre de cahiers), en tenant compte des 100 stylos et des 20 rames déjà prévus.
\item Résoudre cette inéquation. La contrainte des 80 cahiers minimum est-elle satisfaisable ? Quel est le nombre maximal de cahiers commandables ?
\item Rédiger un \textbf{rapport d'une demi-page} à l'attention du président de la coopérative : prix unitaires retrouvés, commande optimale proposée, dépense totale exacte et reliquat éventuel du budget. La démarche doit être justifiée étape par étape.
\end{enumerate}

\courssub{Ressources à mobiliser}

\begin{aretenir}[Boîte à outils de la leçon]
\begin{itemize}
\item \textbf{Système $3\times 3$ par substitution} : on exprime une inconnue (par exemple $z$) en fonction des autres dans une équation, on la remplace dans les deux autres équations ; on obtient un système $2\times 2$ que l'on résout, puis on remonte.
\item \textbf{Pivot de Gauss} : par les opérations autorisées $L_i \leftarrow L_i + \lambda L_j$ et $L_i \leftrightarrow L_j$, on élimine $x$ des équations 2 et 3, puis $y$ de l'équation 3 ; le système triangulaire se résout par remontée.
\item \textbf{Inéquation du premier degré} $an + b \le c$ : on isole $n$ ; attention, la division par un nombre \emph{négatif} change le sens de l'inégalité.
\item \textbf{Solution entière} : dans un contexte concret, on ne retient que les valeurs entières compatibles avec la situation (on ne commande pas une fraction de cahier).
\end{itemize}
\end{aretenir}

\courssub{Démarche et corrigé commenté}

\begin{exoresolu}[Le budget de la coopérative scolaire — corrigé complet]
\textbf{Question 1.} Traduire les trois commandes par un système $(S)$.

\textbf{Question 2.} Résoudre $(S)$ par substitution.

\textbf{Question 3.} Vérifier par le pivot de Gauss.

\textbf{Questions 5 et 6.} Modéliser la contrainte budgétaire et la résoudre.

\textbf{Question 7.} Rédiger le rapport.

\tcblower

\textbf{Étape 1 : mise en équations.} Chaque bordereau donne une équation : quantité $\times$ prix unitaire, sommée sur les trois articles, égale au montant total.
\[
(S)\quad \begin{cases}
100x + 50y + 20z = 90\,000 & (L_1)\\
60x + 100y + 10z = 65\,000 & (L_2)\\
40x + 30y + 40z = 86\,000 & (L_3)
\end{cases}
\]

\textbf{Étape 2 : résolution par substitution.} Simplifions d'abord chaque équation en divisant par 10 pour alléger les calculs :
\[
\begin{cases}
10x + 5y + 2z = 9\,000 & (L_1')\\
6x + 10y + z = 6\,500 & (L_2')\\
4x + 3y + 4z = 8\,600 & (L_3')
\end{cases}
\]
Dans $(L_2')$, le coefficient de $z$ est 1 : on exprime $z$ en fonction de $x$ et $y$.
\[
z = 6\,500 - 6x - 10y
\]
On substitue cette expression dans $(L_1')$ et $(L_3')$ :
\begin{align*}
(L_1') : \quad & 10x + 5y + 2(6\,500 - 6x - 10y) = 9\,000 \\
& 10x + 5y + 13\,000 - 12x - 20y = 9\,000 \\
& -2x - 15y = -4\,000 \quad \Longrightarrow \quad 2x + 15y = 4\,000 \quad (E_1)\\[4pt]
(L_3') : \quad & 4x + 3y + 4(6\,500 - 6x - 10y) = 8\,600 \\
& 4x + 3y + 26\,000 - 24x - 40y = 8\,600 \\
& -20x - 37y = -17\,400 \quad \Longrightarrow \quad 20x + 37y = 17\,400 \quad (E_2)
\end{align*}
On résout le système $2\times 2$ $(E_1)$-$(E_2)$. De $(E_1)$, on tire $x = \dfrac{4\,000 - 15y}{2}$. On reporte dans $(E_2)$ :
\begin{align*}
20 \times \frac{4\,000 - 15y}{2} + 37y &= 17\,400 \\
10(4\,000 - 15y) + 37y &= 17\,400 \\
40\,000 - 150y + 37y &= 17\,400 \\
-113y &= -22\,600 \\
y &= 200
\end{align*}
On remonte vers $x$ :
\[
x = \frac{4\,000 - 15 \times 200}{2} = \frac{4\,000 - 3\,000}{2} = 500
\]
On remonte vers $z$ avec l'expression trouvée au départ :
\[
z = 6\,500 - 6 \times 500 - 10 \times 200 = 6\,500 - 3\,000 - 2\,000 = 1\,500
\]

\textbf{Conclusion de l'étape 2 :} un cahier coûte $x = 500$ FCFA, un stylo $y = 200$ FCFA, une rame $z = 1\,500$ FCFA.

\textbf{Étape 3 : vérification par pivot de Gauss.} On repart du système simplifié $(L_1')-(L_2')-(L_3')$. On élimine $x$ des lignes 2 et 3.
\begin{itemize}
\item $L_2 \leftarrow 5 \times L_2' - 3 \times L_1'$ (pour annuler $10x$ et $6x$ via $30x$) :
\begin{align*}
5(6x + 10y + z) - 3(10x + 5y + 2z) &= 5 \times 6\,500 - 3 \times 9\,000 \\
30x + 50y + 5z - 30x - 15y - 6z &= 32\,500 - 27\,000 \\
35y - z &= 5\,500 \quad (L_2'')
\end{align*}
\item $L_3 \leftarrow 5 \times L_3' - 2 \times L_1'$ (pour annuler $10x$ et $4x$ via $20x$) :
\begin{align*}
5(4x + 3y + 4z) - 2(10x + 5y + 2z) &= 5 \times 8\,600 - 2 \times 9\,000 \\
20x + 15y + 20z - 20x - 10y - 4z &= 43\,000 - 18\,000 \\
5y + 16z &= 25\,000 \quad (L_3'')
\end{align*}
\end{itemize}
On élimine $y$ de la nouvelle ligne 3 : $L_3 \leftarrow 7 \times L_3'' - L_2''$ (coefficients de $y$ : $35$ et $35$).
\begin{align*}
7(5y + 16z) - (35y - z) &= 7 \times 25\,000 - 5\,500 \\
35y + 112z - 35y + z &= 175\,000 - 5\,500 \\
113z &= 169\,500 \\
z &= 1\,500
\end{align*}
Remontée :
\begin{itemize}
\item $(L_2'')$ : $35y = 5\,500 + z = 5\,500 + 1\,500 = 7\,000 \quad \Rightarrow \quad y = 200$.
\item $(L_1')$ : $10x = 9\,000 - 5y - 2z = 9\,000 - 1\,000 - 3\,000 = 5\,000 \quad \Rightarrow \quad x = 500$.
\end{itemize}
\textbf{Le pivot de Gauss confirme exactement les mêmes prix :} $x=500$, $y=200$, $z=1\,500$.

\textbf{Étape 4 : cohérence des prix.}
\begin{itemize}
\item Cahier (500 FCFA) $>$ Stylo (200 FCFA) : cohérent.
\item Rame (1\,500 FCFA) $>$ Cahier (500 FCFA) : cohérent (une rame = 500 feuilles).
\item Vérification des bordereaux :
\begin{align*}
\text{Commande 1 : } & 100 \times 500 + 50 \times 200 + 20 \times 1\,500 = 50\,000 + 10\,000 + 30\,000 = 90\,000 \quad \checkmark \\
\text{Commande 2 : } & 60 \times 500 + 100 \times 200 + 10 \times 1\,500 = 30\,000 + 20\,000 + 15\,000 = 65\,000 \quad \checkmark \\
\text{Commande 3 : } & 40 \times 500 + 30 \times 200 + 40 \times 1\,500 = 20\,000 + 6\,000 + 60\,000 = 86\,000 \quad \checkmark
\end{align*}
\end{itemize}

\textbf{Étapes 5 et 6 : inéquation budgétaire.}
Dépense prévue pour les stylos et rames (prix nouveaux du 2\textsuperscript{e} trimestre) :
\[
100 \times 500 + 20 \times 2\,000 = 50\,000 + 40\,000 = 90\,000 \text{ FCFA}.
\]
Soit $n$ le nombre de cahiers à commander (prix unitaire $x = 500$ FCFA). Le budget total ne doit pas dépasser 150\,000 FCFA :
\[
500n + 90\,000 \le 150\,000
\]
Résolution :
\[
500n \le 60\,000 \quad \Longleftrightarrow \quad n \le \frac{60\,000}{500} = 120.
\]
Contrainte du minimum : $n \ge 80$.
L'ensemble des solutions entières est $\{ n \in \mathbb{N} \mid 80 \le n \le 120 \}$.
La contrainte des 80 cahiers minimum est \textbf{satisfaisable} (l'intervalle n'est pas vide).
La commande optimale (maximisant le nombre de cahiers) est $n = 120$.

\textbf{Étape 7 : rapport (modèle de rédaction).}
\begin{quote}
\textbf{Objet :} Proposition de commande pour le 2\textsuperscript{e} trimestre — Rapport du trésorier.

Monsieur le Président,

À partir des trois bordereaux du premier trimestre, nous avons posé et résolu un système de trois équations à trois inconnues par la méthode de substitution, résultat confirmé par le pivot de Gauss. Les prix unitaires retrouvés sont :
\begin{itemize}
\item Cahier 200 pages : 500 FCFA
\item Stylo bleu : 200 FCFA
\item Ramette A4 : 1\,500 FCFA
\end{itemize}
Pour le deuxième trimestre, le bureau a arrêté l'achat de 100 stylos à 500 FCFA et 20 rames à 2\,000 FCFA, soit un engagement de 90\,000 FCFA. Le budget restant pour les cahiers est de 60\,000 FCFA. L'inéquation $500n + 90\,000 \le 150\,000$ donne $n \le 120$. La contrainte d'au moins 80 cahiers est respectée.

Nous proposons de commander \textbf{120 cahiers}. La dépense totale s'élèvera à $120 \times 500 + 90\,000 = 150\,000$ FCFA, soit exactement le budget, sans dépassement ni reliquat.

Cordialement, \hfill Le Trésorier
\end{quote}
\end{exoresolu}

\begin{attention}[Pièges fréquents rencontrés pendant l'activité]
\begin{itemize}
\item Oublier de simplifier les équations (division par 10) avant de substituer : les calculs deviennent vite ingérables et sources d'erreurs.
\item Se tromper de signe en isolant $z$ : $z = 6\,500 - 6x - 10y$, et non $z = 6\,500 + 6x + 10y$.
\item Conclure $n \le 120$ sans vérifier la contrainte $n \ge 80$ : une solution mathématique doit toujours être confrontée aux contraintes concrètes du problème.
\item Proposer $n = 120{,}5$ ou une valeur non entière : le nombre de cahiers est un entier naturel.
\item Ne pas vérifier la cohérence des prix trouvés avec les montants initiaux (test des trois commandes).
\end{itemize}
\end{attention}

\courssub{Bilan de l'activité}

Cette activité a permis de mettre en évidence que :
\begin{itemize}
\item un même problème concret (retrouver des prix unitaires) se traduit naturellement par un système $3\times 3$, et que les deux méthodes officielles du programme — \cle{substitution} et \cle{pivot de Gauss} — conduisent au même triplet solution, ce qui constitue un moyen d'auto-vérification puissant ;
\item le pivot de Gauss est plus systématique et plus sûr dès que les coefficients sont grands, tandis que la substitution est rapide lorsqu'une inconnue a un coefficient égal à 1 (ou $-1$) ;
\item une contrainte budgétaire se modélise par une \cle{inéquation}, dont l'ensemble des solutions doit être interprété dans le contexte (valeurs entières, bornes minimales et maximales) ;
\item la résolution d'un problème ne s'achève pas au calcul : il faut \textbf{rédiger une conclusion justifiée}, exploitable par un non-mathématicien, compétence essentielle dans les métiers techniques et de gestion.
\end{itemize}

\begin{savaistu}[Les coopératives scolaires au Cameroun]
Les coopératives scolaires, très présentes dans les lycées techniques camerounais, apprennent aux élèves la gestion d'une petite entreprise : achats, stocks, trésorerie, bénéfices réinvestis. Derrière chaque bordereau de commande se cachent des systèmes d'équations : les mathématiques de la Terminale technique sont déjà des outils de gestion professionnelle.
\end{savaistu}

\newpage

\coursec{Méthodes et exercices résolus}

\coursec{Équations, inéquations, systèmes}

\courssub{Équations du second degré}

\begin{definition}[Équation du second degré]
Une \cle{équation du second degré} à une inconnue $x$ s'écrit sous la forme canonique :
\[
ax^2 + bx + c = 0 \qquad (a \neq 0,\; a,b,c \in \mathbb{R}).
\]
Le nombre $\Delta = b^2 - 4ac$ s'appelle le \cle{discriminant} de l'équation.
\end{definition}

\begin{propriete}[Résolution d'une équation du second degré]
Soit $ax^2+bx+c=0$ ($a \neq 0$) et $\Delta = b^2-4ac$.
\begin{itemize}
  \item Si $\Delta < 0$ : l'équation n'admet \textbf{aucune solution réelle} ($S = \varnothing$).
  \item Si $\Delta = 0$ : l'équation admet une \textbf{solution double} $x_0 = -\dfrac{b}{2a}$.
  \item Si $\Delta > 0$ : l'équation admet \textbf{deux solutions distinctes}
  \[
  x_1 = \dfrac{-b - \sqrt{\Delta}}{2a}, \qquad x_2 = \dfrac{-b + \sqrt{\Delta}}{2a}.
  \]
\end{itemize}
\emph{Preuve exigible au Probatoire :} factorisation $a\left[\left(x+\frac{b}{2a}\right)^2 - \frac{\Delta}{4a^2}\right]=0$ et discussion du signe de $\Delta$.
\end{propriete}

\begin{methode}[Résoudre une équation du second degré $ax^2+bx+c=0$ ($a \neq 0$)]
\begin{enumerate}
\item \textbf{Mettre sous forme canonique} (second membre nul) et identifier $a$, $b$, $c$.
\item \textbf{Calculer le discriminant} $\Delta = b^2 - 4ac$.
\item \textbf{Étudier le signe de $\Delta$} et appliquer la propriété ci-dessus.
\item \textbf{Simplifier} les expressions (racines carrées, fractions) si possible.
\item \textbf{Rédiger la conclusion} : $S = \{ \dots \}$ ou $S = \varnothing$.
\end{enumerate}
\end{methode}

\begin{aretenir}[Équation bicarrée]
Une équation de la forme $ax^4+bx^2+c=0$ ($a \neq 0$) se résout par le changement de variable $X = x^2$ ($X \ge 0$).
On résout $aX^2+bX+c=0$ ; pour chaque solution positive $X_i$, on obtient $x = \pm\sqrt{X_i}$.
\end{aretenir}

\begin{exoresolu}[Équation du second degré et équation bicarrée]
Résoudre dans $\mathbb{R}$ les équations suivantes :
\begin{enumerate}
  \item $3x^2 - 7x + 2 = 0$
  \item $x^4 - 13x^2 + 36 = 0$
\end{enumerate}
\tcblower
\textbf{1.} L'équation est sous forme canonique : $a=3$, $b=-7$, $c=2$.
$\Delta = (-7)^2 - 4 \times 3 \times 2 = 49 - 24 = 25 = 5^2 > 0$.
Deux solutions distinctes :
$x_1 = \dfrac{7 - 5}{2 \times 3} = \dfrac{2}{6} = \dfrac{1}{3}$,
$x_2 = \dfrac{7 + 5}{6} = \dfrac{12}{6} = 2$.
$S = \left\{ \dfrac{1}{3} ; 2 \right\}$.

\textbf{2.} Équation bicarrée. On pose $X = x^2$ ($X \ge 0$).
$X^2 - 13X + 36 = 0$.
$\Delta_X = 169 - 144 = 25$. $\sqrt{\Delta_X} = 5$.
$X_1 = \dfrac{13 - 5}{2} = 4$, $X_2 = \dfrac{13 + 5}{2} = 9$.
Les deux sont positives, on revient à $x$ :
$x^2 = 4 \Rightarrow x = \pm 2$ ; $x^2 = 9 \Rightarrow x = \pm 3$.
$S = \{-3 ; -2 ; 2 ; 3\}$.
\end{exoresolu}

\begin{exercice}[Entraînement : équations du second degré]
Résoudre dans $\mathbb{R}$ :
\begin{enumerate}
  \item $2x^2 - 5x - 3 = 0$
  \item $x^2 + 4x + 7 = 0$
  \item $4x^4 - 17x^2 + 4 = 0$
  \item $(x-1)(2x+3) = 5$
\end{enumerate}
\end{exercice}

\courssub{Signe du trinôme et inéquations du second degré}

\begin{propriete}[Signe du trinôme $ax^2+bx+c$]
Soit $P(x) = ax^2+bx+c$ ($a \neq 0$) de discriminant $\Delta$ et de racines $x_1 \le x_2$ (si $\Delta \ge 0$).
\begin{itemize}
  \item \textbf{$\Delta < 0$} : $P(x)$ a le \textbf{signe constant de $a$} sur $\mathbb{R}$.
  \item \textbf{$\Delta = 0$} : $P(x)$ a le signe de $a$ sur $\mathbb{R}$ et s'annule en $x_0 = -\frac{b}{2a}$.
  \item \textbf{$\Delta > 0$} : $P(x) = a(x-x_1)(x-x_2)$.
    \begin{itemize}
      \item $P(x)$ a le signe de $a$ \textbf{en dehors} de l'intervalle $[x_1; x_2]$ (c'est-à-dire sur $]-\infty; x_1[ \cup ]x_2; +\infty[$).
      \item $P(x)$ a le signe \textbf{contraire} à $a$ \textbf{à l'intérieur} de l'intervalle $]x_1; x_2[$.
      \item $P(x) = 0$ pour $x = x_1$ ou $x = x_2$.
    \end{itemize}
\end{itemize}
\end{propriete}

\begin{methode}[Résoudre une inéquation du second degré $ax^2+bx+c > 0$ (ou $\ge, <, \le$)]
\begin{enumerate}
\item \textbf{Ramener le second membre à 0} et factoriser si possible (ou utiliser $\Delta$).
\item \textbf{Trouver les racines} $x_1 \le x_2$ (réelles) de l'équation associée $ax^2+bx+c=0$.
\item \textbf{Construire un tableau de signes} (obligatoire au Probatoire pour justifier) :
  \begin{itemize}
    \item Colonnes : intervalles et valeurs racines.
    \item Ligne pour le trinôme (ou ses facteurs).
    \item Appliquer la règle du signe (signe de $a$ à $+\infty$).
  \end{itemize}
\item \textbf{Lire la solution} selon l'inégalité (strict ou large) et conclure par $S = \dots$
\end{enumerate}
\end{methode}

\begin{exoresolu}[Inéquation du second degré et produit de facteurs]
Résoudre dans $\mathbb{R}$ :
\begin{enumerate}
  \item $-2x^2 + 5x + 3 \ge 0$
  \item $(2x-1)(x+3) < 0$
\end{enumerate}
\tcblower
\textbf{1.} $a=-2$, $b=5$, $c=3$. $\Delta = 25 - 4(-2)(3) = 25 + 24 = 49$. $\sqrt{\Delta}=7$.
Racines : $x_1 = \dfrac{-5 - 7}{-4} = 3$, $x_2 = \dfrac{-5 + 7}{-4} = -\dfrac{1}{2}$.
Ordre : $-\frac{1}{2} < 3$. $a < 0 \Rightarrow$ trinôme $\ge 0$ \textbf{entre} les racines.
\begin{center}
\begin{tabular}{|c|ccccc|}\hline
$x$ & $]-\infty; -\frac{1}{2}[$ & $-\frac{1}{2}$ & $]-\frac{1}{2}; 3[$ & $3$ & $]3; +\infty[$ \\ \hline
$-2x^2+5x+3$ & $-$ & $0$ & $+$ & $0$ & $-$ \\ \hline
\end{tabular}
\end{center}
$S = \left[ -\dfrac{1}{2} ; 3 \right]$.

\textbf{2.} Forme factorisée. Racines : $x = \frac{1}{2}$ et $x = -3$. Coefficient directeur du produit développé : $2 \times 1 = 2 > 0$.
Le produit est $< 0$ (négatif) \textbf{entre} les racines.
\begin{center}
\begin{tabular}{|c|ccccc|}\hline
$x$ & $]-\infty; -3[$ & $-3$ & $]-3; \frac{1}{2}[$ & $\frac{1}{2}$ & $]\frac{1}{2}; +\infty[$ \\ \hline
$2x-1$ & $-$ & $-$ & $-$ & $0$ & $+$ \\
$x+3$ & $-$ & $0$ & $+$ & $+$ & $+$ \\ \hline
$(2x-1)(x+3)$ & $+$ & $0$ & $-$ & $0$ & $+$ \\ \hline
\end{tabular}
\end{center}
$S = \left] -3 ; \dfrac{1}{2} \right[$.
\end{exoresolu}

\begin{attention}[Pièges classiques sur les inéquations du second degré]
\begin{enumerate}
\item \textbf{Oublier le signe de $a$} : croire que $ax^2+bx+c > 0$ donne toujours les valeurs \textbf{extérieures} aux racines.
\item \textbf{Multiplier une inéquation par une expression contenant l'inconnue} sans discussion de signe (ex: $\frac{x+1}{x-2} > 0 \not\Rightarrow x+1 > 0$). \textbf{Méthode sûre :} tableau de signes du numérateur et du dénominateur.
\item \textbf{Ne pas simplifier} les fractions ou les racines dans les expressions finales des racines.
\item \textbf{Oublier l'étude du discriminant} avant d'écrire les formules des racines.
\end{enumerate}
\end{attention}

\begin{exercice}[Entraînement : inéquations du second degré]
Résoudre dans $\mathbb{R}$ :
\begin{enumerate}
  \item $x^2 - 4x - 5 < 0$
  \item $-3x^2 + 6x - 2 \ge 0$
  \item $\dfrac{x-2}{x+1} \le 0$
  \item $(3x+1)^2 - 4 \le 0$
\end{enumerate}
\end{exercice}

\courssub{Inéquations de degré 3}

\begin{propriete}[Règle des racines multiples]
Dans un tableau de signes, pour un facteur $(x - \alpha)^k$ :
\begin{itemize}
  \item Si $k$ est \textbf{impair} : le signe \textbf{change} en $\alpha$.
  \item Si $k$ est \textbf{pair} : le signe \textbf{ne change pas} en $\alpha$ (le facteur garde le signe de $(x-\alpha)^k \ge 0$).
\end{itemize}
\end{propriete}

\begin{methode}[Résoudre une inéquation polynomiale de degré 3 $P(x) > 0$]
\begin{enumerate}
\item \textbf{Factoriser $P(x)$} au maximum (recherche d'une racine évidente $\pm 1, \pm 2, \dots$ divisateurs du terme constant, puis division euclidienne / factorisation par groupe).
\item \textbf{Identifier toutes les racines réelles} (simples ou multiples) et les classer par ordre croissant.
\item \textbf{Construire le tableau de signes} complet :
  \begin{itemize}
    \item Une ligne par facteur (linéaire ou trinôme sans racine réelle).
    \item Une ligne pour $P(x)$ (produit des signes).
    \item Appliquer la règle des racines multiples.
  \end{itemize}
\item \textbf{Lire les intervalles} solutions selon l'inégalité.
\item \textbf{Conclure} par l'ensemble solution $S$.
\end{enumerate}
\end{methode}

\begin{exoresolu}[Inéquation de degré 3 factorisable]
Résoudre dans $\mathbb{R}$ : $x^3 - 4x^2 - 7x + 10 > 0$.
\tcblower
On cherche une racine évidente : $P(1) = 1 - 4 - 7 + 10 = 0$. Donc $x=1$ est racine.
Division de $P(x)$ par $(x-1)$ (ou identification) :
$P(x) = (x-1)(x^2 - 3x - 10)$.
Factorisation du trinôme : $\Delta = 9 + 40 = 49$. Racines : $\frac{3 \pm 7}{2} \Rightarrow 5$ et $-2$.
$P(x) = (x-1)(x-5)(x+2)$.
Racines classées : $-2 < 1 < 5$. Toutes simples (multiplicité 1 $\Rightarrow$ changement de signe).
\begin{center}
\begin{tabular}{|c|ccccccc|}\hline
$x$ & $]-\infty; -2[$ & $-2$ & $]-2; 1[$ & $1$ & $]1; 5[$ & $5$ & $]5; +\infty[$ \\ \hline
$x+2$ & $-$ & $0$ & $+$ & $+$ & $+$ & $+$ & $+$ \\
$x-1$ & $-$ & $-$ & $-$ & $0$ & $+$ & $+$ & $+$ \\
$x-5$ & $-$ & $-$ & $-$ & $-$ & $-$ & $0$ & $+$ \\ \hline
$P(x)$ & $-$ & $0$ & $+$ & $0$ & $-$ & $0$ & $+$ \\ \hline
\end{tabular}
\end{center}
$P(x) > 0$ sur $]-2; 1[$ et $]5; +\infty[$.
$S = ]-2; 1[ \cup ]5; +\infty[$.
\end{exoresolu}

\begin{experience}[Découverte : signe d'un polynôme de degré 3]
Soit $P(x) = (x+1)^2(x-3)$.
\begin{enumerate}
  \item Quel est le degré de $P$ ? Quel est son coefficient directeur ?
  \item Construire le tableau de signes de $P$ en notant la multiplicité de la racine $-1$.
  \item En déduire le signe de $P(x)$ sur $\mathbb{R}$.
  \item Vérifier par le calcul de $P(-2)$, $P(0)$, $P(4)$.
\end{enumerate}
\tcblower
\textbf{1.} Degré 3, coefficient directeur $+1$.
\textbf{2.} Racines : $-1$ (double, multiplicité 2), $3$ (simple).
\begin{center}
\begin{tabular}{|c|ccccc|}\hline
$x$ & $]-\infty; -1[$ & $-1$ & $]-1; 3[$ & $3$ & $]3; +\infty[$ \\ \hline
$(x+1)^2$ & $+$ & $0$ & $+$ & $+$ & $+$ \\
$x-3$ & $-$ & $-$ & $-$ & $0$ & $+$ \\ \hline
$P(x)$ & $-$ & $0$ & $-$ & $0$ & $+$ \\ \hline
\end{tabular}
\end{center}
\textbf{3.} $P(x) \le 0$ sur $]-\infty; 3]$, $P(x) > 0$ sur $]3; +\infty[$.
\textbf{4.} $P(-2)=(-1)^2(-5)=-5<0$, $P(0)=1^2(-3)=-3<0$, $P(4)=5^2(1)=25>0$. OK.
\end{experience}

\begin{exercice}[Entraînement : inéquations de degré 3]
Résoudre dans $\mathbb{R}$ :
\begin{enumerate}
  \item $x^3 - 3x^2 - 4x + 12 \le 0$
  \item $x^3 + 2x^2 - x - 2 > 0$
  \item $(x-1)^2(x+2) \ge 0$
  \item $2x^3 - 5x^2 - x + 6 < 0$ (on note que $x=2$ est racine)
\end{enumerate}
\end{exercice}

\courssub{Systèmes 3×3 : méthode de substitution}

\begin{definition}[Système linéaire 3×3]
Un \cle{système linéaire de trois équations à trois inconnues} $(x,y,z)$ s'écrit :
\[
\begin{cases}
a_{11}x + a_{12}y + a_{13}z = b_1 \\
a_{21}x + a_{22}y + a_{23}z = b_2 \\
a_{31}x + a_{32}y + a_{33}z = b_3
\end{cases}
\]
On cherche les triplets $(x;y;z) \in \mathbb{R}^3$ satisfaisant les trois équations.
\end{definition}

\begin{methode}[Résoudre un système linéaire 3×3 par substitution]
\begin{enumerate}
\item \textbf{Choisir une équation} où une inconnue a pour coefficient $\pm 1$ (idéalement) et \textbf{l'isoler} (ex: $z = \dots$).
\item \textbf{Substituer} cette expression dans les deux autres équations $\rightarrow$ système de \textbf{2 équations à 2 inconnues}.
\item \textbf{Résoudre ce système 2×2} (par combinaison ou substitution) pour trouver les valeurs des deux inconnues.
\item \textbf{Remonter} : remplacer ces valeurs dans l'expression de l'inconnue isolée à l'étape 1.
\item \textbf{Vérifier} les trois valeurs dans les trois équations initiales.
\item \textbf{Conclure} : $S = \{(x;y;z)\}$ (unique), $S = \varnothing$ (incompatible) ou $S$ infini (dépendant, paramétrer).
\end{enumerate}
\end{methode}

\begin{exoresolu}[Système 3×3 par substitution]
Résoudre :
$\begin{cases}
x + y + z = 6 & (E_1)\\
2x - y + 3z = 9 & (E_2)\\
3x + 2y - z = 4 & (E_3)
\end{cases}$
\tcblower
D'après $(E_1)$ : $z = 6 - x - y$.
Substitution dans $(E_2)$ :
$2x - y + 3(6 - x - y) = 9 \iff 2x - y + 18 - 3x - 3y = 9 \iff -x - 4y = -9 \iff x + 4y = 9 \quad (E'_2)$.
Substitution dans $(E_3)$ :
$3x + 2y - (6 - x - y) = 4 \iff 3x + 2y - 6 + x + y = 4 \iff 4x + 3y = 10 \quad (E'_3)$.
Système 2×2 :
$\begin{cases} x + 4y = 9 \\ 4x + 3y = 10 \end{cases}$
De $(E'_2)$ : $x = 9 - 4y$.
Dans $(E'_3)$ : $4(9 - 4y) + 3y = 10 \iff 36 - 16y + 3y = 10 \iff -13y = -26 \iff y = 2$.
$x = 9 - 4(2) = 1$.
$z = 6 - 1 - 2 = 3$.
Vérification : $(E_2) : 2 - 2 + 9 = 9$ OK. $(E_3) : 3 + 4 - 3 = 4$ OK.
$S = \{(1; 2; 3)\}$.
\end{exoresolu}

\begin{exercice}[Entraînement : substitution]
Résoudre par substitution :
\begin{enumerate}
  \item $\begin{cases} x+y+z=3 \\ 2x-y+z=0 \\ x+2y-z=1 \end{cases}$
  \item $\begin{cases} x+2y-z=4 \\ 3x-y+2z=1 \\ 2x+y+z=5 \end{cases}$
\end{enumerate}
\end{exercice}

\courssub{Systèmes 3×3 : pivot de Gauss}

\begin{methode}[Résoudre un système linéaire 3×3 par pivot de Gauss (matrice augmentée)]
\begin{enumerate}
\item \textbf{Écrire la matrice augmentée} $\left( \begin{array}{ccc|c} a_{11} \& a_{12} \& a_{13} \& b_1 \\ a_{21} \& a_{22} \& a_{23} \& b_2 \\ a_{31} \& a_{32} \& a_{33} \& b_3 \end{array} \right)$.
\item \textbf{Mettre en forme échelonnée} (triangulaire supérieure) par \cle{opérations élémentaires sur les lignes (OEL)} :
  \begin{enumerate}
    \item Permuter deux lignes ($L_i \leftrightarrow L_j$).
    \item Multiplier une ligne par un non-nul ($L_i \leftarrow k L_i$).
    \item Ajouter un multiple d'une ligne à une autre ($L_i \leftarrow L_i + k L_j$).
  \end{enumerate}
  Objectif : obtenir des zéros sous la diagonale principale. Les pivots (coefficients diagonaux) doivent être non nuls.
\item \textbf{Analyser la dernière ligne} :
  \begin{itemize}
    \item $(0\;0\;0 \mid c \neq 0)$ : Système \textbf{incompatible} ($S = \varnothing$).
    \item $(0\;0\;0 \mid 0)$ : \textbf{Infinité de solutions} (paramétrer l'inconnue libre).
    \item Pivot non nul : \textbf{Solution unique}.
  \end{itemize}
\item \textbf{Remonter} (substitution à rebours) pour trouver $z$, puis $y$, puis $x$.
\item \textbf{Conclure}.
\end{enumerate}
\end{methode}

\begin{exoresolu}[Système 3×3 par pivot de Gauss]
Résoudre par la méthode du pivot de Gauss :
$\begin{cases}
x + 2y - z = 3 & (L_1)\\
2x - y + 2z = 6 & (L_2)\\
3x + y + z = 8 & (L_3)
\end{cases}$
\tcblower
Matrice augmentée :
$\left( \begin{array}{ccc|c}
1 & 2 & -1 & 3 \\
2 & -1 & 2 & 6 \\
3 & 1 & 1 & 8
\end{array} \right)$

$L_2 \leftarrow L_2 - 2L_1$ : $\left( \begin{array}{ccc|c} 1 & 2 & -1 & 3 \\ 0 & -5 & 4 & 0 \\ 3 & 1 & 1 & 8 \end{array} \right)$
$L_3 \leftarrow L_3 - 3L_1$ : $\left( \begin{array}{ccc|c} 1 & 2 & -1 & 3 \\ 0 & -5 & 4 & 0 \\ 0 & -5 & 4 & -1 \end{array} \right)$
$L_3 \leftarrow L_3 - L_2$ : $\left( \begin{array}{ccc|c} 1 & 2 & -1 & 3 \\ 0 & -5 & 4 & 0 \\ 0 & 0 & 0 & -1 \end{array} \right)$

Dernière ligne : $0x + 0y + 0z = -1$, impossible.
Le système est \textbf{incompatible}.
$S = \varnothing$.
\end{exoresolu}

\begin{exoresolu}[Système 3×3 : cas indéterminé (infini de solutions)]
Résoudre par Gauss :
$\begin{cases}
x + y + z = 3 \\
2x + 2y + 2z = 6 \\
x - y + z = 1
\end{cases}$
\tcblower
Matrice :
$\left( \begin{array}{ccc|c} 1 & 1 & 1 & 3 \\ 2 & 2 & 2 & 6 \\ 1 & -1 & 1 & 1 \end{array} \right)$
$L_2 \leftarrow L_2 - 2L_1$ : $\left( \begin{array}{ccc|c} 1 & 1 & 1 & 3 \\ 0 & 0 & 0 & 0 \\ 1 & -1 & 1 & 1 \end{array} \right)$
$L_3 \leftarrow L_3 - L_1$ : $\left( \begin{array}{ccc|c} 1 & 1 & 1 & 3 \\ 0 & 0 & 0 & 0 \\ 0 & -2 & 0 & -2 \end{array} \right)$
$L_2 \leftrightarrow L_3$ : $\left( \begin{array}{ccc|c} 1 & 1 & 1 & 3 \\ 0 & -2 & 0 & -2 \\ 0 & 0 & 0 & 0 \end{array} \right)$
Ligne 2 : $-2y = -2 \Rightarrow y = 1$.
Ligne 1 : $x + 1 + z = 3 \Rightarrow x + z = 2 \Rightarrow x = 2 - z$.
$z$ est paramètre libre ($z \in \mathbb{R}$).
$S = \{(2-t; 1; t) \mid t \in \mathbb{R}\}$.
\end{exoresolu}

\begin{exercice}[Entraînement : pivot de Gauss]
Résoudre par la méthode du pivot de Gauss :
\begin{enumerate}
  \item $\begin{cases} 2x - y + z = 1 \\ x + y - z = 2 \\ 3x + 2z = 5 \end{cases}$
  \item $\begin{cases} x + 2y - z = 4 \\ 2x - y + 3z = 1 \\ 3x + y + 2z = 5 \end{cases}$
  \item $\begin{cases} x + y + z = 2 \\ 2x + 2y + 2z = 4 \\ x - y + z = 0 \end{cases}$
\end{enumerate}
\end{exercice}

\courssub{Modélisation et résolution de problèmes concrets}

\begin{methode}[Modéliser et résoudre un problème technique (équations/inéquations/systèmes)]
\begin{enumerate}
\item \textbf{Lire attentivement} : identifier les grandeurs connues, inconnues, les contraintes (unités, positivité, entiers).
\item \textbf{Choisir les variables} (souvent $x, y, z$ ou $x_1, x_2, x_3$) et \textbf{noter leur unité}.
\item \textbf{Traduire} chaque phrase/condition en langage mathématique (équation, inéquation, système).
\item \textbf{Résoudre} le modèle mathématique avec la méthode adaptée (2nd degré, système 3×3, inéquation).
\item \textbf{Interpréter} les solutions mathématiques dans le contexte réel (rejeter les solutions non physiques : négatives, non entières si comptage, etc.).
\item \textbf{Répondre précisément} à la question posée par une phrase complète.
\end{enumerate}
\end{methode}

\begin{exoresolu}[Problème de mélange / achats (Système 3×3)]
Un chef de chantier achète trois types de matériaux : ciment, sable, gravier.
\begin{itemize}
  \item 2 sacs de ciment, 3 m$^3$ de sable, 1 m$^3$ de gravier coûtent 185 000 F.
  \item 1 sac de ciment, 2 m$^3$ de sable, 2 m$^3$ de gravier coûtent 150 000 F.
  \item 3 sacs de ciment, 1 m$^3$ de sable, 3 m$^3$ de gravier coûtent 255 000 F.
\end{itemize}
Déterminer le prix unitaire de chaque matériau.
\tcblower
\textbf{Variables :} $x$ = prix sac ciment (en 1000 F), $y$ = prix m$^3$ sable, $z$ = prix m$^3$ gravier. ($x,y,z > 0$).
\textbf{Système :}
$\begin{cases}
2x + 3y + z = 185 & (1)\\
x + 2y + 2z = 150 & (2)\\
3x + y + 3z = 255 & (3)
\end{cases}$

\textbf{Résolution par Gauss :}
Matrice : $\left( \begin{array}{ccc|c} 2 & 3 & 1 & 185 \\ 1 & 2 & 2 & 150 \\ 3 & 1 & 3 & 255 \end{array} \right)$
$L_1 \leftrightarrow L_2$ (pivot 1) : $\left( \begin{array}{ccc|c} 1 & 2 & 2 & 150 \\ 2 & 3 & 1 & 185 \\ 3 & 1 & 3 & 255 \end{array} \right)$
$L_2 \leftarrow L_2 - 2L_1$ : $\left( \begin{array}{ccc|c} 1 & 2 & 2 & 150 \\ 0 & -1 & -3 & -115 \\ 3 & 1 & 3 & 255 \end{array} \right)$
$L_3 \leftarrow L_3 - 3L_1$ : $\left( \begin{array}{ccc|c} 1 & 2 & 2 & 150 \\ 0 & -1 & -3 & -115 \\ 0 & -5 & -3 & -195 \end{array} \right)$
$L_2 \leftarrow -L_2$ : $\left( \begin{array}{ccc|c} 1 & 2 & 2 & 150 \\ 0 & 1 & 3 & 115 \\ 0 & -5 & -3 & -195 \end{array} \right)$
$L_3 \leftarrow L_3 + 5L_2$ : $\left( \begin{array}{ccc|c} 1 & 2 & 2 & 150 \\ 0 & 1 & 3 & 115 \\ 0 & 0 & 12 & 380 \end{array} \right)$

Remontée :
$12z = 380 \Rightarrow z = \frac{380}{12} = \frac{95}{3} \approx 31,67$ (en 1000 F) $\Rightarrow 31\,667$ F.
$y + 3z = 115 \Rightarrow y = 115 - 95 = 20$ (en 1000 F) $\Rightarrow 20\,000$ F.
$x + 2y + 2z = 150 \Rightarrow x + 40 + \frac{190}{3} = 150 \Rightarrow x = 110 - \frac{190}{3} = \frac{140}{3} \approx 46,67$ (en 1000 F) $\Rightarrow 46\,667$ F.

\textbf{Vérification :} (1) $2\times\frac{140}{3} + 60 + \frac{95}{3} = \frac{280+95}{3}+60 = 125+60=185$ OK.
\textbf{Réponse :} Le sac de ciment coûte 46 667 F, le m$^3$ de sable 20 000 F, le m$^3$ de gravier 31 667 F.
\end{exoresolu}

\begin{exoresolu}[Problème géométrique (Inéquation du second degré)]
Un rectangle a une longueur supérieure de 5 cm à sa largeur. Son aire doit être comprise entre 50 cm$^2$ et 150 cm$^2$. Déterminer les valeurs possibles de la largeur.
\tcblower
\textbf{Variables :} $x$ = largeur en cm ($x > 0$). Longueur = $x+5$.
Aire $A = x(x+5)$.
Contrainte : $50 \le x(x+5) \le 150$.
Système d'inéquations :
$\begin{cases} x^2 + 5x - 50 \ge 0 \\ x^2 + 5x - 150 \le 0 \end{cases}$

1) $x^2 + 5x - 50 \ge 0$. $\Delta = 25 + 200 = 225 = 15^2$. Racines : $\frac{-5 \pm 15}{2} \Rightarrow -10$ et $5$.
$a=1>0 \Rightarrow$ solution : $]-\infty; -10] \cup [5; +\infty[$.
Avec $x>0$ : $x \in [5; +\infty[$.

2) $x^2 + 5x - 150 \le 0$. $\Delta = 25 + 600 = 625 = 25^2$. Racines : $\frac{-5 \pm 25}{2} \Rightarrow -15$ et $10$.
$a=1>0 \Rightarrow$ solution : $[-15; 10]$.
Avec $x>0$ : $x \in ]0; 10]$.

Intersection : $[5; +\infty[ \cap ]0; 10] = [5; 10]$.
\textbf{Réponse :} La largeur doit être comprise entre 5 cm et 10 cm.
\end{exoresolu}

\begin{exercice}[Problèmes de modélisation]
\begin{enumerate}
  \item \textbf{Électricité :} Trois résistances $R_1, R_2, R_3$ sont montées en série. La tension totale est $U = 12$ V. On mesure $U_1 = 3$ V aux bornes de $R_1$ et $U_2 = 4$ V aux bornes de $R_2$. Le courant est $I = 0,5$ A. Calculer $R_1, R_2, R_3$ (Loi d'Ohm : $U = RI$).
  \item \textbf{Gestion de stock :} Un magasin vend trois articles A, B, C. Le stock initial total est 100 unités. On vend 2 A, 3 B, 1 C pour 185 € ; 1 A, 2 B, 2 C pour 150 € ; 3 A, 1 B, 3 C pour 255 €. Trouver le prix unitaire de A, B, C.
  \item \textbf{Trajectoire :} La hauteur $h$ (en m) d'un projectile en fonction du temps $t$ (en s) est $h(t) = -5t^2 + 20t + 2$. Pendant combien de temps le projectile est-il à une hauteur supérieure à 17 m ?
\end{enumerate}
\end{exercice}

\begin{attention}[Les 5 erreurs qui coûtent des points au Probatoire]
\begin{enumerate}
\item \textbf{Oublier l'étude du discriminant} avant d'appliquer les formules des racines (écrire $x = \frac{-b \pm \sqrt{\Delta}}{2a}$ sans dire $\Delta \ge 0$). $\rightarrow$ Perte de points « Raisonnement ».
\item \textbf{Confondre le signe du trinôme} : croire que $ax^2+bx+c > 0$ donne toujours les valeurs \textbf{extérieures} aux racines (oublier le signe de $a$). $\rightarrow$ Solution fausse.
\item \textbf{Multiplier une inéquation par une expression contenant l'inconnue} sans faire de discussion de signe (cas positif / cas négatif). Ex: $\frac{x+1}{x-2} > 0 \not\Rightarrow x+1 > 0$. $\rightarrow$ Erreur majeure, 0 point à l'exercice.
\item \textbf{Erreurs de signes dans le pivot de Gauss} : soustraire au lieu d'ajouter, oublier de changer les signes de toute la ligne, erreur sur le second membre. $\rightarrow$ Vérifier systématiquement en fin de résolution.
\item \textbf{Ne pas conclure / Ne pas répondre à la question} : donner $x=3$ au lieu de « La longueur est 3 m », ou oublier de rejeter la racine négative dans un problème géométrique. $\rightarrow$ Point « Conclusion / Communication » perdu.
\end{enumerate}
\end{attention}

\begin{savaistu}[Histoire et applications]
\textbf{Al-Khwarizmi (vers 780-850)}, mathématicien perse, a donné les premières méthodes systématiques de résolution des équations du second degré dans son ouvrage \emph{Al-Kitab al-mukhtasar fi hisab al-jabr wa'l-muqabala} (d'où le mot « algèbre »). Il classait les équations en six types et les résolvait par des raisonnements géométriques (complétion du carré).

\textbf{Carl Friedrich Gauss (1777-1855)} a formalisé la méthode d'élimination (pivot de Gauss) pour résoudre les systèmes linéaires, bien que des méthodes similaires existaient en Chine dès le 1er siècle av. J.-C. (\emph{Neuf chapitres sur l'art mathématique}). Cette méthode est la base de l'algèbre linéaire numérique utilisée aujourd'hui en ingénierie (calcul de structures, circuits électriques, modélisation 3D).

\textbf{Au Cameroun :} Les ingénieurs des Travaux Publics utilisent des systèmes 3×3 (et bien plus grands) pour calculer les efforts dans les poutres des ponts (méthode des déplacements) ou pour dimensionner les réseaux d'adduction d'eau (lois de Kirchhoff hydrauliques).
\end{savaistu}

\newpage

\coursec{Exercices}

\courssub{Je vérifie mes connaissances}

\begin{exercice}[QCM et Vrai/Faux justifié]
\textbf{1.} Soit l'équation $3x^2 - 5x + 2 = 0$. Son discriminant $\Delta$ vaut~:
\begin{enumerate}
\item $1$ \quad \item $49$ \quad \item $13$ \quad \item $-1$
\end{enumerate}
\textbf{2.} Le trinôme $f(x) = -2x^2 + 4x - 5$ est~:
\begin{enumerate}
\item toujours positif \quad \item toujours négatif \quad \item de signe variable \quad \item nul pour $x=1$
\end{enumerate}
\textbf{3.} Vrai ou Faux ? « Une inéquation du type $(x-2)(x+3) \ge 0$ a pour ensemble de solutions $[-3; 2]$. » Justifiez.
\textbf{4.} Vrai ou Faux ? « La méthode du pivot de Gauss consiste à remplacer une équation par la somme de cette équation et d'un multiple d'une autre. » Justifiez.
\end{exercice}

\begin{exercice}[Vocabulaire et définitions]
Complétez les phrases suivantes en utilisant le terme exact :
\begin{enumerate}
\item Pour l'équation $ax^2+bx+c=0$ ($a \neq 0$), l'expression $b^2-4ac$ s'appelle le \ldots
\item Si $\Delta > 0$, l'équation admet \ldots solutions réelles distinctes.
\item La forme \ldots d'un trinôme $ax^2+bx+c$ est $a\left[\left(x+\frac{b}{2a}\right)^2 - \frac{\Delta}{4a^2}\right]$.
\item Dans la résolution d'un système $3 \times 3$ par pivot de Gauss, on cherche à obtenir une matrice \ldots (forme échelonnée).
\item Le théorème qui permet d'affirmer qu'un polynôme $P(x)$ est divisible par $(x-a)$ si $P(a)=0$ s'appelle le \ldots
\end{enumerate}
\end{exercice}

\begin{exercice}[Calculs directs]
\begin{enumerate}
\item Calculez le discriminant de $2x^2 - 8x + 5 = 0$ et donnez ses solutions exactes.
\item Résolvez dans $\mathbb{R}$ : $x^2 - 9 = 0$.
\item Factorisez l'expression $x^3 - 3x^2 - 4x + 12$ en remarquant une racine évidente.
\item Soit le système : $\begin{cases} x + y + z = 6 \\ 2x - y + z = 3 \\ x + 2y - z = 2 \end{cases}$. Donnez la matrice augmentée associée.
\end{enumerate}
\end{exercice}

\begin{exercice}[Résolutions rapides]
Résolvez dans $\mathbb{R}$ :
\begin{enumerate}
\item $4x^2 - 12x + 9 = 0$
\item $x^2 - 4x - 5 < 0$
\item $\dfrac{x+1}{x-2} \le 0$ (utilisez un tableau de signes)
\item $x^3 - 4x \ge 0$
\end{enumerate}
\end{exercice}

\courssub{Je m'entraîne}

\begin{exercice}[Équations du second degré et bicarrées]
Résolvez dans $\mathbb{R}$ les équations suivantes :
\begin{enumerate}
\item $3x^2 - 7x + 2 = 0$
\item $2x^2 + 5x - 3 = 0$
\item $x^4 - 5x^2 + 4 = 0$ \quad (équation bicarrée)
\item $9x^4 - 6x^2 + 1 = 0$
\end{enumerate}
\end{exercice}

\begin{exercice}[Inéquations du second degré et signes]
\begin{enumerate}
\item Résolvez l'inéquation $2x^2 - 5x - 3 < 0$.
\item Résolvez l'inéquation $-x^2 + 6x - 9 \ge 0$.
\item Résolvez dans $\mathbb{R}$ : $\dfrac{x^2 - 4}{x - 1} \ge 0$. Présentez le tableau de signes complet.
\item Déterminez l'ensemble des réels $x$ tels que $x^2 + 2x + 3 > 0$.
\end{enumerate}
\end{exercice}

\begin{exercice}[Inéquations de degré 3]
Soit le polynôme $P(x) = x^3 - 4x^2 + x + 6$.
\begin{enumerate}
\item Vérifiez que $2$ est une racine de $P$.
\item Factorisez $P(x)$ sous la forme $(x-2)Q(x)$ où $Q$ est un trinôme du second degré.
\item Factorisez complètement $P(x)$.
\item Résolvez l'inéquation $P(x) \le 0$ à l'aide d'un tableau de signes.
\end{enumerate}
\end{exercice}

\begin{exercice}[Systèmes 3×3 : Substitution et Pivot de Gauss]
Résolvez le système suivant par les \textbf{deux méthodes} et comparez la rapidité de chacune :
\[
(S) : \begin{cases}
x + y - z = 0 \\
2x - y + z = 5 \\
x + 2y + z = 8
\end{cases}
\]
\begin{enumerate}
\item Méthode par substitution (exprimez $z$ par la première, substituez dans les deux autres, etc.).
\item Méthode du pivot de Gauss (écrivez la matrice augmentée, pivotez pour obtenir un échelon, remontez).
\item Vérifiez que le triplet solution est le même. Quelle méthode vous a semblé la plus adaptée ici ? Justifiez.
\end{enumerate}
\end{exercice}

\courssub{Je me prépare au Probatoire}

\begin{exercice}[Problème d'atelier : Coûts de production]
\label{exo:atelier}
Un atelier de menuiserie fabrique trois types de meubles : des chaises (C), des tables (T) et des bancs (B).
Le coût de revient unitaire (en FCFA) de chaque meuble est inconnu : $x$ pour une chaise, $y$ pour une table, $z$ pour un banc.
Sur trois semaines de production, l'atelier a enregistré les bilans suivants :
\begin{itemize}
\item Semaine 1 : 10 chaises, 5 tables, 3 bancs $\rightarrow$ Coût total : 1\,250\,000 FCFA.
\item Semaine 2 : 8 chaises, 7 tables, 2 bancs $\rightarrow$ Coût total : 1\,310\,000 FCFA.
\item Semaine 3 : 12 chaises, 4 tables, 5 bancs $\rightarrow$ Coût total : 1\,580\,000 FCFA.
\end{itemize}
\begin{enumerate}
\item Modélisez la situation par un système de trois équations à trois inconnues $(S)$.
\item Résolvez $(S)$ par la méthode du pivot de Gauss (présentez les matrices successives).
\item Donnez le coût unitaire de chaque meuble.
\item L'atelier reçoit une commande de 20 chaises, 15 tables et 10 bancs. Calculez le coût total de cette commande.
\item \textbf{Analyse critique :} Si le prix de vente d'une chaise est fixé à 60\,000 FCFA, d'une table à 120\,000 FCFA et d'un banc à 80\,000 FCFA, l'atelier réalise-t-il un bénéfice sur cette commande ? Calculez le bénéfice ou la perte.
\end{enumerate}
\end{exercice}

\begin{exercice}[Problème de géométrie : Dimensions d'un terrain]
\label{exo:terrain}
Un entrepreneur doit clôturer un terrain rectangulaire. Il dispose de 60 mètres de grillage pour le périmètre. Le cahier des charges impose que l'aire du terrain soit \textbf{au moins} de 200 m$^2$.
On note $x$ la longueur (en mètres) et $y$ la largeur (en mètres) du terrain.
\begin{enumerate}
\item Exprimez $y$ en fonction de $x$ à l'aide de la condition sur le périmètre.
\item Traduisez la condition sur l'aire par une inéquation en $x$ uniquement.
\item Résolvez cette inéquation dans $\mathbb{R}$.
\item En déduire les dimensions possibles (longueur et largeur) pour le terrain. Précisez si la longueur et la largeur sont interchangeables dans votre réponse.
\item L'entrepreneur souhaite minimiser la longueur de la diagonale du terrain. Quelle dimension doit-il choisir ? (On admet que pour un périmètre fixe, le rectangle de plus petite diagonale est le carré).
\end{enumerate}
\end{exercice}

\begin{exercice}[Problème d'électricité : Loi des nœuds et des mailles]
\label{exo:electricite}
Dans le circuit électrique ci-dessous, on cherche à déterminer les intensités $I_1$, $I_2$, $I_3$ (en Ampères) circulant dans trois branches.
\begin{center}
\begin{tikzpicture}[scale=0.9, >=Stealth]
% Nœuds
\node[circle,fill,inner sep=2pt,label=above:$A$] (A) at (0,2) {};
\node[circle,fill,inner sep=2pt,label=below:$B$] (B) at (0,0) {};
\node[circle,fill,inner sep=2pt,label=above:$C$] (C) at (4,2) {};
\node[circle,fill,inner sep=2pt,label=below:$D$] (D) at (4,0) {};
% Branches
\draw[thick] (A) -- (B) node[midway, left] {$E_1=12V$};
\draw[thick] (C) -- (D) node[midway, right] {$E_2=6V$};
\draw[thick] (A) -- (C) node[midway, above] {$R_1=2\Omega$};
\draw[thick] (B) -- (D) node[midway, below] {$R_2=4\Omega$};
\draw[thick] (A) -- (D) node[midway, above, sloped] {$R_3=1\Omega$};
% Courants
\draw[->, popblue, thick] (0.5,2.2) -- (1.5,2.2) node[midway, above] {$I_1$};
\draw[->, poporange, thick] (0.5,-0.2) -- (1.5,-0.2) node[midway, below] {$I_2$};
\draw[->, popgreen, thick] (0.2,1.8) -- (0.2,0.2) node[midway, left] {$I_3$};
% Loi des nœuds
\node[align=center] at (2,-1.5) {Schéma : Pont de Wheatstone simplifié};
\end{tikzpicture}
\end{center}
En appliquant la loi des nœuds au point $A$ et la loi des mailles aux deux mailles indépendantes (gauche et droite), on obtient le système :
\[
\begin{cases}
I_1 - I_2 - I_3 = 0 & \text{(Nœud A)} \\
2I_1 + I_3 = 12 & \text{(Maille gauche : } E_1 = R_1 I_1 + R_3 I_3\text{)} \\
4I_2 - I_3 = 6 & \text{(Maille droite : } E_2 = R_2 I_2 - R_3 I_3\text{)}
\end{cases}
\]
\begin{enumerate}
\item Vérifiez que le système s'écrit sous la forme matricielle $A X = B$ en précisant $A$, $X$, $B$.
\item Résolvez ce système par la méthode du pivot de Gauss.
\item Donnez les valeurs des intensités $I_1$, $I_2$, $I_3$. Les sens choisis initialement sur le schéma sont-ils les bons ? Justifiez.
\item Calculez la puissance dissipée par la résistance $R_3$ (formule $P = R I^2$).
\end{enumerate}
\end{exercice}

\newpage

\coursec{Corrigés des exercices}

\begin{corrige}[Exercice 1 — QCM et Vrai/Faux justifié]
\textbf{1.} $\Delta = b^2 - 4ac = (-5)^2 - 4 \times 3 \times 2 = 25 - 24 = 1 > 0$. L'équation admet \textbf{deux solutions réelles distinctes}. \textbf{Réponse : item 1.}

\textbf{2.} $\Delta = 4^2 - 4 \times (-2) \times (-5) = 16 - 40 = -24 < 0$. Le trinôme ne s'annule pas et garde le signe de $a = -2 < 0$ : il est \textbf{toujours négatif} sur $\mathbb{R}$. \textbf{Réponse : item 2.}

\textbf{3.} \textbf{Faux.} Le trinôme $(x-2)(x+3) = x^2+x-6$ a pour racines $-3$ et $2$, avec $a = 1 > 0$ : il est positif \emph{à l'extérieur} des racines. L'ensemble des solutions de $(x-2)(x+3) \ge 0$ est $S = ]-\infty\,; -3] \cup [2\,; +\infty[$, et non $[-3\,;2]$.

\textbf{4.} \textbf{Vrai.} Le \cle{pivot de Gauss} repose précisément sur des opérations élémentaires du type $L_i \leftarrow L_i + \lambda L_j$ (remplacer une ligne par la somme de cette ligne et d'un multiple d'une autre), qui conservent les solutions du système.
\end{corrige}

\begin{corrige}[Exercice 2 — Vocabulaire et définitions]
\begin{enumerate}
\item L'expression $b^2 - 4ac$ s'appelle le \cle{discriminant} (noté $\Delta$).
\item Si $\Delta > 0$, l'équation admet \textbf{deux} solutions réelles distinctes.
\item La forme \cle{canonique} d'un trinôme est $a\left[\left(x+\dfrac{b}{2a}\right)^2 - \dfrac{\Delta}{4a^2}\right]$.
\item On cherche à obtenir une matrice \textbf{triangulaire supérieure} (forme échelonnée).
\item Il s'agit du \cle{théorème de factorisation} (conséquence du théorème de Bézout) : $P(a) = 0 \iff P(x)$ est divisible par $(x-a)$.
\end{enumerate}
\end{corrige}

\begin{corrige}[Exercice 3 — Calculs directs]
\textbf{1.} $2x^2 - 8x + 5 = 0$. $\Delta = (-8)^2 - 4 \times 2 \times 5 = 64 - 40 = 24 > 0$. Deux solutions :
\[ x_1 = \frac{8 - \sqrt{24}}{4} = \frac{8 - 2\sqrt{6}}{4} = 2 - \frac{\sqrt{6}}{2} \quad ; \quad x_2 = 2 + \frac{\sqrt{6}}{2}. \]

\textbf{2.} $x^2 - 9 = 0 \iff (x-3)(x+3) = 0 \iff x = 3$ ou $x = -3$. $S = \{-3\,; 3\}$.

\textbf{3.} $x = 3$ est racine évidente : $27 - 27 - 12 + 12 = 0$. Par groupement :
\[ x^3 - 3x^2 - 4x + 12 = x^2(x-3) - 4(x-3) = (x-3)(x^2 - 4) = (x-3)(x-2)(x+2). \]

\textbf{4.} Matrice augmentée associée au système :
\[ \begin{pmatrix} 1 & 1 & 1 & \big| & 6 \\ 2 & -1 & 1 & \big| & 3 \\ 1 & 2 & -1 & \big| & 2 \end{pmatrix}. \]
\end{corrige}

\begin{corrige}[Exercice 4 — Résolutions rapides]
\textbf{1.} $4x^2 - 12x + 9 = (2x-3)^2 = 0 \iff x = \dfrac{3}{2}$ (racine double, $\Delta = 144 - 144 = 0$). $S = \left\{\dfrac{3}{2}\right\}$.

\textbf{2.} $x^2 - 4x - 5 < 0$. $\Delta = 16 + 20 = 36$, racines $x_1 = -1$, $x_2 = 5$. Comme $a = 1 > 0$, le trinôme est négatif \emph{entre} les racines : $S = ]-1\,; 5[$.

\textbf{3.} $\dfrac{x+1}{x-2} \le 0$. Zéro : $x = -1$ ; valeur interdite : $x = 2$.
\begin{center}\begin{tabular}{|c|ccccccc|}\hline $x$ & $-\infty$ & & $-1$ & & $2$ & & $+\infty$ \\ \hline $x+1$ & & $-$ & $0$ & $+$ & & $+$ & \\ \hline $x-2$ & & $-$ & & $-$ & $0$ & $+$ & \\ \hline quotient & & $+$ & $0$ & $-$ & $\|$ & $+$ & \\ \hline\end{tabular}\end{center}
Le quotient est négatif ou nul sur $S = [-1\,; 2[$ (la valeur $2$ est exclue).

\textbf{4.} $x^3 - 4x \ge 0 \iff x(x^2-4) \ge 0 \iff x(x-2)(x+2) \ge 0$. Racines : $-2, 0, 2$. Le signe de ce produit (coefficient dominant positif) donne :
\[ S = [-2\,; 0] \cup [2\,; +\infty[. \]
\end{corrige}

\begin{corrige}[Exercice 5 — Équations du second degré et bicarrées]
\textbf{1.} $3x^2 - 7x + 2 = 0$ : $\Delta = 49 - 24 = 25$, $x_1 = \dfrac{7-5}{6} = \dfrac{1}{3}$, $x_2 = \dfrac{7+5}{6} = 2$. $S = \left\{\dfrac{1}{3}\,; 2\right\}$.

\textbf{2.} $2x^2 + 5x - 3 = 0$ : $\Delta = 25 + 24 = 49$, $x_1 = \dfrac{-5-7}{4} = -3$, $x_2 = \dfrac{-5+7}{4} = \dfrac{1}{2}$. $S = \left\{-3\,; \dfrac{1}{2}\right\}$.

\textbf{3.} Posons $t = x^2$ ($t \ge 0$) : $t^2 - 5t + 4 = 0$, $\Delta = 25 - 16 = 9$, $t = 1$ ou $t = 4$. Donc $x^2 = 1$ ou $x^2 = 4$ : $S = \{-2\,; -1\,; 1\,; 2\}$.

\textbf{4.} Posons $t = x^2$ : $9t^2 - 6t + 1 = (3t-1)^2 = 0 \iff t = \dfrac{1}{3}$. Alors $x^2 = \dfrac{1}{3} \iff x = \pm \dfrac{1}{\sqrt{3}} = \pm \dfrac{\sqrt{3}}{3}$. $S = \left\{-\dfrac{\sqrt{3}}{3}\,; \dfrac{\sqrt{3}}{3}\right\}$.
\end{corrige}

\begin{corrige}[Exercice 6 — Inéquations du second degré et signes]
\textbf{1.} $2x^2 - 5x - 3 < 0$ : $\Delta = 25 + 24 = 49$, $x_1 = \dfrac{5-7}{4} = -\dfrac{1}{2}$, $x_2 = \dfrac{5+7}{4} = 3$. Comme $a = 2 > 0$, le trinôme est strictement négatif \emph{entre} les racines : $S = \left]-\dfrac{1}{2}\,; 3\right[$.

\textbf{2.} $-x^2 + 6x - 9 = -(x^2 - 6x + 9) = -(x-3)^2$. Or $-(x-3)^2 \ge 0 \iff (x-3)^2 \le 0 \iff x = 3$. $S = \{3\}$.

\textbf{3.} $\dfrac{x^2-4}{x-1} = \dfrac{(x-2)(x+2)}{x-1} \ge 0$. Zéros : $-2$ et $2$ ; valeur interdite : $1$.
\begin{center}\begin{tabular}{|c|ccccccccc|}\hline $x$ & $-\infty$ & & $-2$ & & $1$ & & $2$ & & $+\infty$ \\ \hline $x+2$ & & $-$ & $0$ & $+$ & & $+$ & & $+$ & \\ \hline $x-1$ & & $-$ & & $-$ & $0$ & $+$ & & $+$ & \\ \hline $x-2$ & & $-$ & & $-$ & & $-$ & $0$ & $+$ & \\ \hline quotient & & $-$ & $0$ & $+$ & $\|$ & $-$ & $0$ & $+$ & \\ \hline\end{tabular}\end{center}
Le quotient est positif ou nul sur $S = [-2\,; 1[ \cup [2\,; +\infty[$.

\textbf{4.} $x^2 - 2x + 3 > 0$. $\Delta = 4 - 12 = -8 < 0$ et $a = 1 > 0$ : le trinôme est strictement positif pour tout réel. $S = \mathbb{R}$.
\end{corrige}

\begin{corrige}[Exercice 7 — Inéquations de degré 3]
Soit $P(x) = x^3 - 4x^2 + x + 6$.

\textbf{1.} $P(2) = 8 - 16 + 2 + 6 = 0$ : $2$ est bien une racine de $P$.

\textbf{2.} On cherche $Q(x) = x^2 + bx + c$ tel que $P(x) = (x-2)(x^2+bx+c)$. Par identification (ou division euclidienne) : le terme constant donne $-2c = 6$, donc $c = -3$ ; le terme en $x^2$ donne $b - 2 = -4$, donc $b = -2$.
\[ P(x) = (x-2)(x^2 - 2x - 3). \]
Vérification : $(x-2)(x^2-2x-3) = x^3 - 2x^2 - 3x - 2x^2 + 4x + 6 = x^3 - 4x^2 + x + 6$. \checkmark

\textbf{3.} $x^2 - 2x - 3 = 0$ : $\Delta = 4 + 12 = 16$, racines $x = -1$ et $x = 3$. Donc factorisation complète :
\[ P(x) = (x-2)(x+1)(x-3). \]

\textbf{4.} Tableau de signes :
\begin{center}\begin{tabular}{|c|ccccccccc|}\hline $x$ & $-\infty$ & & $-1$ & & $2$ & & $3$ & & $+\infty$ \\ \hline $x+1$ & & $-$ & $0$ & $+$ & & $+$ & & $+$ & \\ \hline $x-2$ & & $-$ & & $-$ & $0$ & $+$ & & $+$ & \\ \hline $x-3$ & & $-$ & & $-$ & & $-$ & $0$ & $+$ & \\ \hline $P(x)$ & & $-$ & $0$ & $+$ & $0$ & $-$ & $0$ & $+$ & \\ \hline\end{tabular}\end{center}
$P(x) \le 0$ sur $S = ]-\infty\,; -1] \cup [2\,; 3]$.
\end{corrige}

\begin{corrige}[Exercice 8 — Systèmes 3×3 : Substitution et Pivot de Gauss]
Système :
\[ (S) : \begin{cases} x + y - z = 0 \\ 2x - y + z = 5 \\ x + 2y + z = 8 \end{cases} \]

\textbf{1. Méthode par substitution.} La première équation donne $z = x + y$. On substitue dans les deux autres :
\begin{align*}
2x - y + (x+y) &= 5 \iff 3x = 5 \iff x = \frac{5}{3}, \\
x + 2y + (x+y) &= 8 \iff 2x + 3y = 8 \iff \frac{10}{3} + 3y = 8 \iff y = \frac{14}{9}.
\end{align*}
Puis $z = x + y = \dfrac{15}{9} + \dfrac{14}{9} = \dfrac{29}{9}$. Solution : $\left(\dfrac{5}{3}\,; \dfrac{14}{9}\,; \dfrac{29}{9}\right)$.

\textbf{2. Méthode du pivot de Gauss.}
\[ \begin{pmatrix} 1 & 1 & -1 & \big| & 0 \\ 2 & -1 & 1 & \big| & 5 \\ 1 & 2 & 1 & \big| & 8 \end{pmatrix}
\xrightarrow[L_3 \leftarrow L_3 - L_1]{L_2 \leftarrow L_2 - 2L_1}
\begin{pmatrix} 1 & 1 & -1 & \big| & 0 \\ 0 & -3 & 3 & \big| & 5 \\ 0 & 1 & 2 & \big| & 8 \end{pmatrix}
\xrightarrow{L_3 \leftarrow 3L_3 + L_2}
\begin{pmatrix} 1 & 1 & -1 & \big| & 0 \\ 0 & -3 & 3 & \big| & 5 \\ 0 & 0 & 9 & \big| & 29 \end{pmatrix} \]
Remontée : $9z = 29 \Rightarrow z = \dfrac{29}{9}$ ; $-3y + 3 \times \dfrac{29}{9} = 5 \Rightarrow -3y = 5 - \dfrac{29}{3} = -\dfrac{14}{3} \Rightarrow y = \dfrac{14}{9}$ ; $x = z - y = \dfrac{15}{9} = \dfrac{5}{3}$.

\textbf{3. Vérification.} Les deux méthodes donnent le même triplet $\left(\dfrac{5}{3}\,; \dfrac{14}{9}\,; \dfrac{29}{9}\right)$. Contrôle dans $(S)$ : $\frac{15}{9}+\frac{14}{9}-\frac{29}{9} = 0$ ; $\frac{30}{9}-\frac{14}{9}+\frac{29}{9} = \frac{45}{9} = 5$ ; $\frac{15}{9}+\frac{28}{9}+\frac{29}{9} = \frac{72}{9} = 8$. \checkmark

\begin{attention}[Choix de la méthode]
Ici, la substitution est plus rapide car la première équation permet d'isoler $z$ immédiatement (coefficient $-1$). Le \cle{pivot de Gauss} reste plus systématique et plus sûr lorsqu'aucune inconnue ne s'isole facilement.
\end{attention}
\end{corrige}

\begin{corrige}[Exercice 9 — Problème d'atelier : Coûts de production]
\textbf{Barème indicatif (sur 10) :} modélisation 2 pts ; pivot de Gauss 4 pts ; coûts unitaires 1 pt ; coût de la commande 1,5 pt ; analyse critique 1,5 pt.

\textbf{1. Modélisation.} En exprimant les coûts en FCFA :
\[ (S) : \begin{cases} 10x + 5y + 3z = 1\,250\,000 \\ 8x + 7y + 2z = 1\,310\,000 \\ 12x + 4y + 5z = 1\,580\,000 \end{cases} \]

\textbf{2. Pivot de Gauss.} Matrice augmentée :
\[ \begin{pmatrix} 10 & 5 & 3 & \big| & 1\,250\,000 \\ 8 & 7 & 2 & \big| & 1\,310\,000 \\ 12 & 4 & 5 & \big| & 1\,580\,000 \end{pmatrix} \]
$L_2 \leftarrow 5L_2 - 4L_1$ : $(0 \;;\; 15 \;;\; -2 \;|\; 1\,550\,000)$. \quad $L_3 \leftarrow 5L_3 - 6L_1$ : $(0 \;;\; -10 \;;\; 7 \;|\; 400\,000)$.
\[ \begin{pmatrix} 10 & 5 & 3 & \big| & 1\,250\,000 \\ 0 & 15 & -2 & \big| & 1\,550\,000 \\ 0 & -10 & 7 & \big| & 400\,000 \end{pmatrix}
\xrightarrow{L_3 \leftarrow 3L_3 + 2L_2}
\begin{pmatrix} 10 & 5 & 3 & \big| & 1\,250\,000 \\ 0 & 15 & -2 & \big| & 1\,550\,000 \\ 0 & 0 & 17 & \big| & 4\,300\,000 \end{pmatrix} \]
Remontée :
\begin{align*}
17z &= 4\,300\,000 \Rightarrow z = \dfrac{4\,300\,000}{17} \approx 252\,941,18 \text{ FCFA}, \\
15y &= 1\,550\,000 + 2z = \dfrac{26\,350\,000 + 8\,600\,000}{17} = \dfrac{34\,950\,000}{17} \Rightarrow y = \dfrac{2\,330\,000}{17} \approx 137\,058,82 \text{ FCFA}, \\
10x &= 1\,250\,000 - 5y - 3z = \dfrac{21\,250\,000 - 11\,650\,000 - 12\,900\,000}{17} = -\dfrac{3\,300\,000}{17} \Rightarrow x = -\dfrac{330\,000}{17} \approx -19\,411,76 \text{ FCFA}.
\end{align*}

\textbf{3. Coûts unitaires.} Chaise : $x \approx -19\,411,76$ FCFA ; table : $y \approx 137\,058,82$ FCFA ; banc : $z \approx 252\,941,18$ FCFA.

\begin{attention}[Vigilance économique]
Le coût d'une chaise est \emph{négatif}, ce qui est impossible en réalité. Les données de l'énoncé sont donc incohérentes économiquement : en situation réelle, l'atelier devrait vérifier ses bilans. Le candidat doit signaler cette anomalie — c'est une attitude attendue dans « rédiger et justifier une conclusion ».
\end{attention}

\textbf{4. Coût de la commande.} $20x + 15y + 10z = \dfrac{-6\,600\,000 + 34\,950\,000 + 43\,000\,000}{17} = \dfrac{71\,350\,000}{17} \approx 4\,197\,058,82$ FCFA.

\textbf{5. Analyse critique.} Recette de vente : $20 \times 60\,000 + 15 \times 120\,000 + 10 \times 80\,000 = 1\,200\,000 + 1\,800\,000 + 800\,000 = 3\,800\,000$ FCFA.
Comme $3\,800\,000 < 4\,197\,058,82$, l'atelier réalise une \textbf{perte} de :
\[ \frac{71\,350\,000}{17} - 3\,800\,000 = \frac{6\,750\,000}{17} \approx 397\,058,82 \text{ FCFA}. \]
\textbf{Justifications attendues :} citer les opérations élémentaires sur les lignes, vérifier la cohérence des signes des solutions, comparer coût et recette avant de conclure.
\end{corrige}

\begin{corrige}[Exercice 10 — Problème de géométrie : Dimensions d'un terrain]
\textbf{Barème indicatif (sur 10) :} question 1 : 1 pt ; question 2 : 2 pts ; question 3 : 3 pts ; question 4 : 2 pts ; question 5 : 2 pts.

\textbf{1.} Périmètre : $2x + 2y = 60 \iff x + y = 30 \iff y = 30 - x$.

\textbf{2.} Aire : $xy \ge 200 \iff x(30 - x) \ge 200 \iff -x^2 + 30x - 200 \ge 0$, soit encore (en multipliant par $-1$ et en inversant le sens) :
\[ x^2 - 30x + 200 \le 0. \]

\textbf{3.} $\Delta = 900 - 800 = 100$ ; racines : $x_1 = \dfrac{30 - 10}{2} = 10$, $x_2 = \dfrac{30 + 10}{2} = 20$. Comme $a = 1 > 0$, le trinôme est négatif \emph{entre} les racines : $S = [10\,; 20]$.

\textbf{4.} La longueur $x$ peut prendre toute valeur de $[10\,; 20]$ m, la largeur étant alors $y = 30 - x$ (comprise entre 10 et 20 m). Oui, longueur et largeur sont \textbf{interchangeables} : le couple $(x\,; y) = (12\,; 18)$ et le couple $(18\,; 12)$ décrivent le même terrain. Si l'on impose la convention $x \ge y$ (longueur $\ge$ largeur), alors $x \in [15\,; 20]$.

\textbf{5.} Pour un périmètre fixe, la diagonale est minimale pour le \textbf{carré} : $x = y = 15$ m. L'aire vaut alors $225 \text{ m}^2 \ge 200 \text{ m}^2$ : la contrainte est respectée. La diagonale vaut $15\sqrt{2} \approx 21,21$ m.

\begin{attention}[Points de vigilance]
Ne pas oublier de transformer l'inéquation pour avoir un coefficient de $x^2$ positif avant d'appliquer la règle des signes ; vérifier que la solution retenue respecte la contrainte d'aire ; justifier l'appartenance du carré à l'ensemble des solutions ($15 \in [10\,;20]$).
\end{attention}
\end{corrige}

\begin{corrige}[Exercice 11 — Problème d'électricité : Loi des nœuds et des mailles]
\textbf{Barème indicatif (sur 10) :} écriture matricielle 2 pts ; pivot de Gauss 4 pts ; intensités et interprétation des signes 2 pts ; puissance 2 pts.

\textbf{1.} Le système s'écrit $AX = B$ avec :
\[ A = \begin{pmatrix} 1 & -1 & -1 \\ 2 & 0 & 1 \\ 0 & 4 & -1 \end{pmatrix}, \quad X = \begin{pmatrix} I_1 \\ I_2 \\ I_3 \end{pmatrix}, \quad B = \begin{pmatrix} 0 \\ 12 \\ 6 \end{pmatrix}. \]

\textbf{2. Pivot de Gauss.}
\[ \begin{pmatrix} 1 & -1 & -1 & \big| & 0 \\ 2 & 0 & 1 & \big| & 12 \\ 0 & 4 & -1 & \big| & 6 \end{pmatrix}
\xrightarrow{L_2 \leftarrow L_2 - 2L_1}
\begin{pmatrix} 1 & -1 & -1 & \big| & 0 \\ 0 & 2 & 3 & \big| & 12 \\ 0 & 4 & -1 & \big| & 6 \end{pmatrix}
\xrightarrow{L_3 \leftarrow L_3 - 2L_2}
\begin{pmatrix} 1 & -1 & -1 & \big| & 0 \\ 0 & 2 & 3 & \big| & 12 \\ 0 & 0 & -7 & \big| & -18 \end{pmatrix} \]
Remontée : $-7I_3 = -18 \Rightarrow I_3 = \dfrac{18}{7}$ A. Puis $2I_2 = 12 - 3I_3 = 12 - \dfrac{54}{7} = \dfrac{30}{7} \Rightarrow I_2 = \dfrac{15}{7}$ A. Enfin $I_1 = I_2 + I_3 = \dfrac{15}{7} + \dfrac{18}{7} = \dfrac{33}{7}$ A.

\textbf{3. Intensités.} $I_1 = \dfrac{33}{7} \approx 4,71$ A ; $I_2 = \dfrac{15}{7} \approx 2,14$ A ; $I_3 = \dfrac{18}{7} \approx 2,57$ A.
Vérification (maille gauche) : $2 \times \dfrac{33}{7} + \dfrac{18}{7} = \dfrac{84}{7} = 12$ \checkmark ; (maille droite) : $4 \times \dfrac{15}{7} - \dfrac{18}{7} = \dfrac{42}{7} = 6$ \checkmark.
Les trois intensités sont \textbf{positives} : les courants circulent bien dans les sens choisis sur le schéma. (Une intensité négative aurait signifié que le sens réel est opposé au sens fléché.)

\textbf{4. Puissance dissipée par $R_3$.}
\[ P = R_3 I_3^2 = 1 \times \left(\frac{18}{7}\right)^2 = \frac{324}{49} \approx 6,61 \text{ W}. \]

\begin{attention}[Points de vigilance]
Respecter l'ordre des inconnues dans la matrice ; effectuer les opérations $L_i \leftarrow L_i + \lambda L_j$ sur \emph{toute} la ligne, y compris le second membre ; interpréter physiquement le signe des intensités ; ne pas oublier l'unité (A, W) dans la conclusion.
\end{attention}
\end{corrige}

\newpage

\coursec{Fiche bilan}

\begin{popfigure}
\begin{center}\tcbox[colback=poporange,colframe=poporange,coltext=white,arc=9pt,boxsep=4pt,left=6pt,right=6pt]{\bfseries\small Équations, Inéquations, Systèmes (Terminale Tech)}\end{center}
\vspace{-2pt}
\begin{tcbraster}[raster columns=3,raster equal height=rows,raster column skip=6pt,raster row skip=6pt,enhanced,blankest]
\begin{tcolorbox}[enhanced,colback=white,colframe=poporange,boxrule=0.9pt,arc=3pt,title={Équations 2\textsuperscript{nd} degré},fonttitle=\bfseries\scriptsize,coltitle=white,colbacktitle=poporange,left=4pt,right=4pt,top=2pt,bottom=2pt,boxsep=1pt,toptitle=1.5pt,bottomtitle=1.5pt,halign=left]
\begin{itemize}[nosep,leftmargin=*,itemsep=1pt,font=\scriptsize]
\item Δ = b²-4ac
\item Racines : $\frac{-b\pm\sqrt{\Delta}}{2a}$
\item Forme canonique
\end{itemize}
\end{tcolorbox}
\begin{tcolorbox}[enhanced,colback=white,colframe=popgreen,boxrule=0.9pt,arc=3pt,title={Inéquations 2\textsuperscript{nd} degré},fonttitle=\bfseries\scriptsize,coltitle=white,colbacktitle=popgreen,left=4pt,right=4pt,top=2pt,bottom=2pt,boxsep=1pt,toptitle=1.5pt,bottomtitle=1.5pt,halign=left]
\begin{itemize}[nosep,leftmargin=*,itemsep=1pt,font=\scriptsize]
\item Signe du trinôme
\item Tableau de signes
\item f(x) > 0, < 0, ≥ 0
\end{itemize}
\end{tcolorbox}
\begin{tcolorbox}[enhanced,colback=white,colframe=poppurple,boxrule=0.9pt,arc=3pt,title={Inéquations 3\textsuperscript{rd} degré},fonttitle=\bfseries\scriptsize,coltitle=white,colbacktitle=poppurple,left=4pt,right=4pt,top=2pt,bottom=2pt,boxsep=1pt,toptitle=1.5pt,bottomtitle=1.5pt,halign=left]
\begin{itemize}[nosep,leftmargin=*,itemsep=1pt,font=\scriptsize]
\item Factorisation (racine évidente)
\item Produit de facteurs
\item Tableau de signes complet
\end{itemize}
\end{tcolorbox}
\begin{tcolorbox}[enhanced,colback=white,colframe=poppink,boxrule=0.9pt,arc=3pt,title={Systèmes 3×3 : Substitution},fonttitle=\bfseries\scriptsize,coltitle=white,colbacktitle=poppink,left=4pt,right=4pt,top=2pt,bottom=2pt,boxsep=1pt,toptitle=1.5pt,bottomtitle=1.5pt,halign=left]
\begin{itemize}[nosep,leftmargin=*,itemsep=1pt,font=\scriptsize]
\item Isoler une variable
\item Substituer dans les 2 autres
\item Résoudre le 2×2 obtenu
\item Remonter pour le triplet
\end{itemize}
\end{tcolorbox}
\begin{tcolorbox}[enhanced,colback=white,colframe=popgold,boxrule=0.9pt,arc=3pt,title={Systèmes 3×3 : Pivot de Gauss},fonttitle=\bfseries\scriptsize,coltitle=white,colbacktitle=popgold,left=4pt,right=4pt,top=2pt,bottom=2pt,boxsep=1pt,toptitle=1.5pt,bottomtitle=1.5pt,halign=left]
\begin{itemize}[nosep,leftmargin=*,itemsep=1pt,font=\scriptsize]
\item Matrice augmentée
\item Opérations sur lignes
\item Forme échelonnée
\item Remontée (back-substitution)
\end{itemize}
\end{tcolorbox}
\begin{tcolorbox}[enhanced,colback=white,colframe=popblue!70!popgreen,boxrule=0.9pt,arc=3pt,title={Modélisation \& Problèmes},fonttitle=\bfseries\scriptsize,coltitle=white,colbacktitle=popblue!70!popgreen,left=4pt,right=4pt,top=2pt,bottom=2pt,boxsep=1pt,toptitle=1.5pt,bottomtitle=1.5pt,halign=left]
\begin{itemize}[nosep,leftmargin=*,itemsep=1pt,font=\scriptsize]
\item Variables = inconnues réelles
\item Traduire contraintes
\item Résoudre \& Interpréter
\item Rédaction soignée
\end{itemize}
\end{tcolorbox}
\end{tcbraster}

\legende{Carte mentale de la leçon : 6 branches principales structurant le programme.}
\end{popfigure}

\begin{bilan}

\courssub{Définitions et notations clés}

\begin{itemize}
  \item \textbf{Trinôme du second degré} : $f(x)=ax^2+bx+c$ ($a\neq0$). \textbf{Discriminant} : $\Delta=b^2-4ac$.
  \item \textbf{Racines} : $x_1=\frac{-b-\sqrt{\Delta}}{2a},\; x_2=\frac{-b+\sqrt{\Delta}}{2a}$ (si $\Delta\ge0$). Forme factorisée : $a(x-x_1)(x-x_2)$.
  \item \textbf{Signe du trinôme} : même signe que $a$ en dehors des racines, signe contraire entre les racines (si $\Delta>0$).
  \item \textbf{Polynôme degré 3} : $P(x)=ax^3+bx^2+cx+d$. Recherche d'une racine évidente ($\pm1,\pm\frac{\text{diviseurs de }d}{\text{diviseurs de }a}$) puis division euclidienne pour factoriser.
  \item \textbf{Système linéaire 3×3} : $\begin{cases} a_1x+b_1y+c_1z=d_1\\ a_2x+b_2y+c_2z=d_2\\ a_3x+b_3y+c_3z=d_3 \end{cases}$. \textbf{Matrice augmentée} : $\left(\begin{array}{ccc|c} a_1\&b_1\&c_1\&d_1\\ a_2\&b_2\&c_2\&d_2\\ a_3\&b_3\&c_3\&d_3 \end{array}\right)$.
  \item \textbf{Pivot} : coefficient non nul utilisé pour annuler les coefficients situés en dessous (ou au-dessus) dans sa colonne.
\end{itemize}

\courssub{Théorèmes, formules et règles à connaître par cœur}

\begin{center}
\begin{tabularx}{\linewidth}{|>{\bfseries}l|X|}\hline
\textbf{Objet} & \textbf{Règle / Formule} \\ \hline
Racines trinôme & $\Delta=b^2-4ac$ ; $x_{1,2}=\frac{-b\mp\sqrt{\Delta}}{2a}$ ; $x_1+x_2=-\frac{b}{a},\; x_1x_2=\frac{c}{a}$ \\ \hline
Signe trinôme ($a>0$) & $\Delta<0$ : $>0\ \forall x$ ; $\Delta=0$ : $\ge0$, nul en $x_0$ ; $\Delta>0$ : $<0$ sur $]x_1,x_2[$, $>0$ ailleurs \\ \hline
Résolution $P(x)>0$ (deg 3) & 1. Factoriser $P(x)=a(x-r_1)(x-r_2)(x-r_3)$ (racines réelles). 2. Tableau de signes (règle du produit). 3. Lire les intervalles selon l'inégalité. \\ \hline
Substitution 3×3 & Isoler $x$ (ou $y,z$) dans l'équation la plus simple $\to$ substituer dans les 2 autres $\to$ système 2×2 $\to$ résoudre $\to$ remonter. \\ \hline
Pivot de Gauss & Opérations autorisées : (L1) Permuter 2 lignes, (L2) Multiplier une ligne par $k\neq0$, (L3) Ajouter à une ligne un multiple d'une autre. But : matrice triangulaire supérieure. \\ \hline
Nature solutions 3×3 & \textbf{Unique} : pivot sur chaque colonne variable. \textbf{Infinité} : au moins une colonne sans pivot (variable libre). \textbf{Aucune} : ligne $[0\ 0\ 0\ |\ k\neq0]$. \\ \hline
\end{tabularx}
\end{center}

\courssub{Méthodes express (étapes clés)}

\begin{methode}[Résoudre une inéquation du second degré $ax^2+bx+c > 0$]
\begin{enumerate}
  \item Calculer $\Delta = b^2-4ac$.
  \item Trouver les racines $x_1 \le x_2$ (si $\Delta\ge0$).
  \item Tracer le tableau de signes (parabole : signe de $a$ à l'extérieur, contraire à l'intérieur).
  \item Lire l'ensemble des solutions $S$ selon l'inégalité stricte ou large.
  \item Écrire la conclusion : $S = \dots$
\end{enumerate}
\end{methode}

\begin{methode}[Résoudre une inéquation du troisième degré $P(x) \le 0$]
\begin{enumerate}
  \item Chercher une racine évidente $r$ (tester $\pm1, \pm2, \dots$).
  \item Effectuer la division euclidienne de $P(x)$ par $(x-r)$ pour obtenir $Q(x)$ (deg 2).
  \item Factoriser $Q(x)$ si possible. Écrire $P(x)=a(x-r)(x-r_2)(x-r_3)$.
  \item Ordonner les racines $r_1 < r_2 < r_3$. Construire le tableau de signes complet (4 lignes : chaque facteur, produit $P$).
  \item Lire les intervalles où $P(x)$ a le signe demandé. Conclure.
\end{enumerate}
\end{methode}

\begin{methode}[Pivot de Gauss sur système 3×3]
\begin{enumerate}
  \item Écrire la matrice augmentée.
  \item \textbf{Colonne 1} : Pivot en (1,1). Si nul, permuter avec une ligne ayant un non-nul en col 1. Annuler les coeffs (2,1) et (3,1) par combinaisons linéaires.
  \item \textbf{Colonne 2} : Pivot en (2,2) (après étape 1). Si nul, permuter lignes 2 et 3. Annuler coeff (3,2).
  \item \textbf{Colonne 3} : Vérifier pivot (3,3).
  \item \textbf{Remontée} : Exprimer $z$ depuis ligne 3, puis $y$ ligne 2, puis $x$ ligne 1.
  \item Discuter les cas particuliers (ligne nulle = paramètre ; ligne impossible = $\varnothing$).
\end{enumerate}
\end{methode}

\end{bilan}

\begin{aretenir}[Checklist avant le Probatoire]
\begin{enumerate}
  \item $\square$ Je sais calculer $\Delta$ et donner les racines exactes d'un trinôme (forme simplifiée si $\Delta$ n'est pas un carré parfait).
  \item $\square$ Je résous correctement $f(x)=0$, $f(x)>0$, $f(x)\ge0$, $f(x)<0$ pour un trinôme (tableau de signes + conclusion ensembliste).
  \item $\square$ Je factorise un polynôme de degré 3 après avoir trouvé une racine évidente (division euclidienne ou identité remarquable).
  \item $\square$ Je construis un tableau de signes complet pour un polynôme de degré 3 factorisé (4 lignes : 3 facteurs + produit).
  \item $\square$ Je résous une inéquation de degré 3 (ex: $x^3-2x^2-x+2 \le 0$) et j'écris la solution sous forme d'intervalles ou d'union d'intervalles.
  \item $\square$ Je résous un système 3×3 par la méthode de substitution (choix judicieux de l'équation de départ, substitution soignée, remontée).
  \item $\square$ Je résous un système 3×3 par le pivot de Gauss (matrice augmentée, opérations sur lignes notées $L_i \leftarrow \dots$, forme échelonnée, remontée).
  \item $\square$ J'identifie et j'exprime correctement les 3 cas : solution unique (triplet), infinité de solutions (paramétrage avec $t\in\mathbb{R}$), aucune solution (incompatibilité).
  \item $\square$ Je sais modéliser un problème concret (mélange, géométrie, électricité, économie) par un système 3×3 ou une inéquation.
  \item $\square$ Je rédige une conclusion finale en français, en lien avec l'énoncé (ex: « La longueur est de 5 m », « L'ensemble solution est $[-1;2]$ »).
  \item $\square$ Je vérifie mes solutions en les remplaçant dans les équations/équations initiales (contrôle numérique rapide).
  \item $\square$ Je présente ma copie proprement : alignement des égalités, flèches du tableau de variations/signes correctes ($\nearrow, \searrow$), encadrement des résultats.
\end{enumerate}
\end{aretenir}

\findecours

\end{document}
